\documentclass[a4paper]{article}

\usepackage{amsfonts, amsmath, amssymb, amsthm}
\usepackage{fullpage}
\usepackage{hyperref}
\usepackage{booktabs}

\newtheorem{theorem}{Theorem}[section]
\newtheorem{lemma}[theorem]{Lemma}
\newtheorem{corollary}[theorem]{Corollary}
\newtheorem{proposition}[theorem]{Proposition}
\newtheorem{remark}[theorem]{Remark}
\newtheorem{definition}[theorem]{Definition}
\newtheorem{example}[theorem]{Example}
\newtheorem{problem}[theorem]{Problem}
\newtheorem{heuristic}[theorem]{Heuristic}

\newcommand{\F}{\mathbb{F}}
\newcommand{\Z}{\mathbb{Z}}
\newcommand{\Q}{\mathbb{Q}}
\newcommand{\N}{\mathbb{N}}
\newcommand{\ord}[2][]{\operatorname{ord}_{#1}\!\left(#2\right)}
\newcommand{\abs}[1]{\left|#1\right|}
\newcommand{\wt}{\operatorname{wt}}
\newcommand{\Sp}{\operatorname{Sp}}
\newcommand{\SL}{\operatorname{SL}}
\newcommand{\PSL}{\operatorname{PSL}}
\DeclareMathOperator{\Res}{Res}
\DeclareMathOperator{\Tr}{Tr}
\DeclareMathOperator{\Nm}{N}
\DeclareMathOperator{\Gal}{Gal}
\renewcommand{\qedsymbol}{$\blacksquare$}

\setlength{\parindent}{0em}
\setlength{\parskip}{1em}

\begin{document}

\title{Reciprocal Five-Term Relations among $p^j$-th Roots of Unity\\
in Characteristic Two: Reductions, Obstructions, and the Wieferich Plateau}
\author{William Boyles}
\date{\today}
\maketitle

\begin{abstract}
	Let $p>5$ be prime, let $j\geq2$, and let $\theta\in\overline{\F}_2$ have
	exact order $p^j$.
	We study the equation
	\[
		1+\theta^{a}+\theta^{-a}+\theta^{b}+\theta^{-b}=0,
	\]
	whose only solutions are conjectured to have $p\mid a$ and $p\mid b$.
	This is equivalent to the stabilization
	$d(p^k-1)=d(p-1)$ of the nullity of the bounded
	\textit{Lights Out} grid.  This article consolidates the algebraic,
	cyclotomic, local, additive, coding-theoretic, dynamical, elliptic, and
	symplectic attacks on that conjecture, including the points at which
	each method becomes sharp but inconclusive.

	The equation is encoded by the reciprocal quartic
	$Q_Y(T)=T^4+T^3+YT^2+T+1$ and by a shifted-Dickson sequence
	$C_n(S)$ whose repeated roots are exactly the quartics dividing
	$T^n+1$.  We prove that every repeated root of $C_n$ has multiplicity
	exactly two, with $C_n^{[2]}(Y)=Y^{-2}$, and obtain an exact bridge to
	Sutner's nullity formula.  An integer lift factors over the maximal real
	cyclotomic field and turns stabilization into the simultaneous oddness
	of an internal and a mixed collision norm.  Arithmetic dynamics gives
	the same off-plateau descent through the degree-two covering
	$u\mapsto u+u^{-1}$ and identifies irreducible factors with doubling
	orbits of collision pairs.  The equation also becomes a translated
	coordinate intersection on an elliptic curve and the trace-one locus of
	a regular semisimple element of $\Sp_4$; both reformulations clarify the
	geometry but return to the original incidence problem on the plateau.

	If
	$c_p=v_p(2^{\ord[p]{2}}-1)$, every solution has
	$\ord{\theta^a},\ord{\theta^b}\leq p^{c_p}$, so only
	$2\leq j\leq c_p$ remains.  The whole theorem is therefore equivalent
	to one endpoint assertion at order $p^{c_p}$.  Characteristic-zero
	vanishing sums, cyclotomic units, ray classes, power residues, modular
	vanishing-sum theorems, cyclic-code bounds, and maximal-torus
	classification all stop at that same endpoint.  The semi-local residue
	module is explicitly
	$\Z[\bar G]/(\sigma_2-2)$, and its Frobenius, sign, distribution, and
	norm relations leave the primitive orbit values freely parameterized;
	explicit phantom systems show that these relations, even together with
	the known trace constraints, do not force avoidance.

	Genuinely additive input does prove new cases.  The Artin--Schreier law
	$\Tr(\bar\lambda_x^{-1})=\delta$ and the global law
	$\Tr(\theta^x)+\Tr(\theta^y)=\ord[p^j]{2}\bmod2$ exclude several
	families, including $d(3510)=d(72)=d(88)=d(232)=0$.  At the opposite
	extreme, the Mersenne and Fermat cases admit exact binary Kloosterman
	formulas.  The statement is false in other residue characteristics, so
	any proof must use characteristic-two additive structure.  Exact
	searches verify the endpoint for both known base-$2$ Wieferich primes,
	$1093$ and $3511$, and a complete resultant computation settles
	$p^j=49$.  No proof is presently known for a hypothetical new
	Wieferich endpoint.
\end{abstract}

\tableofcontents

\section{The problem}\label{sec:problem}

Let $F_n(x)\in\F_2[x]$ be the Fibonacci polynomials
\[
	F_0(x)=0,\qquad F_1(x)=1,\qquad F_n(x)=xF_{n-1}(x)+F_{n-2}(x),
\]
and let $d(n)$ be the nullity over $\F_2$ of the $n\times n$ bounded
\textit{Lights Out} grid, so that
\begin{equation}\label{eq:sutner}
	d(n)=\deg\gcd\!\left(F_{n+1}(x),\,F_{n+1}(x+1)\right)
\end{equation}
by Sutner's formula \cite{Sutner00sigma-automataand}.
For odd $N>1$ the distinct roots of $F_N$ are
\begin{equation}\label{eq:strata}
	\bigcup_{\substack{M\mid N\\M>1}}R_M,
	\qquad
	R_M=\left\{\zeta+\zeta^{-1}:\ \ord{\zeta}=M\right\},
\end{equation}
and every root has multiplicity two \cite{HUNZIKER2004465}.

A companion manuscript at \path{tex/finite_fields/finite_fields.tex} proves
the \emph{prime stabilization theorem}: with
\begin{equation}\label{eq:c}
	c=c_p=v_p\!\left(2^{\ord[p]{2}}-1\right),
\end{equation}
one has $d(p^k-1)=d(p^c-1)$ for all $k\geq c$.
Since $c=1$ unless $p$ is a Wieferich prime, this settles
$d(p^k-1)=d(p-1)$ for every non-Wieferich prime, and exact computations
settle the two known Wieferich primes $1093$ and $3511$.
The residual obstruction is a single, purely finite-field statement.

\begin{problem}\label{prob:main}
	Let $p>5$ be prime, let $j\geq2$, and let $\theta\in\overline{\F}_2$
	have exact order $p^j$.
	Suppose that
	\begin{equation}\label{eq:main}
		1+\theta^{a}+\theta^{-a}+\theta^{b}+\theta^{-b}=0 .
	\end{equation}
	Must $p\mid a$ and $p\mid b$?
\end{problem}

The hypothesis $p>5$ is genuinely needed: for $p=5$ and $j=1$ the relation
$1+\theta+\theta^{4}+\theta^{2}+\theta^{3}=0$ holds in $\F_{16}$.
Section~\ref{sec:charzero} explains why $5$ is the only obstruction of this
kind.

Throughout, $u=\theta^a$ and $v=\theta^b$, and we write
\[
	S_{N}=\left\{(a,b)\in(\Z/N\Z)^2:\ 1+\theta^{a}+\theta^{-a}+\theta^{b}+\theta^{-b}=0\right\}
\]
for $\theta$ of exact order $N$.

The organization follows the successive changes of language used in the
attack.  Section~\ref{sec:reformulations} develops the shifted-Dickson,
reciprocal-quartic, elliptic, and nullity dictionaries.
Section~\ref{sec:degree} proves descent off the Wieferich plateau and states
the single endpoint lemma equivalent to the full theorem.
Sections~\ref{sec:collision-norms} and~\ref{sec:dynamics} give the integer,
real-cyclotomic, and arithmetic-dynamical forms of the same obstruction.
Section~\ref{sec:sp4} records the symplectic reformulation and explains why
maximal-torus and character theory do not decide the trace-one incidence.
Sections~\ref{sec:charzero}--\ref{sec:units} treat characteristic zero and
the ratios eliminated by cyclotomic units.
Sections~\ref{sec:reduction} and~\ref{sec:power} develop the local
circular-unit, semi-local module, ray-class, power-residue, and additive
trace attacks.
Sections~\ref{sec:modular} and~\ref{sec:coding} explain the exact limitations
of modular vanishing sums and cyclic/Zetterberg codes.
Sections~\ref{sec:counterexamples} and~\ref{sec:quantitative} record the
counterexamples in other characteristics and the size heuristic.
Section~\ref{sec:computations} lists the exact certificates, and
Section~\ref{sec:summary} collects what is proved, what failed, and what
remains.
\subsection{Elementary normalizations}

\begin{lemma}\label{lem:degenerate}
	Let $N>1$ be odd with $3\nmid N$, let $\theta\in\overline{\F}_2$ have
	exact order $N$, and let $(a,b)\in S_N$.
	Then $a\not\equiv0$, $b\not\equiv0$ and $b\not\equiv\pm a$ modulo $N$.
	In particular the five roots of unity $1,\theta^{\pm a},\theta^{\pm b}$
	are pairwise distinct.
\end{lemma}

\begin{proof}
	If $a\equiv0$ then $\theta^a+\theta^{-a}=1+1=0$, so
	$\theta^b+\theta^{-b}=1$ and $\theta^{b}$ is a root of $T^2+T+1$,
	hence of exact order $3$.
	Then $3\mid N$, a contradiction.
	The case $b\equiv0$ is symmetric.
	If $b\equiv\pm a$ then $\theta^{b}+\theta^{-b}=\theta^{a}+\theta^{-a}$
	and the left side of \eqref{eq:main} equals $1$.
\end{proof}

\begin{lemma}\label{lem:symmetry}
	$S_{N}$ is stable under $(a,b)\mapsto(2a,2b)$, under
	$(a,b)\mapsto(\pm a,\pm b)$ and under $(a,b)\mapsto(b,a)$.
	If $\theta$ is replaced by $\theta^{t}$ with $\gcd(t,N)=1$ then
	$S_{N}$ is replaced by $t^{-1}S_{N}$.
	In particular whether $S_{N}$ is empty does not depend on the choice of
	$\theta$.
\end{lemma}

\begin{proof}
	The first claim is the Frobenius $x\mapsto x^2$ applied to
	\eqref{eq:main}; the second and third are the substitutions
	$a\mapsto-a$ and $a\leftrightarrow b$, which permute the summands.
	The last claim is immediate from
	$1+(\theta^{t})^{a}+\cdots=1+\theta^{ta}+\cdots$.
\end{proof}

\subsection{Why the problem is exactly the residual obstruction}

\begin{lemma}\label{lem:d-count}
	Let $N>1$ be odd with $3\nmid N$ and let $\theta\in\overline{\F}_2$
	have exact order $N$.
	Then
	\[
		d(N-1)=\tfrac12\abs{S_N}.
	\]
\end{lemma}

\begin{proof}
	By \eqref{eq:strata} the roots of $F_N$ are exactly the elements
	$u+u^{-1}$ with $u\in\mu_N\setminus\{1\}$, each of multiplicity two.
	A common root of $F_N(x)$ and $F_N(x+1)$ is an $\alpha$ with
	$\alpha=u+u^{-1}$ and $\alpha+1=v+v^{-1}$ for some
	$u,v\in\mu_N\setminus\{1\}$, i.e.\ with
	$1+u+u^{-1}+v+v^{-1}=0$.
	Writing $u=\theta^a$ and $v=\theta^b$ identifies the set of such
	$\alpha$ with $S_N$ modulo the action of
	$\{\pm1\}\times\{\pm1\}$ on $(a,b)$; the action is free because
	$a\equiv-a$ would force $a\equiv0$, which Lemma~\ref{lem:degenerate}
	excludes.
	Hence the number of common roots is $\abs{S_N}/4$, and
	\eqref{eq:sutner} together with multiplicity two gives
	$d(N-1)=2\cdot\abs{S_N}/4$.
\end{proof}

\begin{corollary}\label{cor:reduction}
	If Problem~\ref{prob:main} has an affirmative answer for a prime $p>5$
	and all $j\geq2$, then $d(p^k-1)=d(p-1)$ for all $k\geq1$.
\end{corollary}

\begin{proof}
	Let $k\geq2$ and let $\theta$ have exact order $p^k$.
	Any element of $S_{p^k}$ has $p\mid a$ and $p\mid b$, so
	$\theta^a,\theta^b\in\mu_{p^{k-1}}$ and the same relation is an element
	of $S_{p^{k-1}}$ after replacing $\theta$ by $\theta^{p}$.
	Iterating, $\abs{S_{p^k}}=\abs{S_{p}}$, and Lemma~\ref{lem:d-count}
	gives $d(p^k-1)=d(p-1)$.
\end{proof}

Lemma~\ref{lem:d-count} was verified numerically against the independent
grid-nullity code in \path{code/kernel_size.py}: for $p=5,17,31$ the counts
$\abs{S_p}=8,16,40$ match $d(4)=4$, $d(16)=8$, $d(30)=20$, and for
$p=7,11,41,43,73$ the empty solution sets match
$d(6)=d(10)=d(40)=d(42)=d(72)=0$.

\subsection{The reciprocal quartic}

Set
\[
	Q_Y(T)=T^4+T^3+YT^2+T+1\in\overline{\F}_2[Y][T],
\]
a monic reciprocal quartic whose roots have product $1$.

\begin{lemma}\label{lem:quartic}
	Let $\alpha\in\overline{\F}_2$ and $Y=\alpha^2+\alpha$.
	Then
	\[
		Q_Y(T)=\left(T^2+\alpha T+1\right)\left(T^2+(\alpha+1)T+1\right).
	\]
	Consequently, for $u,v\in\overline{\F}_2^{\times}$ the following are
	equivalent.
	\begin{enumerate}
		\item $1+u+u^{-1}+v+v^{-1}=0$.
		\item $u,u^{-1},v,v^{-1}$ are exactly the roots of $Q_Y$ for
		      $Y=\alpha^2+\alpha$ with $\alpha=u+u^{-1}$.
	\end{enumerate}
	Moreover Problem~\ref{prob:main} has a negative answer for the pair
	$(p,j)$ if and only if there is $Y\in\overline{\F}_2$ such that
	$Q_Y(T)\mid T^{p^j}+1$ and at least one root of $Q_Y$ has exact order
	$p^{j}$.
\end{lemma}

\begin{proof}
	Expanding the product gives
	$T^4+(\alpha+\alpha+1)T^3+(1+\alpha(\alpha+1)+1)T^2+(\alpha+\alpha+1)T+1$,
	which is $Q_{\alpha^2+\alpha}(T)$.
	The quadratic $T^2+\alpha T+1$ has root set $\{u,u^{-1}\}$ precisely
	when $u+u^{-1}=\alpha$; likewise $T^2+(\alpha+1)T+1$ has root set
	$\{v,v^{-1}\}$ precisely when $v+v^{-1}=\alpha+1$, which is
	statement~(1).
	For the last assertion, $T^{p^j}+1$ has exactly the elements of
	$\mu_{p^j}$ as roots, each simple, so $Q_Y\mid T^{p^j}+1$ says
	precisely that $u,v\in\mu_{p^j}$; and a counterexample to
	Problem~\ref{prob:main} is exactly a relation in which $\theta^a$ or
	$\theta^b$ has exact order $p^j$.
\end{proof}

The next lemma identifies the resultant that appears in
Lemma~\ref{lem:quartic} with a Fibonacci-polynomial expression, and is the
bridge between the quartic picture and \eqref{eq:sutner}.
Let $D_N$ denote the Dickson polynomial of the first kind
\cite[Ch.~7]{lidl-niederreiter}, determined by
$D_N(t+t^{-1})=t^{N}+t^{-N}$.

\begin{lemma}\label{lem:dickson}
	Over $\F_2$ one has $D_N(s)=s\,F_N(s)$ for every $N\geq0$.
	Consequently, for odd $N$,
	\[
		\Res_T\!\left(Q_Y(T),\,T^{N}+1\right)
		=
		Y\cdot\Res_s\!\left(s^2+s+Y,\,F_N(s)\right),
	\]
	and for $Y=\alpha^2+\alpha$ with $\alpha\ne0,1$,
	\[
		\Res_s\!\left(s^2+s+Y,\,F_N(s)\right)=F_N(\alpha)F_N(\alpha+1).
	\]
\end{lemma}

\begin{proof}
	Both $D_N$ and $sF_N$ satisfy the recursion $X_N=sX_{N-1}+X_{N-2}$ over
	$\F_2$; the Dickson recursion is $D_N=sD_{N-1}-D_{N-2}$, and $-1=1$.
	The initial values agree, since $D_0=2=0=s\cdot F_0$ and
	$D_1=s=s\cdot F_1$.

	For the resultant, let $t_1,\ldots,t_4$ be the roots of $Q_Y$, paired
	as $\{t,t^{-1}\}$ by Lemma~\ref{lem:quartic}.
	In characteristic two,
	$(t^{N}+1)(t^{-N}+1)=t^{N}+t^{-N}=D_N(t+t^{-1})$, so
	\[
		\Res_T\!\left(Q_Y,\,T^{N}+1\right)
		=\prod_{i=1}^{4}\left(t_i^{N}+1\right)
		=D_N(\alpha)D_N(\alpha+1)
		=\alpha(\alpha+1)F_N(\alpha)F_N(\alpha+1),
	\]
	and $\alpha(\alpha+1)=\alpha^2+\alpha=Y$.
	The roots of $s^2+s+Y$ are $\alpha$ and $\alpha+1$, which gives the
	second displayed identity.
\end{proof}

\section{Shifted-Dickson, nullity, and elliptic reformulations}
\label{sec:reformulations}

The reciprocal quartic can be packaged into one explicit polynomial
sequence.  This is useful for two reasons: it records the repeated-root
locus with its exact multiplicity, and it turns the desired stabilization
into ordinary square-freeness and coprimality statements.

\begin{definition}\label{def:Cn}
	Since $D_n(X)+D_n(X+1)$ is invariant under $X\mapsto X+1$, there is a
	unique polynomial $C_n(S)\in\F_2[S]$ such that
	\begin{equation}\label{eq:Cn-def}
		C_n(X^2+X)=D_n(X)+D_n(X+1).
	\end{equation}
\end{definition}

\begin{proposition}[Power sums and generating function]
\label{prop:Cn-powersum}
	The polynomial $C_n(S)$ is the $n$-th Newton power sum of the four roots
	of $Q_S(T)$:
	\[
		C_n(S)=\sum_{Q_S(\rho)=0}\rho^n.
	\]
	Moreover
	\begin{equation}\label{eq:Cn-gf}
		\sum_{n\geq0}C_n(S)t^n
		=\frac{t(1+t)^2}{Q_S(t)}
		=\frac{t+t^3}{1+t+St^2+t^3+t^4}.
	\end{equation}
	Hence
	\[
		C_{n+4}=C_{n+3}+S C_{n+2}+C_{n+1}+C_n,
		\qquad
		(C_0,C_1,C_2,C_3)=(0,1,1,S).
	\]
\end{proposition}

\begin{proof}
	Write $S=X^2+X$.  By Lemma~\ref{lem:quartic}, the roots of $Q_S$ are
	$\zeta,\zeta^{-1},\eta,\eta^{-1}$ with
	$\zeta+\zeta^{-1}=X$ and $\eta+\eta^{-1}=X+1$.  Their $n$-th power
	sum is therefore $D_n(X)+D_n(X+1)=C_n(S)$.
	If $g(t)=Q_S(t)$, reciprocity gives
	$\prod_\rho(1-\rho t)=g(t)$, and
	\[
		-\frac{t g'(t)}{g(t)}
		=\sum_{\rho}\frac{\rho t}{1-\rho t}
		=\sum_{n\geq1}\left(\sum_\rho\rho^n\right)t^n.
	\]
	In characteristic two, $g'(t)=1+t^2=(1+t)^2$ and $C_0=4=0$.
	This proves \eqref{eq:Cn-gf}; clearing the denominator gives the
	recurrence.
\end{proof}

\begin{proposition}[Elementary identities]\label{prop:Cn-identities}
	For all $m,n\geq1$:
	\[
		C_{2n}=C_n^2,
		\qquad
		m\mid n\Longrightarrow C_m\mid C_n.
	\]
\end{proposition}

\begin{proof}
	The first identity is Frobenius applied to the power sum.  For the
	second write $n=mk$ and use Dickson composition:
	\[
		D_n(X)+D_n(X+1)
		=D_k(D_m(X))+D_k(D_m(X+1)).
	\]
	The difference $A+B$ divides $D_k(A)+D_k(B)$ in characteristic two,
	so the right side is divisible by
	$D_m(X)+D_m(X+1)$.  The substitution
	$\F_2[S]\hookrightarrow\F_2[X]$, $S\mapsto X^2+X$, is injective.
\end{proof}

For $f(S)=\sum_i a_iS^i$, write
$f^{[r]}=\sum_i a_i\binom{i}{r}S^{i-r}$ for its $r$-th Hasse derivative.

\begin{proposition}[Hasse-derivative generating function]
\label{prop:Cn-hasse}
	For every $r\geq0$,
	\begin{equation}\label{eq:Cn-hasse}
		\sum_{n\geq0}C_n^{[r]}(S)t^n
		=\frac{t^{2r+1}(1+t)^2}{Q_S(t)^{r+1}}.
	\end{equation}
\end{proposition}

\begin{proof}
	The polynomial $Q_S(t)=t^2S+(1+t+t^3+t^4)$ is affine in $S$.
	The Hasse expansion of its reciprocal gives
	$\partial_S^{[r]}Q_S(t)^{-1}=t^{2r}Q_S(t)^{-r-1}$.
	Apply this to \eqref{eq:Cn-gf}.
\end{proof}

For odd $n$, the polynomial $F_n$ is the square of a square-free
polynomial \cite[Prop.~2.11]{HUNZIKER2004465}; write
\[
	F_n(X)=H_n(X)^2.
\]
Then $D_n(X)=XH_n(X)^2$ and
\begin{equation}\label{eq:Cn-H}
	C_n(X^2+X)=XH_n(X)^2+(X+1)H_n(X+1)^2.
\end{equation}

\begin{lemma}[Repeated roots and dividing quartics]\label{lem:Cn-repeated}
	Let $n$ be odd.  A value $Y\in\overline{\F}_2$ is a repeated root of
	$C_n$ if and only if
	\[
		Y=X_0^2+X_0
	\]
	for a common root $X_0$ of $F_n(X)$ and $F_n(X+1)$.  Equivalently,
	$Y\neq0$ and $Q_Y(T)\mid T^n+1$.
\end{lemma}

\begin{proof}
	Differentiating \eqref{eq:Cn-H} gives
	\[
		\frac{d}{dX}C_n(X^2+X)
		=H_n(X)^2+H_n(X+1)^2.
	\]
	At a repeated root $X_0$, the two values
	$H_n(X_0)$ and $H_n(X_0+1)$ are equal.  Substitution in
	\eqref{eq:Cn-H} then shows that their common value is zero.  The
	converse is immediate.  The map $X\mapsto X^2+X$ identifies the pair
	$\{X_0,X_0+1\}$, and Lemma~\ref{lem:quartic} together with the standard
	root parametrization of $F_n$ gives the quartic criterion.
	Since $F_n(0)=1$ for odd $n$, neither $X_0=0$ nor $X_0=1$ occurs, so
	$Y\neq0$.
\end{proof}

\begin{theorem}[Multiplicity Two]\label{thm:Cn-mult2}
	If $n$ is odd and $Y$ is a repeated root of $C_n$, then
	\begin{equation}\label{eq:Cn-second}
		C_n^{[2]}(Y)=Y^{-2}\neq0.
	\end{equation}
	Thus every repeated root has multiplicity exactly two, and
	\begin{equation}\label{eq:Cn-factor}
		C_n=A_nB_n^2,\qquad
		\gcd(C_n,C_n')=B_n^2,
	\end{equation}
	where $A_n$ and $B_n$ are square-free and coprime and the roots of
	$B_n$ are precisely the values in Lemma~\ref{lem:Cn-repeated}.
\end{theorem}

\begin{proof}
	Let $Y=X_0^2+X_0$ be repeated, and write the two reciprocal root pairs
	of $Q_Y$ as
	$\{\zeta,\zeta^{-1}\}$ and $\{\eta,\eta^{-1}\}$, where
	$X_0=\zeta+\zeta^{-1}$ and
	$X_0+1=\eta+\eta^{-1}$.  All four roots are distinct $n$-th roots of
	unity.  Put $m=(n-1)/2$.  If $X=t+t^{-1}$, then
	\[
		H_n(X)=\frac{t^n+1}{t^m(t+1)}.
	\]
	Differentiating at $t=\zeta$ and using
	$dX/dt=(t+1)^2/t^2$ gives
	\[
		H_n'(X_0)^2=\frac{\zeta^3}{(\zeta+1)^6},
	\]
	and the analogous identity holds with $\eta$ and $X_0+1$.
	The coefficient of $(X-X_0)^2$ in \eqref{eq:Cn-H} is therefore
	\[
		X_0H_n'(X_0)^2+(X_0+1)H_n'(X_0+1)^2
		=\frac{\zeta^2}{(\zeta+1)^4}
		 \frac{\eta^2}{(\eta+1)^4}.
	\]
	Since $\zeta/(\zeta+1)^2=X_0^{-1}$ and
	$\eta/(\eta+1)^2=(X_0+1)^{-1}$, this is
	\[
		\left(\frac1{X_0}+\frac1{X_0+1}\right)^2
		=(X_0^2+X_0)^{-2}=Y^{-2}.
	\]
	The substitution $S=X^2+X$ has derivative $1$, so this coefficient is
	exactly $C_n^{[2]}(Y)$.  The factorization follows.
\end{proof}

\begin{theorem}[Bridge to Sutner's nullity]\label{thm:Cn-bridge}
	For odd $n$,
	\begin{equation}\label{eq:Cn-bridge}
		\gcd\!\left(F_n(X),F_n(X+1)\right)
		=\gcd(C_n,C_n')(X^2+X).
	\end{equation}
	In particular
	\[
		d(n-1)=4\deg B_n=2\deg\gcd(C_n,C_n').
	\]
\end{theorem}

\begin{proof}
	With $F_n=H_n^2$, the square-free polynomial
	$\gcd(H_n(X),H_n(X+1))$ has roots in translation pairs
	$\{X_0,X_0+1\}$.  Lemma~\ref{lem:Cn-repeated} identifies the polynomial
	obtained after quotienting those pairs as $B_n$, so
	\[
		\gcd(H_n(X),H_n(X+1))=B_n(X^2+X).
	\]
	Square and use \eqref{eq:Cn-factor}.  The degree statement follows
	from Sutner's formula \eqref{eq:sutner}.
\end{proof}

\subsection{Product coordinates and an elliptic translation}

There is also a geometric normalization that retains both multiplicative
coordinates.  Put
\[
	z=uv,\qquad r=u/v.
\]
Then, in characteristic two,
\begin{equation}\label{eq:product-coordinates}
	1+u+u^{-1}+v+v^{-1}=0
	\quad\Longleftrightarrow\quad
	(z+z^{-1})(r+r^{-1})=1.
\end{equation}
Indeed the right side is the square of
$u+u^{-1}+v+v^{-1}$.

\begin{proposition}[Elliptic translation model]\label{prop:elliptic-translation}
	Let
	\[
		E:\quad Y^2+XY=X^3+X,\qquad T=(1,0)\in E(\F_2).
	\]
	The point $T$ has order four.  Away from the degenerate values
	$z,r\in\{0,1\}$, the curve in \eqref{eq:product-coordinates} is
	isomorphic to the open subset of $E$ under
	\[
		P=\left(z,(r+1)(z+1)^2\right).
	\]
	Under this correspondence
	\[
		X(P)=z,\qquad X(P+T)=r.
	\]
	Consequently a five-term relation of conductor $n$ is a point $P$ for
	which both $X(P)$ and $X(P+T)$ lie in $\mu_n$.
\end{proposition}

\begin{proof}
	Set $y=(r+1)(z+1)^2$.  After clearing denominators,
	\eqref{eq:product-coordinates} is
	\[
		z^2r^2+zr+z^2+r^2+1=0,
	\]
	which is equivalent to
	$y^2+zy=z^3+z$.  Thus $(z,y)\in E$.
	The addition law on the binary Weierstrass model gives
	\[
		X(P+T)=1+\frac{Y(P)}{(X(P)+1)^2}=r.
	\]
	Conversely, the curve equation implies
	\[
		(X+X^{-1})\bigl(X(P+T)+X(P+T)^{-1}\bigr)=1.
	\]
	A direct doubling calculation gives $2T=(0,0)$ and $2(0,0)=O$.
\end{proof}

\begin{remark}[What the elliptic model does and does not give]
	The model explains the genus-one geometry behind the equation and
	relates the two multiplicative coordinates by translation by a rational
	four-torsion point.  It does not turn the problem into an ordinary
	torsion question on $E$: the condition is that two \emph{$X$-coordinates}
	lie in the subgroup $\mu_{p^j}\subset\overline{\F}_2^\times$, not that
	$P$ belongs to a specified subgroup of $E$.  Isogenies, complex
	multiplication, and point-count estimates therefore return to the same
	subgroup-incidence problem.  This is why the elliptic attack gives
	exact formulas in full-torus cases but does not resolve a Wieferich
	plateau.
\end{remark}

\section{The degree argument and the Wieferich plateau}\label{sec:degree}

For odd $M>1$ put
\[
	\rho_2(M)=\min\{r>0:\ 2^{r}\equiv\pm1\pmod M\},
\]
so that $[\F_2(\alpha):\F_2]=\rho_2(M)$ for $\alpha\in R_M$; this is the
field-extension-size lemma of the companion manuscript at
\path{tex/finite_fields/finite_fields.tex}.

\begin{lemma}\label{lem:orders}
	Let $p$ be an odd prime and $c=c_p$ as in \eqref{eq:c}.
	Then for every $i\geq1$,
	\[
		\ord[p^i]{2}=\ord[p]{2}\,p^{\max(0,\,i-c)},
		\qquad
		\rho_2(p^{i})=\rho_2(p)\,p^{\max(0,\,i-c)} .
	\]
\end{lemma}

\begin{proof}
	The first identity is the lifting-the-exponent computation carried out
	in \path{tex/finite_fields/finite_fields.tex}.
	For the second, $\rho_2(M)=\ord[M]{2}$ or $\ord[M]{2}/2$ according as
	$-1\notin\langle2\rangle$ or $-1\in\langle2\rangle$ in
	$(\Z/M\Z)^{\times}$.
	The group $(\Z/p^i\Z)^{\times}$ is cyclic and contains a unique element
	of order two, namely $-1$; hence $-1\in\langle2\rangle$ if and only if
	$\ord[p^i]{2}$ is even, i.e.\ if and only if $\ord[p]{2}$ is even,
	because $p^{\max(0,i-c)}$ is odd.
	The criterion is therefore independent of $i$, and the second identity
	follows from the first.
\end{proof}

\begin{definition}
	The \emph{base-$2$ plateau} of $p$ is the set of exponents
	$1\leq i\leq c_p$.
	Equivalently, $i$ lies on the plateau if and only if
	$\mu_{p^{i}}\subseteq\F_{2^{f_1}}$ with $f_1=\ord[p]{2}$, if and only
	if $2$ splits completely in $\Q(\zeta_{p^{i}})/\Q(\zeta_p)$.
	The plateau has length $\geq2$ exactly when $p$ is a Wieferich prime.
\end{definition}

\begin{proposition}[Frobenius-orbit restriction]\label{prop:frobenius-orbit}
	Suppose \eqref{eq:main} holds, $\gcd(a,p)=1$, and
	$b\equiv\pm2^k a\pmod{p^j}$.
	Put $d=\rho_2(p^j)$ and reduce $k$ modulo $d$ to $0\leq k<d$.
	Then $d$ is even and $k=d/2$.
	In particular, when $\rho_2(p^j)$ is odd, a counterexample cannot pair
	two trace values in the same Frobenius orbit.
\end{proposition}

\begin{proof}
	Set $\alpha=\theta^a+\theta^{-a}$. Since $\theta^a$ has exact order
	$p^j$, one has $[\F_2(\alpha):\F_2]=d$. The congruence on $b$ and
	\eqref{eq:main} give
	\[
		\alpha^{2^k}=\theta^b+\theta^{-b}=\alpha+1.
	\]
	Applying the same Frobenius power once more yields
	$\alpha^{2^{2k}}=(\alpha+1)^{2^k}=\alpha$, so $d\mid2k$.
	One cannot have $d\mid k$, since that would give
	$\alpha^{2^k}=\alpha\neq\alpha+1$.
	With $0<k<d$, the divisibility $d\mid2k$ therefore forces
	$d$ even and $k=d/2$.
\end{proof}

\begin{proposition}[Orbit divisibility]\label{prop:orbit-divisibility}
	Let $\Delta_j=\abs{S_{p^j}}-\abs{S_{p^{j-1}}}$, where the lower solution
	set is embedded by $(a,b)\mapsto(pa,pb)$, and put
	$\rho=\rho_2(p^j)$. Then
	\[
		\Delta_j\in
		\begin{cases}
			8\rho\,\Z,&\rho\ \text{odd},\\
			4\rho\,\Z,&\rho\ \text{even}.
		\end{cases}
	\]
	Thus any nonzero new stratum contains at least $8\rho$ ordered solutions
	when $\rho$ is odd and at least $4\rho$ when $\rho$ is even.
\end{proposition}

\begin{proof}
	On solution pairs $(u,v)=(\theta^a,\theta^b)$, independent inversion of
	the two coordinates and transposition generate a dihedral group of order
	$8$. Its action is free for $p>3$: a fixed inversion forces a coordinate
	to be $1$, while a fixed signed transposition forces $v=u^{\pm1}$ and
	then $(u+u^{-1})^2=1$, so $u$ has order $3$.

	Frobenius acts by $(u,v)\mapsto(u^2,v^2)$. If
	$f=\ord[p^j]{2}$ is odd, then $f=\rho$ and its cyclic group meets the
	dihedral group trivially. If $f$ is even, then $f=2\rho$ and
	Frobenius to the power $\rho$ is simultaneous inversion, so the combined
	effective group still has order $8\rho$.

	A nontrivial stabilizer in this effective group must involve a signed
	transposition. With one coordinate primitive, such a stabilizer gives
	$2^{2k}\equiv-1\pmod{p^j}$; the $+1$ alternative again forces
	$v=u^{\pm1}$. Hence it can occur only when $\rho$ is even, with
	$k\equiv\rho/2\pmod\rho$, and then its order is exactly $2$.
	The new solutions therefore split into orbits of size $8\rho$ when
	$\rho$ is odd and of size $8\rho$ or $4\rho$ when $\rho$ is even.
\end{proof}

\begin{theorem}\label{thm:A}
	Let $p>3$ be prime, let $\theta\in\overline{\F}_2$ have exact order
	$p^j$, and suppose \eqref{eq:main} holds.
	Then
	\[
		\ord{\theta^{a}}\leq p^{c},
		\qquad
		\ord{\theta^{b}}\leq p^{c}.
	\]
	In particular $p^{\,j-c}$ divides both $a$ and $b$ whenever $j>c$.
\end{theorem}

\begin{proof}
	By Lemma~\ref{lem:degenerate} the orders $p^{i_1}=\ord{\theta^a}$ and
	$p^{i_2}=\ord{\theta^b}$ satisfy $i_1,i_2\geq1$.
	Put $\alpha=\theta^{a}+\theta^{-a}$, so $\alpha\in R_{p^{i_1}}$ and
	$\alpha+1\in R_{p^{i_2}}$ by \eqref{eq:main}.
	Since $\F_2(\alpha)=\F_2(\alpha+1)$, we get
	$\rho_2(p^{i_1})=\rho_2(p^{i_2})$, whence
	$\max(0,i_1-c)=\max(0,i_2-c)$ by Lemma~\ref{lem:orders}.
	Therefore either $i_1,i_2\leq c$, and we are done, or $i_1=i_2=:i>c$.

	Assume $i>c$.
	Both $\theta^{a}$ and $\theta^{b}$ have exact order $p^{i}$, so writing
	$\zeta=\theta^{a}$ we have $\theta^{b}=\zeta^{s}$ for some $s$ with
	$\gcd(s,p)=1$, and $\zeta$ is a root of
	\begin{equation}\label{eq:Ps}
		P(x)=1+x+x^{p^{i}-1}+x^{s}+x^{p^{i}-s}\in\F_2[x].
	\end{equation}
	Let $m=p^{\,i-c}>1$.
	Then $\zeta^{m}$ has exact order $p^{c}$, and its minimal polynomial
	$g$ over $\F_2$ has degree $\ord[p^{c}]{2}=\ord[p]{2}$ by
	Lemma~\ref{lem:orders}.
	Hence
	\[
		\deg g(x^{m})=\ord[p]{2}\,p^{\,i-c}=\ord[p^{i}]{2},
	\]
	which is the degree of the minimal polynomial of $\zeta$; as $\zeta$ is
	a root of $g(x^m)$, that polynomial \emph{is} the minimal polynomial of
	$\zeta$, and therefore $g(x^{m})\mid P(x)$.

	Decompose $P(x)=\sum_{r=0}^{m-1}x^{r}P_r(x^{m})$, which is possible in
	exactly one way.
	Writing $P=g(x^{m})H$ and $H=\sum_r x^{r}H_r(x^{m})$ and comparing the
	unique decompositions gives $P_r=gH_r$ for every $r$; in particular
	$g\mid P_0$.
	The five exponents in \eqref{eq:Ps} reduce modulo $m$ to
	$0,\ 1,\ -1,\ s,\ -s$.
	Since $m$ is an odd prime power exceeding $1$ and $\gcd(s,p)=1$, none
	of $1,-1,s,-s$ is $\equiv0\pmod m$.
	Hence $P_0=1$ and $g\mid1$, contradicting $\deg g\geq1$.
\end{proof}

\begin{corollary}\label{cor:nonwieferich}
	Problem~\ref{prob:main} has an affirmative answer for every
	non-Wieferich prime $p>3$ and every $j\geq2$, and more generally
	whenever $j>c_p$.
	Consequently a counterexample requires a Wieferich prime $p$ with
	$c_p\geq2$ and an exponent $j$ with $2\leq j\leq c_p$; that is, the
	problem is open \emph{exactly on the base-$2$ Wieferich plateau}.
\end{corollary}

\begin{proof}
	If $p$ is not Wieferich then $c_p=1<2\leq j$, so
	Theorem~\ref{thm:A} gives $p^{\,j-1}\mid a,b$, and in particular
	$p\mid a,b$.
	If $j>c_p$ the same argument applies verbatim.
	A counterexample therefore has $j\leq c_p$, and $j\geq2$ forces
	$c_p\geq2$.
\end{proof}

\begin{corollary}\label{cor:2^64}
	Problem~\ref{prob:main} has an affirmative answer for every prime
	$5<p<2^{64}$ and every $j\geq2$.
	Hence $d(p^k-1)=d(p-1)$ for every such $p$ and every $k$.
\end{corollary}

\begin{proof}
	The only Wieferich primes below $2^{64}$ are $1093$ and $3511$
	\cite{CrandallSearch,KellerRichstein2005,DoraisKlyve2011,PrimeGridWieferich},
	and both have $c_p=2$.
	Corollary~\ref{cor:nonwieferich} handles every other $p$, and handles
	$j\geq3$ for $p\in\{1093,3511\}$.
	The remaining case $j=2$ for $p\in\{1093,3511\}$ is settled by the
	direct searches of Section~\ref{sec:computations}, which find no
	solutions; alternatively it follows from the exact polynomial gcd
	evaluations $d(1093^2-1)=d(3511^2-1)=0$ recorded
	in \path{tex/finite_fields/finite_fields.tex} together with
	Lemma~\ref{lem:d-count}.
	Corollary~\ref{cor:reduction} converts this into the statement about
	$d$.
\end{proof}

Theorem~\ref{thm:A} is a small sharpening of the argument in the companion
manuscript, which treats only the case of equal orders.
The point of stating it in this form is Corollary~\ref{cor:nonwieferich}:
the plateau is not merely where the known proof stops, it is the exact
locus of the residual problem.
Section~\ref{sec:counterexamples} shows that on the analogous plateau in
other residue characteristics the statement can genuinely fail.

\subsection{The single endpoint statement}

The preceding reduction can be sharpened from a finite list of plateau
levels to one statement at the top of the plateau.

\begin{theorem}[Endpoint equivalence]\label{thm:endpoint-equivalence}
	Fix a prime $p>5$ with $c=c_p\geq2$.  The following are equivalent.
	\begin{enumerate}
		\item Problem~\ref{prob:main} has an affirmative answer for every
		      $j\geq2$.
		\item Whenever $\Theta$ has exact order $p^c$ and
		      \[
			      1+\Theta^a+\Theta^{-a}+\Theta^b+\Theta^{-b}=0,
		      \]
		      one has $p^{c-1}\mid a$ and $p^{c-1}\mid b$.
		\item $B_{p^c}=B_p$.
		\item Writing $F_n=H_n^2$ as in Section~\ref{sec:reformulations},
		      \begin{equation}\label{eq:endpoint-gcd}
			      \gcd\!\left(
				      \frac{H_{p^c}(X)}{H_p(X)},
				      H_{p^c}(X+1)
			      \right)=1.
		      \end{equation}
	\end{enumerate}
\end{theorem}

\begin{proof}
	By Theorem~\ref{thm:A}, levels above $c$ introduce no new solution.
	Thus (1) implies (2).  Conversely suppose (2), and let
	$2\leq j\leq c$.  Choose $\Theta$ of order $p^c$ and put
	$\theta=\Theta^{p^{c-j}}$.  A relation at level $p^j$ becomes one at
	level $p^c$ with exponents $p^{c-j}a,p^{c-j}b$.  Statement (2) then
	gives $p^{j-1}\mid a,b$, hence in particular $p\mid a,b$.
	This proves (1).

	Lemma~\ref{lem:Cn-repeated} identifies the roots of $B_{p^k}$ with
	the common-root pairs whose two multiplicative orders divide $p^k$.
	Therefore $B_{p^c}=B_p$ says exactly that every endpoint relation has
	both orders dividing $p$, which is (2).  Finally, the quotient
	$H_{p^c}/H_p$ records the roots whose order is new above level $p$;
	its coprimality with the translate of $H_{p^c}$ is exactly the absence
	of a common-root pair with at least one new coordinate.  This is (3).
\end{proof}

\begin{definition}\label{def:Cn-primitive}
	For $j\geq1$ define the shifted-Dickson primitive quotient
	\[
		\Pi_{p^j}(S)=\frac{C_{p^j}(S)}{C_{p^{j-1}}(S)},
		\qquad C_{p^0}=C_1=1.
	\]
	The letter $\Pi$ is used to avoid confusion with the cyclotomic
	polynomial $\Phi_{p^j}$.
\end{definition}

\begin{proposition}[Polynomial stabilization criterion]
\label{prop:Cn-stabilization}
	For every $j\geq1$ the following are equivalent:
	\begin{enumerate}
		\item $B_{p^j}=B_{p^{j-1}}$;
		\item
		      $\gcd(C_{p^j},C_{p^j}')=
		       \gcd(C_{p^{j-1}},C_{p^{j-1}}')$;
		\item $\Pi_{p^j}$ is square-free and coprime to
		      $C_{p^{j-1}}$.
	\end{enumerate}
	Moreover $B_{p^j}=B_{p^c}$ for every $j\geq c$, so the full theorem
	is equivalent to $B_{p^c}=B_p$.
\end{proposition}

\begin{proof}
	The first two statements are equivalent by
	\eqref{eq:Cn-factor}.  Because every repeated root of every $C_n$ has
	exact multiplicity two, a new repeated root at level $p^j$ either
	appears twice in $\Pi_{p^j}$ or appears once in both
	$\Pi_{p^j}$ and $C_{p^{j-1}}$; these are exactly the two failures in
	(3).  The final statement is Theorem~\ref{thm:A} translated through
	Lemma~\ref{lem:Cn-repeated}.
\end{proof}

\section{Integer lifts and cyclotomic collision norms}
\label{sec:collision-norms}

The shifted-Dickson target has a useful integral lift.  It does not remove
the endpoint obstruction, but it isolates it as the parity of two explicit
integer norms and makes clear that both a primitive--primitive and a
primitive--lower collision must be excluded.

\begin{proposition}[Integer lift]\label{prop:integer-lift}
	For odd $n$ there is a monic polynomial
	$\widetilde C_n\in\Z[Y]$ such that
	\begin{equation}\label{eq:integer-lift}
		D_n(X)+D_n(X+1)
		=(2X+1)\widetilde C_n(X^2+X),
		\qquad
		\widetilde C_n\equiv C_n\pmod2.
	\end{equation}
\end{proposition}

\begin{proof}
	The involution $\iota(X)=-1-X$ fixes $Y=X^2+X$.  For odd $n$, the
	integer Dickson polynomial $D_n$ is odd, so
	$D_n(X)+D_n(X+1)$ is anti-invariant under $\iota$.
	Write an arbitrary element of $\Z[X]$ uniquely as $A(Y)+XB(Y)$.
	Anti-invariance forces $B=2A$, and hence the polynomial is
	$(2X+1)A(Y)$.  Its leading term shows that $A=\widetilde C_n$ is
	monic.  Reducing modulo two gives \eqref{eq:Cn-def}.
\end{proof}

For odd $k$, put
\[
	L_k[A,B]=\frac{D_k(A)+D_k(B)}{A+B}\in\Z[A,B].
\]

\begin{proposition}[Composition and telescoping]\label{prop:integer-telescope}
	If $m$ and $k$ are odd, then
	\[
		\frac{\widetilde C_{mk}}{\widetilde C_m}
		=L_k[D_m(X),D_m(X+1)]\in\Z[Y].
	\]
	In particular $\widetilde C_m\mid\widetilde C_n$ whenever $m\mid n$,
	and
	\begin{equation}\label{eq:integer-telescope}
		\widetilde\Pi_{p^j}(Y)
		:=\frac{\widetilde C_{p^j}(Y)}
		        {\widetilde C_{p^{j-1}}(Y)}
		=L_p[D_{p^{j-1}}(X),D_{p^{j-1}}(X+1)].
	\end{equation}
\end{proposition}

\begin{proof}
	Dickson composition gives
	\[
		D_{mk}(X)+D_{mk}(X+1)
		=D_k(A)+D_k(B)=(A+B)L_k[A,B],
	\]
	with $A=D_m(X)$ and $B=D_m(X+1)$.  The quotient is invariant under
	$X\mapsto-1-X$, hence lies in $\Z[Y]$.  Cancel the common factor
	$2X+1$ from \eqref{eq:integer-lift}.
\end{proof}

\begin{theorem}[Discriminant/resultant certificate]
\label{thm:integer-certificate}
	At level $p^j$, the equality
	$B_{p^j}=B_{p^{j-1}}$ holds if and only if both integers
	\[
		\operatorname{disc}(\widetilde\Pi_{p^j})
		\quad\text{and}\quad
		\Res(\widetilde\Pi_{p^j},\widetilde C_{p^{j-1}})
	\]
	are odd.
\end{theorem}

\begin{proof}
	Reduction modulo two commutes with discriminants and resultants.
	Thus the first integer is odd exactly when $\Pi_{p^j}$ is square-free,
	and the second is odd exactly when $\Pi_{p^j}$ is coprime to
	$C_{p^{j-1}}$.  Apply Proposition~\ref{prop:Cn-stabilization}.
\end{proof}

\subsection{Factorization over the maximal real cyclotomic field}

Let
\[
	\lambda_k=\zeta_n^k+\zeta_n^{-k},
	\qquad
	1\leq k\leq\frac{n-1}{2},
\]
and let $\psi_m$ be the minimal polynomial of
$\zeta_m+\zeta_m^{-1}$.  Put
\[
	\mathcal R_n(\lambda)=\prod_{\substack{d\mid n\\d>1}}\psi_d(\lambda).
\]

\begin{proposition}[Real-cyclotomic factorization]
\label{prop:real-factorization}
	For odd $n$,
	\begin{equation}\label{eq:real-factorization}
		\widetilde C_n(Y)
		=\prod_{k=1}^{(n-1)/2}
		 \bigl((\lambda_k+2)Y+\lambda_k^2-3\bigr)
		=\Res_\lambda\!\left(
			\mathcal R_n(\lambda),
			\lambda^2+Y\lambda+2Y-3
		\right).
	\end{equation}
	Its roots are
	\[
		y(\lambda)=\frac{3-\lambda^2}{\lambda+2},
	\]
	and its primitive quotient is obtained by replacing
	$\mathcal R_n$ with $\psi_{p^j}$.
\end{proposition}

\begin{proof}
	The standard factorization of
	$(D_n(A)+D_n(B))/(A+B)$ over
	$\Q(\zeta_n+\zeta_n^{-1})$ gives, after substituting
	$B=A+1$ and $Y=A^2+A$, the displayed linear factors; equivalently,
	direct substitution shows that every $y(\lambda_k)$ is a root, and
	the degrees agree.
	The leading coefficient is
	$\prod_k(\lambda_k+2)=\prod_{k=1}^{n-1}(1+\zeta_n^k)=1$.
	Grouping the trace values by exact order gives the primitive statement.
\end{proof}

Define the collision factor
\[
	\nu(\lambda,\mu)=(\lambda+2)(\mu+2)-1.
\]
The elementary identity
\begin{equation}\label{eq:y-difference}
	y(\lambda)-y(\mu)
	=\frac{(\mu-\lambda)\nu(\lambda,\mu)}
	       {(\lambda+2)(\mu+2)}
\end{equation}
separates the ordinary cyclotomic discriminant from the new source of
two-adic divisibility.

\begin{theorem}[Two collision norms]\label{thm:two-collision-norms}
	Let $\Lambda_j$ be the primitive $p^j$ trace values and let
	$\Lambda_{<j}$ be the nontrivial trace values of order dividing
	$p^{j-1}$.  Then:
	\begin{enumerate}
		\item the parity of
		      $\operatorname{disc}(\widetilde\Pi_{p^j})$ is the parity of
		      the internal product
		      \[
			      N_j^{\mathrm{int}}
			      =\prod_{\{\lambda,\mu\}\subset\Lambda_j}
			       \nu(\lambda,\mu);
		      \]
		\item the parity of
		      $\Res(\widetilde\Pi_{p^j},\widetilde C_{p^{j-1}})$ is the
		      parity of the mixed product
		      \[
			      N_j^{\mathrm{mix}}
			      =\prod_{\lambda\in\Lambda_j,\ \mu\in\Lambda_{<j}}
			       \nu(\lambda,\mu).
		      \]
	\end{enumerate}
	Consequently stabilization at level $p^j$ is equivalent to the
	simultaneous oddness of the internal and mixed collision norms.
\end{theorem}

\begin{proof}
	Equation~\eqref{eq:y-difference} expresses the discriminant and
	resultant as products of ordinary trace differences times the stated
	collision factors.  Since $2$ is unramified in every odd-conductor
	cyclotomic field, the discriminants of
	$\mathcal R_n$ and $\psi_{p^j}$ are odd.  The factors
	$\lambda+2$ are cyclotomic units of odd norm.  Thus all two-adic
	divisibility is carried by the two displayed products.
	Theorem~\ref{thm:integer-certificate} finishes the proof.
\end{proof}

\begin{remark}[Circular-unit form]
	If
	$\epsilon_a=\zeta^{-a/2}(1+\zeta^a)$, then
	$\epsilon_a^2=\lambda_a+2$, so
	\[
		\nu(\lambda_a,\lambda_b)
		=\epsilon_a^2\epsilon_b^2-1.
	\]
	At a prime above two this vanishes exactly when
	$\epsilon_a\epsilon_b\equiv1$.  With
	$s_a=\zeta^{-a/2}(1-\zeta^a)$ one has
	$\epsilon_a=\sigma_2(s_a)/s_a$, and the same congruence is equivalent
	to $s_as_b\equiv1$.  Thus the two norms here are precisely the
	primitive--primitive and primitive--lower sine-unit obstructions of
	Section~\ref{sec:reduction}.
\end{remark}

\subsection{The collision involution and relative norms}

Set
\[
	\mathfrak m(x)=\frac1{x+2}-2=-\frac{2x+3}{x+2}.
\]
It is an involution, conjugate under $U=x+2$ to $U\mapsto U^{-1}$, and
\[
	\nu(\lambda,\mu)=(\lambda+2)(\mu-\mathfrak m(\lambda)).
\]
Modulo two, $\mathfrak m(x)=x^{-1}$.  Thus the internal norm is odd
exactly when the primitive trace polynomial is coprime modulo two to its
reciprocal, and the mixed norm is odd exactly when no primitive trace is
the inverse of a lower trace.

For a monic polynomial $h$, write
\[
	h^{\mathfrak m}(X)
	=(X+2)^{\deg h}h(\mathfrak m(X)).
\]
The diagonal contribution to
$\Res(\psi_n,\psi_n^{\mathfrak m})$ is odd for odd $n\neq3$, because
$\psi_n(-2)=\pm1$ and $\psi_n(-1),\psi_n(-3)$ are odd.  Hence
\begin{equation}\label{eq:collision-resultants}
	N_j^{\mathrm{mix}}\ \text{is odd}
	\Longleftrightarrow
	\Res\!\left(\psi_{p^j},
		\Bigl(\prod_{i<j}\psi_{p^i}\Bigr)^{\mathfrak m}\right)
	\ \text{is odd},
\end{equation}
and the two-adic valuation of
$\Res(\psi_{p^j},\psi_{p^j}^{\mathfrak m})$ is twice that of the internal
product.

There is also an exact one-step norm identity.  If
$K_i^+=\Q(\zeta_{p^i}+\zeta_{p^i}^{-1})$,
$\lambda=\theta+\theta^{-1}$ has exact level $p^j$,
$\lambda_p=\theta^p+\theta^{-p}$, and $\mu\in K_{j-1}^+$, then
\begin{equation}\label{eq:relative-collision-norm}
	\Nm_{K_j^+/K_{j-1}^+}
	\left((\lambda+2)(\mu+2)-1\right)
	=(\mu+2)^p\left(\lambda_p-D_p(\mathfrak m(\mu))\right).
\end{equation}
Multiplying this identity over lower trace orbits gives a recurrence of
ordinary cyclotomic resultants.  This makes the obstruction computable
without constructing $C_{p^j}$, but it does not force its parity: on the
plateau the right side can vanish modulo a chosen prime above two.

\begin{theorem}[Mixed oddness off the plateau]
\label{thm:mixed-odd}
	If $j>c_p$, then $N_j^{\mathrm{mix}}$ is odd.
\end{theorem}

\begin{proof}
	A mixed collision gives $\theta$ of exact order $p^j$ and $\eta$ of
	order dividing $p^{j-1}$ such that
	\[
		(\theta+\theta^{-1})(\eta+\eta^{-1})=1.
	\]
	The inverse lower trace lies in a field of degree dividing
	$f_{j-1}=\ord[p^{j-1}]{2}$, and $\theta$ satisfies a quadratic over
	that field.  Hence $f_j\mid2f_{j-1}$.  For every $j>c_p$ one has
	$f_j=p f_{j-1}$, which would force the odd prime $p$ to divide $2$.
\end{proof}

The same degree argument says nothing about the internal norm, since both
traces already lie at exact level $p^j$.  The global degree theorem,
Theorem~\ref{thm:A}, rules out new internal collisions off the plateau by
using the five-term sparsity; collision-norm algebra alone does not.

\begin{theorem}[Plateau inversion criterion]\label{thm:plateau-inversion}
	Assume $j\leq c_p$, and in the common residue field put
	\[
		\mathcal T_i
		=\{x+x^{-1}:x\in\mu_{p^i}\}.
	\]
	Then:
	\begin{enumerate}
		\item $N_j^{\mathrm{mix}}$ is even if and only if
		      \[
			      \mathcal T_{j-1}^{-1}
			      \cap(\mathcal T_j\setminus\mathcal T_{j-1})
			      \neq\emptyset;
		      \]
		\item $N_j^{\mathrm{int}}$ is even if and only if
		      $\mathcal T_j\setminus\mathcal T_{j-1}$ meets its own
		      inverse.
	\end{enumerate}
\end{theorem}

\begin{proof}
	On the plateau all the groups $\mu_{p^i}$ lie in one finite field, and
	$\mathfrak m$ reduces to inversion.  Distinct pairs
	$\{x,x^{-1}\}$ give distinct traces, so the primitive trace values are
	exactly $\mathcal T_j\setminus\mathcal T_{j-1}$.  The two assertions
	are therefore direct translations of the mixed and internal collision
	conditions.
\end{proof}

\begin{remark}[What was learned]
	The integral route is an exact reformulation, not a parity theorem.
	It corrects a tempting but insufficient reduction to a single
	discriminant: the mixed resultant is an independent obstruction and
	must also be odd.  Off the plateau its parity follows from degree
	growth, while on the plateau the problem becomes the concrete question
	whether a finite set of primitive trace values meets its own inverse or
	the inverse of the lower trace set.  No general cyclotomic resultant
	formula determines those intersections.
\end{remark}

\section{Arithmetic dynamics of the collision locus}\label{sec:dynamics}

Let $f=D_p$ over $\F_2$.  Dickson composition gives
$D_{p^j}=f^{\circ j}$, so the roots of $D_{p^j}$ form the depth-$j$
preimage tree of the fixed point $0$.

\begin{proposition}[Critically fixed structure]\label{prop:dynamics-critical}
	The map $f$ fixes $0$.  Its finite critical points are exactly the
	nonzero roots of $D_p$, each of ramification index two, and every one
	of them maps to $0$.
\end{proposition}

\begin{proof}
	For odd $p$, $D_p(X)=XH_p(X)^2$, so
	$f'(X)=H_p(X)^2=F_p(X)$.  The square-free roots of $H_p$ are therefore
	the critical points, and $D_p$ sends all of them to zero.  The point
	$0$ is a simple fixed point because $H_p(0)=1$.
\end{proof}

Let
\[
	\mathcal Z_j=(f^{\circ j})^{-1}(0),\qquad
	\operatorname{Coll}_j=\mathcal Z_j\cap(\mathcal Z_j+1).
\]
Lemma~\ref{lem:Cn-repeated} identifies $\operatorname{Coll}_j$, modulo
$X\mapsto X+1$, with the repeated roots of $C_{p^j}$.

\begin{proposition}[Chebyshev linearization]\label{prop:dynamics-linearization}
	The degree-two map
	\[
		\phi:\mathbb G_m\longrightarrow\mathbb A^1,
		\qquad \phi(u)=u+u^{-1},
	\]
	satisfies
	\[
		f^{\circ j}\circ\phi=\phi\circ[u\mapsto u^{p^j}],
	\]
	and maps $\mu_{p^j}$ two-to-one onto $\mathcal Z_j$, except above
	$\phi(1)=0$.
\end{proposition}

\begin{proof}
	This is the Dickson identity
	$D_{p^j}(u+u^{-1})=u^{p^j}+u^{-p^j}$
	\cite{lidl-dickson}.
\end{proof}

Fix a generator $\zeta$ of $\mu_{p^j}$ and write
$\lambda_a=\zeta^a+\zeta^{-a}$.  A \emph{collision pair} is an unordered
pair $\{a,b\}$ with $\lambda_a+\lambda_b=1$, modulo independent sign
changes.  It is new at level $j$ when at least one of $a,b$ is prime to
$p$.

\begin{theorem}[Covering-degree bound]\label{thm:dynamics-covering}
	If $\{a,b\}$ is a mixed collision pair with $p\nmid a$ and $p\mid b$,
	then
	\[
		\ord[p^j]{2}\mid2\ord[p^{j-1}]{2}.
	\]
	Hence no mixed collision exists for $j>c_p$.
\end{theorem}

\begin{proof}
	The lower trace $\lambda_b$ lies in
	$\F_{2^{f_{j-1}}}$, and $\lambda_a=1+\lambda_b$ lies there as well.
	The primitive root $\zeta^a$ satisfies
	$T^2+\lambda_aT+1=0$, a quadratic over that field.  Thus
	$f_j=[\F_2(\zeta^a):\F_2]\mid2f_{j-1}$.
	For $j>c_p$, $f_j=pf_{j-1}$.
\end{proof}

\begin{theorem}[Doubling-orbit factorization]\label{thm:dynamics-orbits}
	The collision pairs decompose into orbits under
	\[
		\{a,b\}\longmapsto\{2a,2b\}.
	\]
	An orbit of size $d$ contributes one irreducible factor of degree $d$
	to $B_{p^j}$, and every irreducible factor arises in this way.
\end{theorem}

\begin{proof}
	Frobenius gives $\lambda_a^2=\lambda_{2a}$ and therefore
	\[
		(\lambda_a^2+\lambda_a)^2
		=\lambda_{2a}^2+\lambda_{2a}.
	\]
	The degree of the minimal polynomial of a repeated root is the size of
	its Frobenius orbit.  The map from a collision pair to
	$Y=\lambda_a^2+\lambda_a$ is injective modulo independent signs and
	interchange, because $T^2+T+Y$ recovers the unordered pair
	$\{\lambda_a,\lambda_b\}$.
\end{proof}

\begin{remark}[Why dynamics also stops]
	On the plateau $f_j=f_{j-1}$, so the covering-degree bound is vacuous.
	The monodromy is the same cyclic-by-inversion group in every residue
	characteristic, while Section~\ref{sec:counterexamples} gives genuine
	plateau failures outside characteristic two.  Dynamics therefore
	identifies the exact finite object to control---doubling orbits of
	collision pairs in one finite field---but monodromy alone cannot prove
	that this object is empty.
\end{remark}

\section{The symplectic and group-theoretic attack}\label{sec:sp4}

A separable reciprocal quartic is naturally a characteristic polynomial in
rank two.  This places the problem inside a well-understood finite group,
but the classification controls the ambient tori rather than the required
trace-one incidence.

\begin{theorem}[$\Sp_4$ reformulation]\label{thm:sp4-equivalence}
	A counterexample at level $p^j$ is equivalent to a regular semisimple
	element $g\in\Sp_4(\F_{2^d})$, for some $d$, such that
	\[
		\operatorname{charpol}(g)
		=Q_Y(T)=T^4+T^3+YT^2+T+1,
	\]
	all eigenvalues lie in $\mu_{p^j}$, and
	$\operatorname{ord}(g)=p^j$.  Equivalently, it is a regular
	semisimple element of pure $p$-power order and trace one.
\end{theorem}

\begin{proof}
	A semisimple symplectic element has eigenvalues
	$\lambda,\lambda^{-1},\mu,\mu^{-1}$, and its characteristic
	polynomial is
	\[
		T^4+(\alpha+\beta)T^3+\alpha\beta T^2
		     +(\alpha+\beta)T+1,
	\]
	where $\alpha=\lambda+\lambda^{-1}$ and
	$\beta=\mu+\mu^{-1}$.  In characteristic two its trace is one exactly
	when $\alpha+\beta=1$, which is the five-term equation.  For $p>5$ the
	four eigenvalues are distinct by Lemma~\ref{lem:degenerate}, and at
	least one has exact order $p^j$.  Conversely, the eigenvalues of such a
	$g$ give the required exponents.  The realization and rational
	conjugacy statement for a separable reciprocal quartic are standard
	for symplectic groups \cite{Wall1963,Enomoto1972}.
\end{proof}

Here $Y$ is a fixed element of $\overline{\F}_2$, not an indeterminate.
Let $k=\F_2(Y)\subseteq\overline{\F}_2$ be the finite field it generates,
put $q=\abs{k}$, and let $\alpha$ be a root of $T^2+T+Y$.

\begin{theorem}[Maximal-torus dichotomy]\label{thm:sp4-tori}
	The regular class in Theorem~\ref{thm:sp4-equivalence} lies in one of
	the following tori.
	\begin{enumerate}
		\item If $\Tr_{k/\F_2}(Y)=0$, then $\alpha\in k$ and the class lies
		      in a product of two like $\SL_2$ tori, of order
		      $(q-1)^2$ or $(q+1)^2$.
		\item If $\Tr_{k/\F_2}(Y)=1$, then $\alpha\notin k$, the two
		      symplectic planes are Galois conjugate, and the class lies
		      in a cyclic torus of order $q^2-1$ or $q^2+1$.
	\end{enumerate}
	The mixed product type $(q-1)(q+1)$ cannot contain the required
	$p^j$-eigenvalue configuration.
\end{theorem}

\begin{proof}
	The Artin--Schreier criterion decides whether $T^2+T+Y$ splits over
	$k$.  In the split case the two reciprocal quadratics are separately
	defined over $k$.  Each is either split, with eigenvalues in
	$k^\times$, or anisotropic, with eigenvalues in the norm-one circle of
	order $q+1$.  An odd prime power cannot divide both $q-1$ and $q+1$,
	so the two factors have the same type.

	In the nonsplit case Frobenius over $k$ interchanges $\alpha$ and
	$\alpha+1$, hence interchanges the two reciprocal pairs.  The four
	eigenvalues form one orbit in a cyclic torus of order $q^2-1$ or
	$q^2+1$.  This is the standard rank-two torus classification
	\cite{Carter1985,Enomoto1972}.
\end{proof}

On a Wieferich plateau, $p^j$ divides the order of one of these tori, so
elements of order $p^j$ are guaranteed.  What is not guaranteed is trace
one.  In product type that condition is
\[
	t_1+t_1^{-1}+t_2+t_2^{-1}=1,
	\qquad t_1,t_2\in\mu_{p^j},
\]
which is the original equation; in cyclic type the same equation holds
with $t_2$ a Frobenius conjugate of $t_1$.  Thus ``order is cheap, trace one
is not.''

\subsection{Why character theory and $\PSL_2$ do not close the gap}

For a regular element in a maximal torus $T$, the exact class count is a
sum over Weyl orbits in $T[p^j]$ satisfying trace one.  Deligne--Lusztig
characters and Enomoto's character table reorganize this sum, but the inner
condition is still the same incidence on $T[p^j]$.  Averaging the trace-one
indicator by additive characters gives the Kloosterman expression of
Section~\ref{sec:quantitative}; its square-root error is much larger than
the expected main term on a plateau.

There is also a superficially similar trace identity in
$\PSL_2(q)=\SL_2(q)$ for $q$ even:
\[
	\operatorname{tr}(A)\operatorname{tr}(C)
	=\operatorname{tr}(AC)+\operatorname{tr}(A^{-1}C).
\]
Requiring $AC$ of order three and $A^{-1}C$ of order two gives
$\operatorname{tr}(A)\operatorname{tr}(C)=1$, whereas the five-term locus
is $\operatorname{tr}(A)+\operatorname{tr}(C)=1$.  The two loci meet only
where $x^2+x+1=0$, namely in the order-three configuration excluded by
$p>5$.  Triangle-group and subgroup-classification results in $\PSL_2$
therefore address a transverse problem.

\begin{remark}[What was learned]
	The group-theoretic reformulation accurately classifies every possible
	ambient torus and rational conjugacy type.  It also shows why merely
	finding $p^j$-torsion in $\Sp_4$ is irrelevant.  The missing statement
	is still that a codimension-one trace condition avoids a particular
	$p$-Sylow subset; neither maximal-torus classification nor general
	character estimates encode the characteristic-two additive rigidity
	needed for that avoidance.
\end{remark}

\section{Characteristic zero: a complete classification}\label{sec:charzero}

Write $\omega=e^{2\pi i/3}$ and $\zeta_n=e^{2\pi i/n}$.

\begin{theorem}\label{thm:charzero}
	Let $u,v$ be complex roots of unity with
	\begin{equation}\label{eq:charzero}
		1+u+u^{-1}+v+v^{-1}=0 .
	\end{equation}
	Then, up to interchanging $u$ and $v$ and replacing either by its
	inverse, exactly one of the following holds.
	\begin{enumerate}
		\item $u=\zeta_5$ and $v=\zeta_5^{2}$; here
		      $\{1,u,u^{-1},v,v^{-1}\}=\mu_5$ and $Q_Y=\Phi_5$ with $Y=1$.
		\item $u=i$ and $v=\omega$; here $Q_Y=\Phi_3\Phi_4$ with $Y=2$.
		\item $u=-1$ and $v=-\omega$; here $Q_Y=\Phi_2^{2}\Phi_6$ with
		      $Y=0$.
	\end{enumerate}
	In particular the order pairs $(\ord{u},\ord{v})$ that occur are
	exactly $(5,5)$, $(4,3)$ and $(2,6)$.
\end{theorem}

\begin{proof}
	The left side of \eqref{eq:charzero} is a vanishing sum of five roots
	of unity counted with multiplicity.
	Decompose it into minimal vanishing subsums.
	By the classification of minimal vanishing sums of weight at most nine
	\cite[Theorem~6]{ConwayJones1976}, restated with the weight bookkeeping
	of \cite[Table~1]{PoonenRubinstein1998}, the possible weights of a
	minimal vanishing sum of roots of unity are
	$2,3,5,6,7,\ldots$; there is none of weight~$4$, the unique similarity
	class of weight~$2$ is $\mathcal{R}_2=1+(-1)$, the unique class of
	weight~$3$ is $\mathcal{R}_3=1+\omega+\omega^2$, and the unique class
	of weight~$5$ is $\mathcal{R}_5=1+\zeta_5+\zeta_5^2+\zeta_5^3+\zeta_5^4$.
	Therefore $5=5$ or $5=2+3$.

	Two degenerate configurations are excluded at once.
	If $u=1$ then \eqref{eq:charzero} reads $3+v+v^{-1}=0$, impossible
	since $\abs{v+v^{-1}}\leq2$.
	If $v=u^{\pm1}$ then \eqref{eq:charzero} reads $1+2(u+u^{-1})=0$, so
	$2\cos\vartheta=-\tfrac12$ for a rational multiple $\vartheta$ of
	$2\pi$; Niven's theorem \cite[Cor.~3.12]{Niven1956} says the only
	rational values of $2\cos\vartheta$ are $0,\pm1,\pm2$.

	\emph{Case $5=5$.}
	The multiset $\{1,u,u^{-1},v,v^{-1}\}$ is $\eta\cdot\mu_5$ for some
	root of unity $\eta$.
	It contains $1$, so $\eta^{-1}\in\mu_5$ and the multiset is $\mu_5$
	itself.
	Thus $\{u,u^{-1},v,v^{-1}\}=\{\zeta_5,\zeta_5^2,\zeta_5^3,\zeta_5^4\}$,
	which is alternative~(1).

	\emph{Case $5=2+3$.}
	Some two-element sub-multiset sums to zero, i.e.\ consists of $x$ and
	$-x$, and the complementary three-element sub-multiset is a rotation of
	$\mathcal{R}_3$.
	We enumerate the possible pairs.
	\begin{itemize}
		\item $\{u,u^{-1}\}$: then $u^{-1}=-u$, so $u^2=-1$ and
		      $u=\pm i$.
		      The remaining sum $1+v+v^{-1}=0$ forces
		      $\{1,v,v^{-1}\}=\{1,\omega,\omega^2\}$, i.e.\
		      $v=\omega^{\pm1}$.
		      This is alternative~(2).
		      The symmetric choice $\{v,v^{-1}\}$ gives the same
		      alternative with $u$ and $v$ interchanged.
		\item $\{1,u\}$: then $u=-1$, and the remaining multiset
		      $\{-1,v,v^{-1}\}$ is a rotation $\eta\{1,\omega,\omega^2\}$.
		      Containing $-1$ forces $\eta\in-\mu_3$, and taking
		      $\eta=-1$ gives $\{v,v^{-1}\}=\{-\omega,-\omega^2\}$,
		      which is consistent because $(-\omega)^{-1}=-\omega^2$.
		      This is alternative~(3); the choice $\{1,v\}$ is symmetric.
		\item $\{u,v\}$, $\{u,v^{-1}\}$, $\{u^{-1},v\}$ or
		      $\{u^{-1},v^{-1}\}$: then $v=\pm u^{\pm1}$ up to sign, and
		      the complementary triple is $\{1,w,-w\}$ for some root of
		      unity $w$, whose sum is $1\ne0$.
	\end{itemize}
	The identification of $Q_Y$ in each case is a direct expansion:
	$(T^2-\alpha T+1)(T^2-\alpha'T+1)$ with $\alpha=-(u+u^{-1})$,
	$\alpha'=-(v+v^{-1})$, $\alpha+\alpha'=1$ and $Y=2+\alpha\alpha'$.
\end{proof}

Theorem~\ref{thm:charzero} was verified by exhaustive search over all pairs
of roots of unity of order dividing $420=2^2\cdot3\cdot5\cdot7$; the only
order pairs found were $(5,5)$, $(3,4)$ and $(2,6)$.

\begin{corollary}\label{cor:charzero-pj}
	Let $p>5$ be prime, $j\geq1$, $\zeta=\zeta_{p^j}$.
	Then
	\[
		\Lambda_{a,b}:=1+\zeta^{a}+\zeta^{-a}+\zeta^{b}+\zeta^{-b}\ne0
		\qquad\text{for all }a,b\in\Z .
	\]
\end{corollary}

\begin{proof}
	By Theorem~\ref{thm:charzero} any vanishing instance has
	$\ord{\zeta^a}\in\{2,3,4,5,6\}$, which is impossible for a divisor of
	$p^j$ with $p>5$ unless $\zeta^a=1$; but the pattern $(u,v)=(1,\ast)$
	does not occur in the classification.
	Alternatively, and more directly, Lam and Leung's main theorem
	\cite[Main Theorem]{LamLeung2000} says that the weights of vanishing
	sums of $p^j$-th roots of unity form the numerical semigroup
	$\N p$; since $0<5<p$, the weight $5$ is not attained.
\end{proof}

\begin{remark}
	Corollary~\ref{cor:charzero-pj} is the reason Problem~\ref{prob:main}
	is not a statement about exact vanishing sums at all.
	The characteristic-zero equation has \emph{no} solutions whatsoever for
	$p>5$; the entire content of the problem is that
	$\Lambda_{a,b}$, a nonzero algebraic integer, might nevertheless be
	divisible by a prime above $2$.
\end{remark}

\section{Cyclotomic units: an unconditional family}\label{sec:units}

Corollary~\ref{cor:charzero-pj} says $\Lambda_{a,b}\ne0$.
If in addition $\Lambda_{a,b}$ is a \emph{unit}, then it is nonzero in every
residue field and \eqref{eq:main} is impossible in every characteristic.
This section determines exactly when that happens for structural reasons.

For $r\geq1$ put
\[
	G_r(x)=x^{2r}+x^{r+1}+x^{r}+x^{r-1}+1\in\Z[x],
	\qquad\text{so}\qquad
	G_r(x)=x^{r}\left(1+x+x^{-1}+x^{r}+x^{-r}\right).
\]

\begin{theorem}\label{thm:units}
	$G_r$ is a product of cyclotomic polynomials if and only if
	$r\in\{2,3\}$, in which case
	\[
		G_2=\Phi_5,
		\qquad
		G_3=\Phi_5\Phi_6 .
	\]
\end{theorem}

\begin{proof}
	The two displayed factorizations are direct expansions.
	Conversely suppose every root of $G_r$ is a root of unity.
	If $u$ is a root, then $u\ne0$ and $1+u+u^{-1}+u^{r}+u^{-r}=0$, so
	$(u,u^{r})$ is one of the pairs of Theorem~\ref{thm:charzero}.
	Alternative~(2) is impossible because $\ord{u}=4$ and $\ord{u^r}=3$
	cannot hold simultaneously with $u^{r}$ a power of $u$, and likewise
	with the roles reversed.
	Alternative~(1) forces $\ord{u}=5$, alternative~(3) forces
	$\ord{u}\in\{2,6\}$; and $u=-1$ is excluded because
	\[
		G_r(-1)=1+(-1)^{r+1}+(-1)^{r}+(-1)^{r-1}+1=2+(-1)^{r-1}\in\{1,3\}.
	\]
	Hence every root of $G_r$ has order $5$ or $6$, so
	$G_r=\Phi_5^{e_5}\Phi_6^{e_6}$ and $2r=4e_5+2e_6$.

	We show $e_5\leq1$ and $e_6\leq1$.
	Set $\Lambda(x)=1+x+x^{-1}+x^{r}+x^{-r}$, so that $G_r=x^{r}\Lambda$ and
	at a root of $G_r$ the derivative $G_r'$ vanishes if and only if
	$\Lambda'$ does.
	Now
	\[
		\Lambda'(x)=1-x^{-2}+r\left(x^{r-1}-x^{-r-1}\right).
	\]
	If $\ord{u}=5$ then, by alternative~(1), $u^{r}=u^{\pm2}$; replacing
	$u$ by $u^{-1}$ if necessary, $u^{r}=u^{2}$, so $u^{r-1}=u$ and
	$u^{-r-1}=u^{-3}=u^{2}$, giving
	$\Lambda'(u)=1-u^{3}+r(u-u^{2})$.
	As $1,u,u^2,u^3$ is a $\Q$-basis of $\Q(\zeta_5)$, this is nonzero.
	If $\ord{u}=6$ then $u^{r}=-1$ by alternative~(3), so
	$u^{r-1}=-u^{-1}=u^{-r-1}$ and the bracket vanishes, leaving
	$\Lambda'(u)=1-u^{-2}\ne0$ because $u^{2}\ne1$.
	Hence $G_r$ has no repeated roots, $e_5,e_6\leq1$, and
	$2r\leq4+2=6$.
	Finally $r=1$ is excluded because $G_1=2x^2+x+2$ is not monic.
\end{proof}

\begin{corollary}\label{cor:units}
	Let $p>5$ be prime and $j\geq1$, and let $a,b$ be integers with
	\[
		b\equiv\pm2a,\ \pm3a
		\quad\text{or}\quad
		a\equiv\pm2b,\ \pm3b
		\pmod{p^j}.
	\]
	Then $\Lambda_{a,b}$ is a unit of $\Z[\zeta_{p^j}]$ if
	$a\not\equiv0$, and $\Lambda_{a,b}=5$ if $a\equiv b\equiv0$.
	In the unit case, \eqref{eq:main} has no solution with these exponents
	in any residue characteristic. In the degenerate case $\Lambda_{a,b}=5$,
	it can vanish only in characteristic $5$. In particular, neither case
	gives a solution in characteristic $2$.
\end{corollary}

\begin{proof}
	Since $\Lambda_{a,b}=\Lambda_{b,a}=\Lambda_{a,-b}=\Lambda_{-a,b}$, it
	suffices to treat $b=2a$ and $b=3a$.
	Put $\xi=\zeta^{a}$, of exact order $p^{i}$ say.
	If $i=0$ then $\Lambda=5$.
	If $i\geq1$ then by Theorem~\ref{thm:units}
	\[
		\Lambda_{a,2a}=\xi^{-2}\Phi_5(\xi),
		\qquad
		\Lambda_{a,3a}=\xi^{-3}\Phi_5(\xi)\Phi_6(\xi).
	\]
	By Apostol's evaluation of resultants of cyclotomic polynomials
	\cite{ApostolResultants}, $\Res(\Phi_m,\Phi_n)=q^{\varphi(n)}$ when
	$m/n$ is a power of a prime $q$ and $\Res(\Phi_m,\Phi_n)=1$ otherwise.
	With $m=p^{i}\geq7$ and $n\in\{5,6\}$ the quotient $m/n$ is not an
	integer, so both resultants equal $1$.
	Hence $\Nm_{\Q(\xi)/\Q}\!\left(\Phi_5(\xi)\right)=\pm1$ and
	$\Nm_{\Q(\xi)/\Q}\!\left(\Phi_6(\xi)\right)=\pm1$, and $\Lambda_{a,b}$
	is a unit.
	A unit has odd norm, so its image in $\Z[\zeta_{p^j}]/\mathfrak{p}$ is
	nonzero for every prime $\mathfrak{p}$.
\end{proof}

\begin{remark}
	The two identities behind Corollary~\ref{cor:units} are the cyclotomic
	unit identities
	\[
		1+\xi+\xi^{-1}+\xi^{2}+\xi^{-2}=\xi^{-2}\,\frac{\xi^{5}-1}{\xi-1},
		\qquad
		1+\xi+\xi^{-1}+\xi^{3}+\xi^{-3}
		=\xi^{-3}\,\frac{\xi^{5}-1}{\xi-1}\cdot\frac{\xi^{6}-1}{\xi^{3}-1}\cdot\frac{\xi-1}{\xi^{2}-1}.
	\]
	Theorem~\ref{thm:units} says that these are the only two members of the
	family, so the cyclotomic unit method removes at most $8$ of the
	$\varphi(p^j)$ possible ratios $b/a$ --- exactly $8$ when the classes
	$\{\pm2^{\pm1}\}$ and $\{\pm3^{\pm1}\}$ are distinct modulo $p^j$,
	and only $4$ when they coincide, as they do modulo $7$ --- and cannot
	remove more.
\end{remark}

\begin{example}
	For $p^j=49$ the ratio classes $\{\pm r^{\pm1}\}$ that consist of units
	are $\{2,24,25,47\}$ and $\{3,16,33,46\}$, exactly as predicted; see
	Table~\ref{tab:res49}.
	The class $\{13,15,34,36\}$ also has relative norm $1$, but
	sporadically: $G_{13}=\Phi_5\cdot H$ with $H$ of degree $22$ and
	Mahler measure $\approx0.2507>0$, and $\Res(\Phi_{n},G_{13})\ne\pm1$
	for $n=13,17,23,29,31$.
\end{example}

\section{Reduction modulo a prime above two}\label{sec:reduction}

Fix $p>5$, $j\geq1$, and set
\[
	K=\Q(\zeta_{p^j}),\qquad
	K^{+}=\Q(\zeta_{p^j}+\zeta_{p^j}^{-1}),\qquad
	f=\ord[p^j]{2},\qquad
	g=\frac{\varphi(p^j)}{f}.
\]
Since $p$ is odd, $2$ is unramified in $K$ and splits into $g$ distinct
primes, each with residue field $\F_{2^{f}}$ and decomposition group
$\langle\sigma_2\rangle$, where $\sigma_t(\zeta)=\zeta^{t}$
\cite[Ch.~2]{Washington1997}.
Reduction modulo any prime $\mathfrak{p}\mid2$ restricts to a group
isomorphism $\mu_{p^j}\to\mu_{p^j}(\F_{2^{f}})$.

\begin{proposition}\label{prop:dictionary}
	Let $\Sigma_{a,b}=\{\mathfrak{p}\mid2:\ \mathfrak{p}\mid\Lambda_{a,b}\}$.
	Then:
	\begin{enumerate}
		\item \eqref{eq:main} holds for some $\theta$ of exact order $p^j$
		      and the given $(a,b)$ if and only if
		      $\Sigma_{a,b}\ne\emptyset$;
		\item $\Lambda_{a,b}\in\Z[\zeta_{p^j}+\zeta_{p^j}^{-1}]$, so
		      $\Nm_{K/\Q}(\Lambda_{a,b})=\Nm_{K^{+}/\Q}(\Lambda_{a,b})^{2}$;
		\item $v_2\!\left(\Nm_{K/\Q}(\Lambda_{a,b})\right)\geq f\cdot\abs{\Sigma_{a,b}}$,
		      with equality when every $v_{\mathfrak{p}}(\Lambda_{a,b})\leq1$;
		\item writing $r=b\,a^{-1}$ when $\gcd(a,p)=1$,
		      \[
			      \Nm_{K/\Q}(\Lambda_{a,b})=\pm\Res\!\left(\Phi_{p^j}(x),\,G_r(x)\right).
		      \]
	\end{enumerate}
	In particular, Problem~\ref{prob:main} has an affirmative answer for
	$(p,j)$ if and only if $\Res(\Phi_{p^j},G_r)$ is odd for every
	nonzero $r\in\Z/p^j\Z$.
	(A counterexample has $p\nmid a$ or $p\nmid b$; after interchanging
	$a$ and $b$ we may assume $\gcd(a,p)=1$, and then $r=ba^{-1}$ is a
	well-defined nonzero element of $\Z/p^j\Z$, not necessarily a unit.)
\end{proposition}

\begin{proof}
	(1) is the reduction isomorphism together with the fact that the $g$
	primes above $2$ are permuted transitively by $\Gal(K/\Q)$, and
	$\sigma_t(\Lambda_{a,b})=\Lambda_{ta,tb}$.
	(2) holds because the five summands are stable under
	$\zeta\mapsto\zeta^{-1}$, and $[K:K^{+}]=2$.
	(3) is the norm formula
	$v_2(\Nm_{K/\Q}(x))=\sum_{\mathfrak{p}\mid2}f\,v_{\mathfrak{p}}(x)$
	valid because $2$ is unramified with residue degree $f$ at every
	$\mathfrak{p}$.
	(4) follows from $\Lambda_{a,b}=\zeta^{-ra}G_r(\zeta^{a})$ and
	$\Nm_{K/\Q}(\zeta)=\pm1$, after the substitution
	$\zeta\mapsto\zeta^{a}$, which permutes the primitive $p^j$-th roots
	of unity; the fact that only the ratio $r$ matters is
	Lemma~\ref{lem:symmetry}.
\end{proof}

\subsection{Circular units and the exact local obstruction}

The parity question can be made slightly sharper by replacing the
five-term sum by products of circular units.  Write
$K_j=\Q(\zeta_j)$, $K_j^+=\Q(\zeta_j+\zeta_j^{-1})$ with
$\zeta_j=\zeta_{p^j}$, put $n=p^j$, and choose the integer
$h_j=(n+1)/2$, so that $2h_j\equiv1\pmod n$.  For
$a\in\Z/n\Z$ define
\[
	\epsilon_{j,a}=\zeta_j^{-h_ja}(1+\zeta_j^a).
\]
Also put
\[
	s_{j,a}=\zeta_j^{-h_ja}(1-\zeta_j^a).
\]
Thus $\epsilon_{j,-a}=\epsilon_{j,a}$ and
\[
	\epsilon_{j,a}=\zeta_j^{a/2}+\zeta_j^{-a/2},
	\qquad
	\epsilon_{j,a}^{\,2}=\lambda_a+2,
	\qquad
	\lambda_a=\zeta_j^a+\zeta_j^{-a}.
\]
If $a\not\equiv0\pmod n$ then $\epsilon_{j,a}$ is a circular unit: indeed
$1+\xi=(1-\xi^2)/(1-\xi)$ for $\xi=\zeta_j^a$ of odd order, and
$\Phi_{\operatorname{ord}(\xi)}(-1)=1$.
For such $a$, the auxiliary element $s_{j,a}$ is a $p$-unit, integral and a unit at
every prime above $2$, and for $\sigma_2(\zeta_j)=\zeta_j^2$ one has the
Hilbert--90 identity
\begin{equation}\label{eq:epsilon-h90}
	\frac{\sigma_2(s_{j,a})}{s_{j,a}}
	=
	\zeta_j^{-h_ja}(1+\zeta_j^a)
	=\epsilon_{j,a}.
\end{equation}
Thus each $\epsilon_{j,a}$ has local norm one for the unramified
decomposition extension at $2$, and is not merely a circular unit but a
specified local coboundary.

\begin{lemma}[norm compatibility]\label{lem:epsilon-norms}
	For $j>1$,
	\[
		\Nm_{K_j/K_{j-1}}(s_{j,a})=s_{j-1,\bar a},
		\qquad
		\Nm_{K_j/K_{j-1}}(\epsilon_{j,a})=\epsilon_{j-1,\bar a}
		\qquad (p\nmid a),
	\]
	where $\bar a$ denotes reduction modulo $p^{j-1}$, and
	\[
		s_{j,pa}=s_{j-1,a},
		\epsilon_{j,pa}=\epsilon_{j-1,a}.
	\]
\end{lemma}

\begin{proof}
	The conjugates over $K_{j-1}$ send
	$\zeta_j$ to $\zeta_j^{1+tp^{j-1}}$, $0\leq t<p$.  Hence
	\[
		\Nm_{K_j/K_{j-1}}(1+\zeta_j^a)=1+\zeta_{j-1}^a
	\]
	for $p\nmid a$, by evaluating $T^p-\zeta_{j-1}^a$ at $T=-1$; the
	same evaluation at $T=1$ gives
	$\Nm_{K_j/K_{j-1}}(1-\zeta_j^a)=1-\zeta_{j-1}^a$.
	Also $\Nm_{K_j/K_{j-1}}(\zeta_j)=\zeta_{j-1}$, and
	$h_j\equiv h_{j-1}\pmod {p^{j-1}}$, giving the first identity.
	The second follows at once from $\zeta_j^p=\zeta_{j-1}$ and the same
	congruence for $h_j$.
\end{proof}

\begin{proposition}[Hilbert--90 reformulation]\label{prop:epsilon-reform}
	Let $\mathfrak{P}\mid2$ be a prime of $K_j$ and let
	$a,b\in\Z/n\Z$ be nonzero.  Put
	$w_{j,a,b}=s_{j,a}s_{j,b}$.  Then the following are equivalent:
	\[
		\epsilon_{j,a}\epsilon_{j,b}\equiv1\pmod{\mathfrak{P}},
		\qquad
		w_{j,a,b}\equiv1\pmod{\mathfrak{P}},
	\]
	and
	\[
		1+\zeta_j^{(a+b)/2}+\zeta_j^{-(a+b)/2}
		 +\zeta_j^{(a-b)/2}+\zeta_j^{-(a-b)/2}
		\equiv0\pmod{\mathfrak{P}} .
	\]
	Consequently a conductor-$p^j$ counterexample to
	Problem~\ref{prob:main} is exactly a congruence
	$\epsilon_{j,a}\epsilon_{j,b}\equiv1$ modulo one prime above $2$,
	with at least one of $a,b$ prime to $p$.
\end{proposition}

\begin{proof}
	Since $\sigma_2$ lies in the decomposition group at every prime above
	$2$, reduction modulo $\mathfrak P$ sends $\sigma_2(x)$ to
	$\bar x^2$.  By \eqref{eq:epsilon-h90},
	\[
		\overline{\epsilon_{j,a}\epsilon_{j,b}}
		=\frac{\overline{w_{j,a,b}}^{\,2}}{\overline{w_{j,a,b}}}
		=\overline{w_{j,a,b}},
	\]
	because $w_{j,a,b}$ is a unit at $\mathfrak P$.  Hence
	$\epsilon_{j,a}\epsilon_{j,b}\equiv1$ if and only if
	$w_{j,a,b}\equiv1$.  Equivalently,
	$\sigma_2(w_{j,a,b})/w_{j,a,b}\equiv1$ forces the residue of
	$w_{j,a,b}$ to be Frobenius-fixed, hence to lie in
	$\F_2^\times=\{1\}$.

	Expanding the product gives
	\[
		\epsilon_{j,a}\epsilon_{j,b}
		=\lambda_{(a+b)/2}+\lambda_{(a-b)/2}.
	\]
	In the residue field of $\mathfrak{P}$, subtracting $1$ is the same as
	adding $1$, so the displayed equivalence is immediate.  The change of
	variables from the original exponents $A,B$ is
	$a=A+B$, $b=A-B$.  Since $2$ is invertible modulo $p^j$, $A$ and $B$
	are both divisible by $p$ if and only if $a$ and $b$ are both divisible
	by $p$.
\end{proof}

For the local cohomological interpretation, let
$L=K_{j,\mathfrak P}$, let $U=\mathcal O_L^\times$ and
$U^1=1+\mathfrak P$, and put $D=\langle\sigma_2\rangle$.
The coboundary map $\delta(u)=\sigma_2(u)/u$ induces on the residue
sequence $1\to U^1\to U\to\F_{2^f}^\times\to1$ the map
$\bar u\mapsto\bar u^2/\bar u=\bar u$.  Therefore
\[
	\delta^{-1}(U^1)=U^1,
\]
and $H^1(D,U)=H^1(D,U^1)=0$ in this unramified local situation.  Thus
Hilbert~90 merely says that the norm-one unit
$\epsilon_{j,a}\epsilon_{j,b}$ is a coboundary; the extra parity
condition is exactly that the chosen $p$-unit primitive
$w_{j,a,b}=s_{j,a}s_{j,b}$ already lies in the principal local unit group
$U^1$.  In ray-class language, if $\mathfrak l_j^+$ is the unique prime
of $K_j^+$ above $p$ and $v_p(a)=\alpha$, $v_p(b)=\beta$, then
\[
	(w_{j,a,b})=(\mathfrak l_j^+)^{(p^\alpha+p^\beta)/2}
	\qquad\text{in }K_j^+.
\]
Thus a primitive--primitive collision is a generator of
$\mathfrak l_j^+$ that is $1$ modulo one prime above $2$; a
primitive--lower collision is the analogous ray-class condition for
$(\mathfrak l_j^+)^{(1+p^\beta)/2}$ with $\beta\geq1$, with the
generator restricted to the sine-circular form $s_{j,a}s_{j,b}$.

This isolates exactly the two collision norms of
Theorem~\ref{thm:two-collision-norms}.  Let
\[
	P_j=(\Z/p^j\Z)^\times/\{\pm1\},
	\qquad
	L_j=\bigl(p\Z/p^j\Z\setminus\{0\}\bigr)/\{\pm1\}.
\]
Define
\[
	N_j^{\rm pp}=
	\prod_{\substack{\{a,b\}\subset P_j\\a\ne b}}
	\Nm_{K_j^+/\Q}\!\left(\epsilon_{j,a}^{\,2}\epsilon_{j,b}^{\,2}-1\right),
	\qquad
	N_j^{\rm pl}=
	\prod_{\substack{a\in P_j\\ b\in L_j}}
	\Nm_{K_j^+/\Q}\!\left(\epsilon_{j,a}^{\,2}\epsilon_{j,b}^{\,2}-1\right).
\]
Because the residue fields above $2$ have characteristic two,
$x^2-1\equiv0$ for a unit $x$ if and only if $x\equiv1$.  Therefore
Problem~\ref{prob:main} at level $p^j$ is equivalent to the simultaneous
oddness of $N_j^{\rm pp}$ and $N_j^{\rm pl}$.  In the notation
$\nu(\lambda,\mu)=(\lambda+2)(\mu+2)-1$ of the reciprocal-quartic note,
these are exactly the primitive--primitive and primitive--lower products,
since $\lambda_a+2=\epsilon_{j,a}^{\,2}$.

Now suppose $2\leq j\leq c_p$, so we are on the base-$2$ plateau.  Then
$\ord[p^j]{2}=\ord[p^{j-1}]{2}$, and every prime above $2$ in
$K_{j-1}$ splits completely in $K_j$.  If $\mathfrak{P}$ is a prime of
$K_{j-1}$ above $2$, then for $p\nmid a$,
\[
	\prod_{\mathfrak{Q}\mid\mathfrak{P}}
		\overline{s_{j,a}}^{\,(\mathfrak Q)}
	=
	\overline{s_{j-1,\bar a}}^{\,(\mathfrak P)},
	\qquad
	\prod_{\mathfrak{Q}\mid\mathfrak{P}}
		\overline{\epsilon_{j,a}}^{\,(\mathfrak Q)}
	=
	\overline{\epsilon_{j-1,\bar a}}^{\,(\mathfrak P)},
\]
and, for lower-level indices,
\[
	\prod_{\mathfrak{Q}\mid\mathfrak{P}}
		\overline{s_{j,pb}}^{\,(\mathfrak Q)}
	=
	\left(\overline{s_{j-1,b}}^{\,(\mathfrak P)}\right)^p,
	\qquad
	\prod_{\mathfrak{Q}\mid\mathfrak{P}}
		\overline{\epsilon_{j,pb}}^{\,(\mathfrak Q)}
	=
	\left(\overline{\epsilon_{j-1,b}}^{\,(\mathfrak P)}\right)^p .
\]
These identities are just Lemma~\ref{lem:epsilon-norms} reduced modulo
$\mathfrak{P}$.  Thus plateau splitting supplies only a product constraint
over the $p$ primes above $\mathfrak{P}$: if
$w_{j,a,b}\equiv1$ (equivalently
$\epsilon_{j,a}\epsilon_{j,b}\equiv1$) at one of them, the remaining
$p-1$ conjugate factors must multiply to the prescribed lower-level
$p$-unit.  There is no contradiction from splitting alone; a proof must
control the \emph{semi-local} reduction of these particular sine and
circular units.

Equivalently, let $\mathcal E_j^+$ be the $p$-unit group of $K_j^+$ and
let $\mathcal C_j^+$ be Sinnott's circular-unit group.  The semi-local
reduction map
\[
	\rho_j:\mathcal E_j^+\longrightarrow
	\prod_{\mathfrak q\mid2}
		(\mathcal O_{K_j^+}/\mathfrak q)^\times
\]
is the relevant finite quotient.  The desired parity theorem says precisely that
no sine $p$-unit $s_{j,a}s_{j,b}$, with $a\in P_j$ and
$b\in P_j\cup L_j$, lies in any one-prime principal-unit kernel
$\mathcal E_j^+\cap(1+\mathfrak q)$.  Applying the coboundary
$\delta=\sigma_2-1$ sends this finite set into
$\mathcal C_j^+$ as the circular units
$\epsilon_{j,a}\epsilon_{j,b}$; the sine lift is essential, because
$\delta^{-1}(1+\mathfrak q)=1+\mathfrak q$ only after using the
characteristic-two residue field.  The standard ray-class sequence
for modulus $2\mathcal O_{K_j^+}$ identifies
\[
	\left(\prod_{\mathfrak q\mid2}
		(\mathcal O_{K_j^+}/\mathfrak q)^\times\right)\big/
	\rho_j(\mathcal O_{K_j^+}^{\times})
\]
with the kernel of the ray class group modulo $2$ over the ordinary class
group.  Replacing full units by circular units gives a natural
surjection from the corresponding circular cokernel to this ray-class
kernel, with kernel a quotient of
$\mathcal O_{K_j^+}^{\times}/\mathcal C_j^+$ (whose size is governed by
Sinnott's formula), while the
Stickelberger--Jacobi-sum formulas describe character sums and ideal
classes.  Those tools do not by themselves decide the required pointwise
avoidance in the finite residue fields: for a chosen
$\mathfrak q$, the congruence is equivalent to
$\chi(\epsilon_{j,a})\chi(\epsilon_{j,b})=1$ for every multiplicative
character $\chi$ of the residue field, and the characters of order
prime to $p$ are needed as well as the $p$-primary Jacobi-sum data.

\subsection{What the reciprocal symmetry buys}

\begin{proposition}\label{prop:sigma}
	Let $\Sigma=\Sigma_{a,b}\ne\emptyset$, and assume $\gcd(a,p)=1$, which
	holds after possibly interchanging $a$ and $b$ in any counterexample to
	Problem~\ref{prob:main}.
	Let
	\[
		H=\left\{t\in(\Z/p^j\Z)^{\times}:\ \{\pm ta,\pm tb\}=\{\pm a,\pm b\}\right\}
	\]
	be the stabilizer of $\Lambda_{a,b}$ in $\Gal(K/\Q)$.
	Then $\Sigma$ is a union of $H$-orbits, $-1\in H$, and
	$H\subseteq\{t:\ t^{4}=1\}$, so $\abs{H}\leq4$.
	Consequently
	\[
		\abs{\Sigma}\ \geq\ \frac{\abs{H}}{\abs{H\cap\langle2\rangle}},
	\]
	and in particular $\abs{\Sigma}\geq2$ whenever
	$-1\notin\langle2\rangle$ in $(\Z/p^j\Z)^{\times}$, i.e.\ whenever
	$\rho_2(p^j)=\ord[p^j]{2}$.
\end{proposition}

\begin{proof}
	If $t\in H$ then $\sigma_t(\Lambda_{a,b})=\Lambda_{a,b}$, so
	$\mathfrak{p}\mid\Lambda_{a,b}$ implies
	$\sigma_t(\mathfrak{p})\mid\Lambda_{a,b}$; thus $\Sigma$ is
	$H$-stable.
	Clearly $-1\in H$.
	Let $t\in H$.
	If $ta\equiv\pm a$, then $t\equiv\pm1$ because $a$ is invertible.
	Otherwise $ta\equiv\pm b$ and $tb\equiv\pm a$; the second relation
	shows $b$ is invertible, and multiplying the two gives
	$t^{2}ab\equiv\pm ab$, hence $t^{2}\equiv\pm1$ and $t^{4}\equiv1$.
	Since $(\Z/p^j\Z)^{\times}$ is cyclic,
	$\abs{\{t:t^4=1\}}=\gcd(4,p-1)\leq4$.
	The orbit of $\mathfrak{p}$ under $H$ has size
	$\abs{H}/\abs{H\cap D_{\mathfrak{p}}}$ with
	$D_{\mathfrak{p}}=\langle2\rangle$.
	If $-1\notin\langle2\rangle$ then $-1\in H\setminus\langle2\rangle$, so
	$\abs{H}/\abs{H\cap\langle2\rangle}\geq2$.
\end{proof}

\begin{example}
	For $p=17$, $j=1$, one has $\ord[17]{2}=8$, $\varphi(17)=16$, $g=2$
	and $-1=2^{4}\in\langle2\rangle$.
	The ratio class of $r=13$ gives
	$\Res(\Phi_{17},G_{13})=256=2^{8}$, i.e.\
	$\abs{\Sigma}=1$; the relation
	$1+\theta^{3}+\theta^{-3}+\theta^{5}+\theta^{-5}=0$ indeed holds in
	$\F_{256}$ for a suitable $\theta$ of order $17$, and $\abs{S_{17}}=16$,
	matching $d(16)=8$.

	For $p=31$, $j=1$, one has $\ord[31]{2}=5$, $\varphi(31)=30$, $g=6$,
	and $-1\notin\langle2\rangle=\{1,2,4,8,16\}$.
	The class of $r=5$ gives $\Res(\Phi_{31},G_5)=1024=2^{10}=2^{5\cdot2}$,
	i.e.\ $\abs{\Sigma}=2$, in exact agreement with
	Proposition~\ref{prop:sigma}.
	Here $\abs{S_{31}}=40$ and $d(30)=20$.
\end{example}

\section{Sine residues, ray classes and power residues}\label{sec:power}

The reformulations of Section~\ref{sec:reduction} say that a counterexample
is a congruence $s_{j,a}s_{j,b}\equiv1$ at one prime above $2$.
This section carries that out concretely.
We compute the residues of the sine $p$-units, prove that the residue map is
injective on $\pm$-classes, and use this to settle unconditionally every
ratio lying in the decomposition group at $2$
(Theorem~\ref{thm:frobratio}).
We then make the ray-class formulation explicit as a family of power-residue
conditions, and determine the $p$-primary member of that family --- the
Stickelberger--Jacobi part --- completely
(Theorem~\ref{thm:pstructure}).
The outcome is a genuine obstruction, strong enough to prove
$d(p^{j}-1)=0$ for every prime in an explicit family, conditional only
on one power-residue nonvanishing that is verified in every computed case
(Theorem~\ref{thm:qr}).  Section~\ref{sec:traceobs} then removes that
hypothesis by an additive (Artin--Schreier) argument
(Theorem~\ref{thm:trace}, Corollary~\ref{cor:trace-family}).
The obstruction is vacuous when $\ord[p]{2}$ is
even, and for odd order its descent fails on the Wieferich plateau
(Proposition~\ref{prop:plateau-degenerate}).

Throughout this section $p>5$, $n=p^{j}$, $f=\ord[n]{2}$, $q=2^{f}$,
$\mathfrak P\mid2$ is a fixed prime of $K_j=\Q(\zeta_n)$, and
$\theta=\zeta_n\bmod\mathfrak P$ has exact order $n$ in $\F_q^\times$.
For $x\in\Z/n\Z$ we write
\[
	\bar\lambda_x=\theta^{x}+\theta^{-x}\in\F_{q},
\]
the reduction of $\lambda_x=\zeta_n^{x}+\zeta_n^{-x}$.

\subsection{The residue of a sine unit is a trace}

\begin{lemma}\label{lem:sine-trace}
	Let $h=(n+1)/2$, so $2h\equiv1\pmod n$.  For every $a\in\Z/n\Z$ with
	$a\not\equiv0$,
	\[
		\overline{s_{j,a}}=\overline{\epsilon_{j,a}}=\bar\lambda_{ha}.
	\]
	Moreover, for all $x,y\in\Z/n\Z$,
	\begin{equation}\label{eq:lam-rules}
		\bar\lambda_{2x}=\bar\lambda_x^{\,2},
		\qquad
		\bar\lambda_{-x}=\bar\lambda_{x},
		\qquad
		\bar\lambda_x\bar\lambda_y=\bar\lambda_{x+y}+\bar\lambda_{x-y}.
	\end{equation}
\end{lemma}

\begin{proof}
	From $2h\equiv1$ we get $1-h\equiv h$, hence
	$\overline{s_{j,a}}=\theta^{-ha}(1+\theta^{a})
	=\theta^{-ha}+\theta^{a-ha}=\theta^{-ha}+\theta^{ha}$.
	Reducing \eqref{eq:epsilon-h90} modulo $\mathfrak P$ and using that
	$\sigma_2$ acts as the Frobenius on the residue field gives
	$\overline{\epsilon_{j,a}}=\overline{s_{j,a}}^{\,2}/\overline{s_{j,a}}
	=\overline{s_{j,a}}$, because $s_{j,a}$ is a unit at $\mathfrak P$.
	The three rules in \eqref{eq:lam-rules} are the Frobenius, the
	substitution $x\mapsto-x$, and multiplying out.
\end{proof}

\begin{lemma}[fibre lemma]\label{lem:fibre}
	For $x,y\in\Z/n\Z$:
	\begin{enumerate}
		\item $\bar\lambda_x=0$ if and only if $x\equiv0$;
		\item $\bar\lambda_x\ne1$ for every $x$;
		\item $\bar\lambda_x=\bar\lambda_y$ if and only if
		      $y\equiv\pm x\pmod n$.
	\end{enumerate}
	Consequently $\F_2(\bar\lambda_x)=\F_{2^{d}}$, where $d$ is the order
	of $2$ in $(\Z/m\Z)^{\times}/\{\pm1\}$ and $m=n/\gcd(x,n)$ is the exact
	order of $\theta^{x}$.
\end{lemma}

\begin{proof}
	In characteristic two
	$(T-\theta^{x})(T-\theta^{-x})=T^{2}+\bar\lambda_xT+1$, so
	$\bar\lambda_x$ determines the unordered pair
	$\{\theta^{x},\theta^{-x}\}$, and (3) follows because reduction is
	injective on $\mu_n$.
	For (1), $\bar\lambda_x=0$ forces $\theta^{2x}=1$, hence $x\equiv0$
	because $n$ is odd.
	For (2), $\bar\lambda_x=1$ forces $\theta^{2x}+\theta^{x}+1=0$, so
	$\theta^{x}$ has order $3$, impossible for $p>3$.
	The last statement is (3) applied to $y=2^{k}x$: one has
	$\bar\lambda_x^{\,2^{k}}=\bar\lambda_{2^{k}x}=\bar\lambda_x$ if and only
	if $2^{k}\equiv\pm1\pmod m$.
\end{proof}

\begin{proposition}[product dictionary]\label{prop:product-form}
	Put $x=h(A+B)$ and $y=h(A-B)$, a bijection of $(\Z/n\Z)^2$ onto itself
	with inverse $A=x+y$, $B=x-y$.  Then
	\[
		\Lambda_{A,B}\equiv0\pmod{\mathfrak P}
		\quad\Longleftrightarrow\quad
		\bar\lambda_x\bar\lambda_y=1 ,
	\]
	the excluded degenerate cases $A\equiv0$, $B\equiv0$, $B\equiv\pm A$
	correspond to $y\equiv-x$, $y\equiv x$, $xy\equiv0$, and $p\mid A,B$
	if and only if $p\mid x,y$.
	The ratios correspond by the Möbius involution
	\[
		r_{\rm sum}=B/A\ \longleftrightarrow\
		r_{\rm prod}=y/x=\frac{1-r_{\rm sum}}{1+r_{\rm sum}} .
	\]
	In particular Problem~\ref{prob:main} at level $p^{j}$ asserts that
	there is no pair $x,y\in\Z/n\Z$, not both divisible by $p$, with
	$\bar\lambda_x\bar\lambda_y=1$.
\end{proposition}

\begin{proof}
	Immediate from \eqref{eq:lam-rules}: $\bar\lambda_x\bar\lambda_y
	=\bar\lambda_{x+y}+\bar\lambda_{x-y}=\bar\lambda_A+\bar\lambda_B$,
	and $\Lambda_{A,B}=1+\lambda_A+\lambda_B$.
	The statement on ratios is the computation
	$y/x=(1-B/A)/(1+B/A)$; $h$ is invertible, so $p\mid x,y$ iff
	$p\mid A,B$.
\end{proof}

We record two exact identities among the residues, both of which are
reductions of norm relations and both of which will be used below.

\begin{lemma}\label{lem:lam-relations}
	\[
		\prod_{t\in(\Z/n\Z)^{\times}}\bar\lambda_t=1,
		\qquad
		\prod_{t\in H_\beta}\bar\lambda_{xt}=\bar\lambda_{p^{\beta}x}
		\quad(1\leq\beta\leq j),
	\]
	where $H_\beta=1+p^{\,j-\beta}\Z/p^{j}\Z$ is the subgroup of order
	$p^{\beta}$ of $(\Z/n\Z)^{\times}$ and $\gcd(x,p)=1$.
\end{lemma}

\begin{proof}
	Write $\bar\lambda_t=\theta^{-t}(1+\theta^{2t})$.  Pairing $t$ with
	$n-t$ gives $\sum_{t}t\equiv0\pmod n$, while
	$\prod_{t}(1+\theta^{2t})=\prod_{u}(1+\theta^{u})
	=\overline{\Phi_{n}(-1)}=1$.  This is the first identity.
	For the second, $2xt$ runs over $2x+p^{\,j-\beta}\Z/p^{j}\Z$ as $t$
	runs over $H_\beta$, so
	$\prod_{t}(1+\theta^{2xt})=\prod_{w\bmod p^{\beta}}
	(1+\theta^{2x}\eta^{w})=1+\theta^{2p^{\beta}x}$ with
	$\eta=\theta^{p^{\,j-\beta}}$ of order $p^{\beta}$;
	and $\sum_{t\in H_\beta}xt\equiv p^{\beta}x\pmod{p^{j}}$.
\end{proof}

\subsection{The semi-local residue module and a no-go theorem}
\label{sec:module}

Put
\[
	\bar G=(\Z/n\Z)^\times/\{\pm1\},\qquad
	\bar f=\rho_2(n),\qquad
	\bar D=\langle2\rangle\subseteq\bar G,\qquad
	g^+=\frac{\varphi(n)}{2\bar f},
\]
and $E=\F_{2^{\bar f}}$.  The $g^+$ primes of $K_j^+$ above two are
permuted simply transitively by $\bar G/\bar D$, and their residue fields
are all $E$.  Let
\[
	M=\prod_{\mathfrak q\mid2}
		(\mathcal O_{K_j^+}/\mathfrak q)^\times.
\]

\begin{theorem}[Semi-local module]\label{thm:semilocal-module}
	There is an isomorphism of $\bar G$-modules
	\[
		M\cong\Z[\bar G]/(\sigma_2-2)\Z[\bar G].
	\]
	Reduction at one chosen prime is evaluation on the decomposition
	group: it kills the cosets outside $\bar D$ and sends
	$\sigma_2^i$ to multiplication by $2^i$ in
	$\Z/(2^{\bar f}-1)\Z$.
\end{theorem}

\begin{proof}
	As a permutation module on the primes above two,
	\[
		M\cong\operatorname{Ind}_{\bar D}^{\bar G}E^\times.
	\]
	Choose a generator $\gamma$ of $E^\times$.  The map
	\[
		\Z[\bar D]\longrightarrow E^\times,\qquad
		\sum_i a_i\sigma_2^i\longmapsto
		\gamma^{\sum_i a_i2^i}
	\]
	is surjective with kernel generated by $\sigma_2-2$, since
	\[
		\Z[X]/(X^{\bar f}-1,X-2)
		\cong\Z/(2^{\bar f}-1)\Z.
	\]
	Inducing to $\bar G$ proves the first assertion.  Projection to the
	trivial coset gives the stated one-prime evaluation map.
\end{proof}

The structural relations satisfied by the true residues are the
Frobenius and sign rules in \eqref{eq:lam-rules}, the distribution
relations for circular units \cite{Sinnott1978,Sinnott1980,KubertLang1981},
and their norm consequence.  Abstract them as follows.

\begin{definition}\label{def:admissible-system}
	An \emph{admissible system} is a map
	\[
		m:(\Z/n\Z)\setminus\{0\}\longrightarrow E^\times
	\]
	satisfying
	\begin{align*}
		m(2x)&=m(x)^2,\\
		m(-x)&=m(x),\\
		\prod_{k=0}^{p-1}m(x+kn/p)&=m(px)
		\qquad(px\not\equiv0\pmod n).
	\end{align*}
\end{definition}

The true system $m(x)=\bar\lambda_x$ is admissible by
Lemma~\ref{lem:sine-trace} and Lemma~\ref{lem:lam-relations}.

\begin{theorem}[Freeness]\label{thm:admissible-freeness}
	Choose one representative $x_i$ from each of the $g^+$ orbits of
	$\langle2,-1\rangle$ on the primitive indices.  Evaluation induces a
	bijection
	\[
		\{\text{admissible systems}\}
		\longrightarrow(E^\times)^{g^+},
		\qquad
		m\longmapsto(m(x_1),\ldots,m(x_{g^+})).
	\]
	Thus the primitive-orbit values may be prescribed arbitrarily and
	independently.
\end{theorem}

\begin{proof}
	Frobenius and sign determine $m$ on each primitive orbit from
	$m(x_i)$.  Every nonzero lower-level index is $pz$, and the
	distribution relation determines $m(pz)$ from values one level higher;
	induction gives injectivity.

	Conversely prescribe arbitrary $v_i\in E^\times$ and set
	$m(\pm2^kx_i)=v_i^{2^k}$.  This is well defined because each orbit has
	length $\bar f$ and $v_i^{2^{\bar f}}=v_i$.  Define the lower levels
	recursively by the distribution product.  Changing a lift of $pz$
	permutes the $p$ factors, and doubling or negating the lift respectively
	squares or permutes the same product.  Hence all three axioms hold.
\end{proof}

\begin{corollary}[Multiplicative no-go]\label{cor:multiplicative-nogo}
	If $g^+\geq2$ and $x,y$ are primitive indices in distinct
	$\langle2,-1\rangle$-orbits, then exactly
	$(2^{\bar f}-1)^{g^+-1}$ admissible systems satisfy
	$m(x)m(y)=1$.  Therefore the Frobenius, sign, distribution, and norm
	consequences encoded above cannot by themselves exclude a primitive
	collision.
\end{corollary}

\begin{proof}
	In the coordinates of Theorem~\ref{thm:admissible-freeness}, prescribe
	$v_2=v_1^{-1}$ and choose all remaining coordinates freely.  The norm
	relation is automatic: on each orbit its product is
	$v_i^{2^{\bar f}-1}=1$.
\end{proof}

On a Wieferich plateau $g^+\geq p^{j-1}$, so this no-go theorem applies
precisely where the universal proof is missing.  The script
\path{code/ai_prototypes/circular_module.py} constructs explicit systems
for $(p,j)=(1093,2)$ and $(3511,2)$ that also remain injective on
$\{\pm x\}$-classes, avoid the value $1$, and contain a primitive
collision.  The strengthened construction in
\path{code/ai_prototypes/plateau_endpoint.py} additionally imposes the
additive distribution law, the Artin--Schreier trace law, and the global
trace law of Theorem~\ref{thm:global-trace} below.  These are formal
countermodels to the listed axioms, not counterexamples to the original
problem: they deliberately omit the full multiplication identity
\[
	\bar\lambda_x\bar\lambda_y
	=\bar\lambda_{x+y}+\bar\lambda_{x-y}.
\]
They show exactly that any successful proof in this framework must use
more of that additive subgroup structure.

\subsubsection{Relation weights and Zech logarithms}

A positive relation of weight $w$ is a multiset of nonzero
$\{\pm x\}$-classes with product of the corresponding
$\bar\lambda_x$ equal to $1$.  The main problem is the assertion that
there is no primitive relation of weight two.  Within one Frobenius orbit
this invariant is explicit.  If $d_x$ is the multiplicative order of
$\bar\lambda_x$, then a relation
\[
	\prod_k\bar\lambda_{2^kx}^{e_k}=1
\]
is equivalent to
\[
	\sum_ke_k2^k\equiv0\pmod{d_x}.
\]
Thus an orbit supports a weight-two relation exactly when
$2^i\equiv-1\pmod{d_x}$ for some $i$.  This recovers the exceptional
decomposition-group ratios but says nothing about collisions between
different primitive orbits, which are exactly the free coordinates in
Theorem~\ref{thm:admissible-freeness}.

Equivalently, if $\gamma$ generates $\F_{2^f}^\times$,
$\theta=\gamma^L$, and $Z$ is the Zech logarithm
$1+\gamma^m=\gamma^{Z(m)}$, then
\[
	\log_\gamma\bar\lambda_x=Z(2xL)-xL,
\]
and a collision is
\[
	Z(2xL)+Z(2yL)\equiv(x+y)L\pmod{2^f-1}.
\]
This is an exact restatement, but available Costas/Sidon properties of
Zech logarithms do not control the restriction to two independent
$p^j$-orbits.  The freeness theorem explains why orbitwise identities
alone cannot do so.

\subsection{Ratios in the decomposition group}

The decomposition group at $\mathfrak P$ is $D=\langle2\rangle\subseteq
(\Z/n\Z)^{\times}$, and by \eqref{eq:lam-rules} a ratio $y/x\in D$ is
exactly the case in which the two sine residues are Frobenius conjugate.
This case can be settled completely and unconditionally.

\begin{theorem}\label{thm:frobratio}
	Let $x,y\in\Z/n\Z$ be nonzero with $y\equiv2^{i}x\pmod n$ for some
	$i\geq0$, and suppose $\bar\lambda_x\bar\lambda_y=1$ and
	$y\not\equiv\pm x$.
	Let $m$ be the exact order of $\theta^{x}$ and $f_m=\ord[m]{2}$.
	Then
	\[
		4\mid f_m,
		\qquad
		2^{f_m/2}\equiv-1\pmod m,
		\qquad
		y\equiv\pm\rho^{\pm1}x\ \ \text{with}\ \ \rho=2^{f_m/4},
		\ \rho^{2}\equiv-1\pmod m ,
	\]
	and, writing $\F=\F_{2^{f_m/2}}$ and $\F_0=\F_{2^{f_m/4}}$,
	\[
		\bar\lambda_x\in\F^{\times},
		\qquad
		\Nm_{\F/\F_0}\!\left(\bar\lambda_x\right)=1 .
	\]
	Equivalently, in the sum form of Proposition~\ref{prop:product-form},
	$B\equiv\mp\rho A$ and $\Tr_{\F/\F_0}(\bar\lambda_A)=1$.
\end{theorem}

\begin{proof}
	Since $y=2^{i}x$, both indices have the same $p$-valuation, so
	replacing $\theta$ by $\theta^{x}$ we may assume $x=1$ and $n=m$; write
	$f=f_m$.  (This is legitimate: $x(2^{i}\mp1)\equiv0\pmod n$ if and only
	if $2^{i}\equiv\pm1\pmod m$, so the excluded case is the same one, and
	the conclusion $2^{i}\equiv\rho^{\pm1}\pmod m$ gives back
	$y\equiv\rho^{\pm1}x\pmod n$.)
	By \eqref{eq:lam-rules} the hypothesis reads
	$\bar\lambda_1^{\,2^{i}+1}=1$, so the multiplicative order of
	$\bar\lambda_1$ divides
	\[
		\gcd\!\left(2^{i}+1,\,2^{f}-1\right)=
		\begin{cases}
			2^{d}+1, & f/d\ \text{even},\\
			1,       & f/d\ \text{odd},
		\end{cases}
		\qquad d=\gcd(i,f),
	\]
	a standard identity.
	If $f/d$ is odd then $\bar\lambda_1=1$, contradicting
	Lemma~\ref{lem:fibre}(2).
	So $f/d$ is even and $\bar\lambda_1\in\mu_{2^{d}+1}\subseteq
	\F_{2^{2d}}^{\times}$.
	Then $\bar\lambda_{2^{2d}}=\bar\lambda_1^{\,2^{2d}}=\bar\lambda_1$, so
	$2^{2d}\equiv\pm1\pmod n$ by Lemma~\ref{lem:fibre}(3).

	Suppose $2^{2d}\equiv1$.  Then $f\mid2d$; since $d\mid f$ and $f/d$ is
	even, $f=2d$.  Now $2^{d}$ is a square root of $1$ in the cyclic group
	$(\Z/n\Z)^{\times}$ and $2^{d}\ne1$, so $2^{d}\equiv-1$.  As
	$d=\gcd(i,f)$ divides $i$, we get $2^{i}\in\{2^{d},2^{2d}\}=\{-1,1\}$,
	i.e.\ $y\equiv\pm x$, which was excluded.

	Hence $2^{2d}\equiv-1$.  Then $f\mid4d$ but $f\nmid2d$, and $f/d$ is
	even, so $f=4d$; thus $\rho:=2^{d}=2^{f/4}$ satisfies
	$\rho^{2}\equiv-1$ and $2^{f/2}\equiv-1$.
	From $\gcd(i,4d)=d$ we get $i=kd$ with $k$ odd, so
	$2^{i}=\rho^{k}=\rho^{\pm1}$.
	Finally $\bar\lambda_1\in\mu_{2^{d}+1}\subseteq\F_{2^{2d}}^{\times}
	=\F^{\times}$ and $\bar\lambda_1^{\,2^{d}+1}=1$ is exactly
	$\Nm_{\F/\F_0}(\bar\lambda_1)=1$.

	For the sum form, Proposition~\ref{prop:product-form} with $y=\rho x$
	gives $A=(1+\rho)x$, $B=(1-\rho)x$ and
	$B/A=(1-\rho)/(1+\rho)=-\rho$ because $(1+\rho)^{2}=2\rho$ and
	$(1+\rho)(1-\rho)=2$.
	Since $2^{f/2}\equiv-1$ we have
	$\bar\lambda_A^{\,2^{f/2}}=\bar\lambda_{-A}=\bar\lambda_A$, so
	$\bar\lambda_A\in\F$, and the relation
	$\bar\lambda_A+\bar\lambda_B=\bar\lambda_A+\bar\lambda_A^{\,2^{f/4}}=1$
	is $\Tr_{\F/\F_0}(\bar\lambda_A)=1$.
\end{proof}

\begin{corollary}\label{cor:frobratio}
	If $-1\notin\langle2\rangle$ in $(\Z/p^{j}\Z)^{\times}$, or if
	$4\nmid\ord[p^j]{2}$, then no solution of \eqref{eq:main} has ratio
	$y/x$ in the decomposition group $\langle2\rangle$.
	In general the only ratio class in $\langle2\rangle$ that can occur is
	$\{\pm\rho^{\pm1}\}$ with $\rho=2^{f/4}$ a primitive fourth root of
	unity modulo $p^{j}$, and in that case the stabilizer of
	Proposition~\ref{prop:sigma} is $H=\{\pm1,\pm\rho\}\subseteq
	\langle2\rangle$, so its orbits on the primes above $2$ are singletons
	and the reciprocal symmetry gives no lower bound on
	$\abs{\Sigma_{a,b}}$ beyond $1$.
	(For $p=17$ one indeed has $\abs{\Sigma}=1$.)
\end{corollary}

Theorem~\ref{thm:frobratio} disposes of the $f$ ratios forming the
decomposition group, against the $8$ ratios that cyclotomic units dispose of
(Theorem~\ref{thm:units}); the two families always share
$r_{\rm prod}=2^{\pm1}$ and generically nothing else.
It is sharp: the exceptional class really occurs.

\begin{example}\label{ex:17-257}
	For $p=17$, $f=8$, $\rho=2^{2}=4$ and $\rho^{2}=16=-1$; all eight
	product-form solutions have ratio in $\{4,13\}=\{\pm\rho^{\pm1}\}$, and
	in fact $\abs{\Sigma}=1$, matching the example after
	Proposition~\ref{prop:sigma}.
	Moreover the eight sum-form indices of those solutions are exactly the
	eight $A$ with $\Tr_{\F_{16}/\F_{4}}(\bar\lambda_A)=1$, as
	Theorem~\ref{thm:frobratio} predicts.
	For $p=257$, $f=16$, $\rho=2^{4}=16$, and of the $144$ product-form
	solutions exactly $16$ have ratio in $\langle2\rangle$; all of these
	have ratio $\pm\rho^{\pm1}$.
	For $p=31$ and $p=127$ one has $-1\notin\langle2\rangle$, and indeed
	none of the $20$, resp.\ $56$, solutions has ratio in $\langle2\rangle$.
	For the Wieferich prime $p=1093$ and $j=2$ the exceptional class is
	present ($f=364$, $\rho=2^{91}$), so Theorem~\ref{thm:frobratio}
	reduces the $364$ ratios of the decomposition group to a single class,
	on which the condition becomes
	$\Tr_{\F_{2^{182}}/\F_{2^{91}}}(\bar\lambda_A)=1$.
	For the Wieferich prime $p=3511$, by contrast, $f=1755$ is odd, so
	Corollary~\ref{cor:frobratio} excludes \emph{every} ratio in the
	decomposition group, at every level $j$ of the plateau.
\end{example}

\subsection{The ray class of the prime above $p$ as a power-residue
condition}

We now make the ray-class statement of Section~\ref{sec:reduction}
explicit.  Recall $\mathfrak l_j^{+}$, the unique prime of $K_j^{+}$ above
$p$, and put $\pi_0=(1-\zeta_j)(1-\zeta_j^{-1})$, a generator of
$\mathfrak l_j^{+}$, and $\xi_a=s_{j,a}/s_{j,1}\in\mathcal C_j^{+}$ for
$p\nmid a$, the standard circular units.

\begin{proposition}\label{prop:rayclass}
	Let $\mathfrak q\mid2$ in $K_j^{+}$ and let $a,b$ be prime to $p$.
	Then $s_{j,a}s_{j,b}$ generates $\mathfrak l_j^{+}$ and
	\[
		s_{j,a}s_{j,b}\equiv1\pmod{\mathfrak q}
		\iff
		\bar\pi_0=\left(\bar\xi_a\bar\xi_b\right)^{-1}
		\ \text{in}\ (\mathcal O_{K_j^{+}}/\mathfrak q)^{\times}.
	\]
	In particular a primitive--primitive collision forces the class of
	$\mathfrak l_j^{+}$ in the ray class group of $K_j^{+}$ of modulus
	$\mathfrak q$ to be trivial, i.e.\ forces $\mathfrak l_j^{+}$ to split
	completely in the ray class field of conductor $\mathfrak q$; and it
	forces $\bar\pi_0\in\rho_{\mathfrak q}(\mathcal C_j^{+})$, hence
	$\bar\pi_0\in\rho_{\mathfrak q}(\mathcal O_{K_j^{+}}^{\times})$.
\end{proposition}

\begin{proof}
	$(s_{j,a})=\mathfrak l_j$ in $K_j$ for $p\nmid a$, so
	$(s_{j,a}s_{j,b})=\mathfrak l_j^{2}=\mathfrak l_j^{+}\mathcal O_{K_j}$,
	and $s_{j,a}s_{j,b}$ is fixed by complex conjugation because
	$s_{j,-a}=-s_{j,a}$.
	Also $s_{j,1}^{2}=-\pi_0$, so
	$s_{j,a}s_{j,b}=-\pi_0\xi_a\xi_b$.  Reducing modulo $\mathfrak q$ and
	using $-1\equiv1$ gives the displayed equivalence.
	The ray-class statement is the exact sequence
	$1\to(\mathcal O/\mathfrak q)^{\times}/
	\rho_{\mathfrak q}(\mathcal O^{\times})\to Cl_{\mathfrak q}\to Cl\to1$
	together with the fact that the class of the principal ideal
	$(\pi_0)$ in $Cl_{\mathfrak q}$ is the class of $\bar\pi_0$.
\end{proof}

Because $(\mathcal O_{K_j^{+}}/\mathfrak q)^{\times}$ is cyclic of odd
order, a membership statement in it is equivalent to the vanishing of
finitely many power-residue symbols.  This is the concrete form of the
``semi-local units modulo circular units'' condition.

\begin{definition}\label{def:psiell}
	For a prime $\ell\mid2^{f}-1$ let
	$\chi_\ell:\F_q^{\times}\to\mu_{\ell^{C_\ell}}$,
	$\chi_\ell(z)=z^{(q-1)/\ell^{C_\ell}}$, with
	$C_\ell=v_\ell(q-1)$, be the $\ell$-primary power-residue symbol, fix a
	generator $\xi$ of $\mu_{\ell^{C_\ell}}$ and define
	$\Psi_\ell:\Z/n\Z\setminus\{0\}\to\Z/\ell^{C_\ell}\Z$ by
	$\chi_\ell(\bar\lambda_x)=\xi^{\Psi_\ell(x)}$.
\end{definition}

\begin{lemma}\label{lem:psi-rules}
	$\Psi_\ell(2x)=2\Psi_\ell(x)$, $\Psi_\ell(-x)=\Psi_\ell(x)$,
	$\sum_{t\in(\Z/n\Z)^{\times}}\Psi_\ell(t)=0$, and
	$\sum_{t\in H_\beta}\Psi_\ell(xt)=\Psi_\ell(p^{\beta}x)$.
	A solution of \eqref{eq:main} with product-form indices $x,y$ forces
	\begin{equation}\label{eq:psi-collision}
		\Psi_\ell(x)+\Psi_\ell(y)\equiv0
		\pmod{\ell^{C_\ell}}\qquad\text{for every prime }\ell\mid2^{f}-1 ,
	\end{equation}
	and conversely \eqref{eq:psi-collision} for all $\ell$ is equivalent to
	$\bar\lambda_x\bar\lambda_y=1$.
\end{lemma}

\begin{proof}
	The first two rules are \eqref{eq:lam-rules}, the last two are
	Lemma~\ref{lem:lam-relations}, and \eqref{eq:psi-collision} is
	$\chi_\ell$ applied to $\bar\lambda_x\bar\lambda_y=1$.
	For the converse, $\F_q^{\times}$ is cyclic of order $q-1$ and the
	$\chi_\ell$ jointly separate its elements.
\end{proof}

\subsection{The $p$-primary symbol is determined by Stickelberger
eigenspaces}

The member $\ell=p$ of the family in Lemma~\ref{lem:psi-rules} is the one
carrying arithmetic: $p^{C}\,\|\,q-1$ with
$C=v_p(q-1)=\max(c_p,j)$, and $\chi_p$ is the $p$-power residue symbol whose
Gauss and Jacobi sums are governed by Stickelberger's theorem.  We determine
it completely.  Write $\Psi=\Psi_p$, $G=(\Z/n\Z)^{\times}=G'\times G_p$ with
$G'$ cyclic of order $p-1$ and $G_p$ the principal units of order
$p^{\,j-1}$, let $\tau:G'\to\mu_{p-1}\subseteq(\Z/p^{C}\Z)^{\times}$ be the
Teichmüller isomorphism, and let $f_1=\ord[p]{2}$.

\begin{theorem}\label{thm:pstructure}
	Let $A=\left\{k\bmod p-1:\ k\ \text{even},\ k\equiv1\pmod{f_1}\right\}$.
	Then:
	\begin{enumerate}
		\item $\Psi=\sum_{k\in A}\Psi_k$, where $\Psi_k$ is the
		      $\tau^{k}$-eigencomponent, i.e.\
		      $\Psi_k(tx)=\tau(t)^{k}\Psi_k(x)$ for $t\in G'$.
		      In particular $\abs{A}=(p-1)/(2f_1)$ if $f_1$ is odd, and
		      $A=\emptyset$ if $f_1$ is even.
		\item If $f_1$ is even then $\Psi\equiv0$: every $\bar\lambda_x$
		      is a $p^{C}$-th power in $\F_q^{\times}$, and the whole
		      $p$-primary obstruction is vacuous.
		\item $\Psi\bmod p$ is invariant under the subgroup
		      $\langle v\rangle\subseteq G_p$ of order
		      $p^{\max(0,\,j-c_p)}$ generated by the $p$-part $v$ of the
		      element $2\in G$; equivalently, $\Psi(x)\bmod p$ depends
		      only on $x$ modulo $p^{\min(j,\,c_p)}$.
		\item If $1\leq\beta\leq j-c_p$ then $\Psi(p^{\beta}x)\equiv0
		      \pmod p$.
	\end{enumerate}
\end{theorem}

\begin{proof}
	Let $M=\mathrm{Fun}(G,\Z/p^{C}\Z)$ with the translation action
	$(t\Psi)(x)=\Psi(tx)$.  Since $\abs{G'}=p-1$ is invertible in
	$\Z/p^{C}\Z$ and $\mu_{p-1}\subseteq(\Z/p^{C}\Z)^{\times}$ by Hensel's
	lemma, the idempotents
	$e_k=(p-1)^{-1}\sum_{t\in G'}\tau(t)^{-k}t$ decompose
	$M=\bigoplus_{k\bmod p-1}M_k$, and each $M_k$ is a module over the
	local ring $\Lambda=\Z/p^{C}\Z[G_p]$ with maximal ideal
	$\mathfrak m=(p,\,g-1:g\in G_p)$ and residue field $\F_p$.

	By Lemma~\ref{lem:psi-rules}, $\Psi$ satisfies $\sigma_{-1}\Psi=\Psi$
	and $\sigma_2\Psi=2\Psi$, where $\sigma_t$ is translation by the group
	element $t$ and the right-hand $2$ is the scalar $2\in\Z/p^{C}\Z$.
	As $\tau(-1)=-1$, the element $\sigma_{-1}$ acts on $M_k$ by
	$(-1)^{k}$, so $\Psi_k=0$ for odd $k$.
	Write $2=\omega v$ with $\omega\in G'$, $v\in G_p$; then
	$\omega\equiv2\pmod p$, so $\sigma_\omega$ acts on $M_k$ by
	$\tau(2)^{k}$ and the condition reads
	\[
		\left(\sigma_v-u_k\right)\Psi_k=0\ \text{in}\ M_k,
		\qquad u_k=2\,\tau(2)^{-k}\in(\Z/p^{C}\Z)^{\times}.
	\]
	Since $\sigma_v\equiv1\pmod{\mathfrak m}$, the element
	$\sigma_v-u_k$ is congruent
	to the scalar $1-u_k$ modulo $\mathfrak m$, hence is a unit of
	$\Lambda$ unless $u_k\equiv1\pmod p$.
	Now $u_k\equiv2^{1-k}\pmod p$, so $u_k\equiv1$ if and only if
	$k\equiv1\pmod{f_1}$.  This proves (1); the count of $A$ is the number
	of even residues in the arithmetic progression $1+f_1\Z$ modulo
	$p-1$, which is $(p-1)/(2f_1)$ when $f_1$ is odd and $0$ when $f_1$ is
	even.  Statement (2) is the case $A=\emptyset$.

	For (3), let $k\in A$.  Then $\tau(2)^{k}=\tau(2)$ because $\tau(2)$
	has order $f_1$, so $u_k=2/\tau(2)=\langle2\rangle$ is the principal
	unit part of $2$, and
	\[
		v_p\!\left(\langle2\rangle-1\right)
		=v_p\!\left(\langle2\rangle^{\,p-1}-1\right)
		=v_p\!\left(2^{\,p-1}-1\right)=c_p ,
	\]
	using that $\langle2\rangle\in1+p\Z_p$ and that $p-1$ is prime to $p$.
	Hence $u_k\equiv1\pmod p$ and the relation
	$(\sigma_v-u_k)\Psi_k=0$ gives
	$(\sigma_v-1)\Psi_k\equiv0\pmod p$, i.e.\ $\Psi_k\bmod p$ is
	$v$-invariant.  Since $\ord_G(2)=f=f_1p^{\max(0,j-c_p)}$ and $\omega$
	has order $f_1$ prime to $p$, the element $v$ has order
	$p^{\max(0,j-c_p)}$, and $\langle v\rangle$ is the kernel of the
	reduction $(\Z/p^{j}\Z)^{\times}\to(\Z/p^{\min(j,c_p)}\Z)^{\times}$,
	being the unique subgroup of that order in the cyclic group $G_p$.

	For (4), apply Lemma~\ref{lem:psi-rules} with $H_\beta\subseteq
	\langle v\rangle$, which holds for $\beta\leq j-c_p$: modulo $p$ the
	left-hand side is $p^{\beta}\Psi(x)\equiv0$.
\end{proof}

Part (2) is already decisive for one of the two known Wieferich primes.

\begin{corollary}\label{cor:1093-ppart}
	For $p=1093$ one has $\ord[p]{2}=364$, which is even, so
	$\Psi\equiv0$ at every level $j$: every $\bar\lambda_x$ is a
	$p^{C}$-th power in $\F_{2^{364}}^{\times}$.
	No Stickelberger, Jacobi-sum or $p$-power-residue argument can
	contribute anything at all to the case $p=1093$.
\end{corollary}

When $f_1$ is odd the symbol is nonzero in every case we have computed, and
in the extreme case $\abs{A}=1$ it has a closed form, which yields a
theorem conditional only on the single nonvanishing $\alpha\ne0$.

\begin{theorem}\label{thm:qr}
	Let $p\equiv7\pmod 8$ with $\ord[p]{2}=(p-1)/2$, so that $f_1$ is odd
	and $A=\{k_0\}$ with $k_0=f_1+1$.
	Let $j\geq1$ satisfy $\min(j,c_p)=1$, i.e.\ $j=1$, or $p$ is not a
	Wieferich prime.
	Assume $\alpha:=\Psi(1)\not\equiv0\pmod p$, i.e.\ $\bar\lambda_1$ is
	not a $p$-th power in $\F_q^{\times}$.
	Then
	\[
		\Psi(x)\equiv\alpha\,\bar x\left(\frac{\bar x}{p}\right)\pmod p
		\quad(p\nmid x),
		\qquad
		\Psi(x)\equiv0\pmod p\quad(p\mid x),
	\]
	where $\bar x=x\bmod p$ and $\left(\frac\cdot p\right)$ is the Legendre
	symbol.  Moreover every pair $x,y$ with
	$\bar\lambda_x\bar\lambda_y=1$ has $p\mid x$ and $p\mid y$; that is,
	Problem~\ref{prob:main} has an affirmative answer for $(p,j)$.
\end{theorem}

\begin{proof}
	The residues modulo $p-1=2f_1$ that are $\equiv1\pmod{f_1}$ are $1$ and
	$1+f_1$; as $f_1$ is odd, exactly the second is even, so $A=\{k_0\}$,
	$k_0=f_1+1=(p+1)/2$.
	By Theorem~\ref{thm:pstructure}(1) and (3), and because
	$\min(j,c_p)=1$, the function $\Psi\bmod p$ factors through
	$(\Z/p\Z)^{\times}$: writing $x=tu$ with $t\in G'$ and $u\in G_p$ one
	has $\Psi(x)=\tau(t)^{k_0}\Psi(u)$ and $\Psi(u)\equiv\Psi(1)$ modulo
	$p$, because $u\equiv1\pmod p$; since $\tau(t)\equiv t\equiv x\pmod p$
	this gives $\Psi(x)\equiv\alpha\bar x^{\,k_0}$.
	Euler's criterion gives
	$\bar x^{\,k_0}=\bar x^{\,(p+1)/2}
	=\bar x\cdot\bar x^{(p-1)/2}=\bar x\left(\frac{\bar x}p\right)$.
	Theorem~\ref{thm:pstructure}(4) with $c_p=1$ gives
	$\Psi(p^{\beta}x)\equiv0$ for $\beta\geq1$.

	Put $F(x)=\bar x\left(\frac{\bar x}p\right)$.  If $\bar x$ is a
	quadratic residue then $F(x)=\bar x$ is a quadratic residue; if
	$\bar x$ is a non-residue then $F(x)=-\bar x$, again a quadratic
	residue because $p\equiv3\pmod4$ makes $-1$ a non-residue.  So $F$
	takes only nonzero square values.

	Suppose $\bar\lambda_x\bar\lambda_y=1$ with $x,y\ne0$ and, after
	interchanging $x$ and $y$, with $p\nmid x$.  By
	Lemma~\ref{lem:psi-rules}, $\Psi(x)+\Psi(y)\equiv0\pmod p$.
	If $p\mid y$ this reads $\alpha F(x)\equiv0$, impossible since
	$\alpha\ne0$ and $F(x)\ne0$.
	If $p\nmid y$ it reads $F(x)=-F(y)$, with $F(x)$ a square and $-F(y)$ a
	non-zero non-square: impossible.
	Thus every solution has $p\mid x$ and $p\mid y$, which by
	Proposition~\ref{prop:product-form} is the assertion of
	Problem~\ref{prob:main}.
\end{proof}

\begin{corollary}\label{cor:qr-nullity}
	Under the hypotheses of Theorem~\ref{thm:qr} for every level
	$1\leq j'\leq j$, one has $S_{p^{j}}=\emptyset$ and $d(p^{j}-1)=0$.
\end{corollary}

\begin{proof}
	Induct on $j$.  For $j=1$ every nonzero index is prime to $p$, so
	Theorem~\ref{thm:qr} gives $S_p=\emptyset$.  For $j>1$,
	Theorem~\ref{thm:qr} leaves only pairs with $p\mid x$ and $p\mid y$;
	such a pair is a solution for $\theta^{p}$, of exact order $p^{j-1}$,
	hence is excluded by the inductive hypothesis.  Now apply
	Lemma~\ref{lem:d-count}.
\end{proof}

\begin{corollary}\label{cor:3511-level1}
	For $p=3511$ one has $p\equiv7\pmod8$ and $\ord[p]{2}=1755=(p-1)/2$,
	and $\bar\lambda_1$ is not a $p$-th power in $\F_{2^{1755}}^{\times}$
	(verified in \path{code/ai_prototypes/sine_unit_obstruction.py}).
	Hence $S_{3511}=\emptyset$ and $d(3510)=0$, by a power-residue argument
	using one power-residue computation but no enumeration of exponent pairs.
	The same conclusion is obtained unconditionally, without a
	power-residue or finite-field computation, in
	Corollary~\ref{cor:trace-family} below.
\end{corollary}

\begin{example}\label{ex:qr-family}
	The hypotheses of Theorem~\ref{thm:qr} hold for
	$p=7,23,47,71,79,103,167,191,199,239,\ldots$, and in every one of these
	cases the computation confirms both $\alpha\ne0$ and the closed form
	$\Psi(x)\equiv\alpha\bar x\left(\frac{\bar x}{p}\right)$ at level
	$j=1$, hence $d(p-1)=0$; the same was confirmed at $j=2,3$ for $p=7$
	and at $j=2$ for $p=23$, so Corollary~\ref{cor:qr-nullity} gives
	$d(7^{j}-1)=0$ for $j\leq3$ and $d(23^{j}-1)=0$ for $j\leq2$.
	For $p=7$ this reproves, in three lines, the case
	$p^{j}=49$ that Table~\ref{tab:res49} settles by fourteen resultant
	factorizations.
	The obstruction is not always this strong: for $p=31$ and $p=127$,
	where $\abs{A}=3$ and $\abs{A}=9$, the admissible ratio classes cut
	out by \eqref{eq:psi-collision} with $\ell=p$ are exactly the ratio
	classes of the solutions that do exist, so the $p$-part is sharp but
	not decisive; and for $p=17$, where $f_1=8$ is even, it is empty.
\end{example}

\subsection{The trace obstruction: an unconditional replacement}
\label{sec:traceobs}

The hypothesis $\alpha\ne0$ of Theorem~\ref{thm:qr} can be removed
entirely, at the cost of replacing the multiplicative power-residue symbol
by an additive one.  Let $\bar f=\rho_2(p^{j})$, let $E=\F_{2^{\bar f}}$ be
the residue field of the prime $\mathfrak q\mid2$ of $K_j^{+}$ --- so that
$\bar\lambda_x\in E$ for every $x$ --- and put
\[
	\delta=
	\begin{cases}
		1,&-1\in\langle2\rangle\subseteq(\Z/p^{j}\Z)^{\times}
		\quad(\text{equivalently }f=2\bar f),\\
		0,&\text{otherwise},
	\end{cases}
	\qquad
	\tau(x)=\Tr_{E/\F_2}\!\left(\bar\lambda_x\right).
\]

\begin{theorem}\label{thm:trace}
	For every $x\not\equiv0\pmod{p^j}$ one has
	$\Tr_{E/\F_2}\!\left(\bar\lambda_x^{-1}\right)=\delta$.
	Consequently, if $\bar\lambda_x\bar\lambda_y=1$ then
	\[
		\tau(x)=\tau(y)=\delta .
	\]
\end{theorem}

\begin{proof}
	Put $u=\theta^{x}$ and $c=\bar\lambda_x$, so $u^{2}+cu+1=0$;
	substituting $u=cS$ gives the Artin--Schreier equation
	$S^{2}+S=c^{-2}$, whose roots are $u/c$ and $u^{-1}/c$.  It has a root
	in $E$ if and only if $\Tr_{E/\F_2}(c^{-2})=\Tr_{E/\F_2}(c^{-1})=0$,
	and also if and only if $u\in E$.  Now $u\in E$ exactly when
	$2^{\bar f}\equiv1$, i.e.\ when $\delta=0$: for $\delta=1$ one has
	$u^{2^{\bar f}}=u^{-1}\ne u$ because $u^{2}\ne1$.
	For the consequence, $\bar\lambda_y=\bar\lambda_x^{-1}$, so
	$\tau(x)=\Tr_{E/\F_2}(\bar\lambda_y^{-1})=\delta$ and symmetrically.
\end{proof}

\begin{theorem}\label{thm:twoorbit}
	Suppose $\ord[p^{j}]{2}=\varphi(p^{j})/2$ is odd, i.e.\ that
	$\langle2\rangle$ has index two in $(\Z/p^{j}\Z)^{\times}$ and
	$-1\notin\langle2\rangle$.  Then $\delta=0$, $E=\F_{2^{f}}$, and
	\[
		\tau(x)=
		\begin{cases}
			1,& v_p(x)=j-1,\\
			0,& v_p(x)\leq j-2 .
		\end{cases}
	\]
	Hence every solution of \eqref{eq:main} has
	$v_p(x),v_p(y)\leq j-2$ in product form; for $j=1$ there is none, so
	$S_p=\emptyset$ and $d(p-1)=0$.
\end{theorem}

\begin{proof}
	Oddness of $\varphi(p^j)/2=p^{\,j-1}(p-1)/2$ forces $p\equiv3\pmod4$,
	so $-1$ is a non-square modulo every power of $p$; and
	$\langle2\rangle$, being of index two in a cyclic group, is the group
	of squares.  Hence $\delta=0$, $\bar f=f$ and $\theta\in E$.  Put
	$t(x)=\Tr_{E/\F_2}(\theta^{x})$, so $\tau(x)=t(x)+t(-x)$ and
	$t(2x)=t(x)$.

	Fix $\beta\leq j-1$.  Reduction
	$(\Z/p^{j}\Z)^{\times}\to(\Z/p^{\,j-\beta}\Z)^{\times}$ carries squares
	onto squares, so the image of $\langle2\rangle$ again has index two and
	misses $-1$; therefore the indices of level $\beta$ split into exactly
	two $\langle2\rangle$-orbits, interchanged by $x\mapsto-x$, and $t$ is
	constant on each, with values $t_\beta^{\pm}$.  Thus $\tau$ is constant
	on level $\beta$, equal to $\tau_\beta=t_\beta^{+}+t_\beta^{-}$.

	Summing $t$ over level $\beta$ and using that in characteristic two the
	sum of all $p^{k}$-th roots of unity vanishes for $k\geq1$, so that the
	sum of the primitive $p^{k}$-th roots of unity is $0$ for $k\geq2$ and
	$1$ for $k=1$,
	\[
		\frac{\varphi(p^{\,j-\beta})}{2}\,\tau_\beta
		=\sum_{v_p(x)=\beta}t(x)
		=\begin{cases}
			\Tr_{E/\F_2}(0)=0,&\beta\leq j-2,\\
			\Tr_{E/\F_2}(1)=f\bmod 2=1,&\beta=j-1 .
		\end{cases}
	\]
	Since $\varphi(p^{\,j-\beta})/2$ is odd, reduction modulo $2$ gives
	$\tau_\beta=0$ for $\beta\leq j-2$ and $\tau_{j-1}=1$.
	Now apply Theorem~\ref{thm:trace}: a collision needs
	$\tau(x)=\tau(y)=\delta=0$.
	Finally $d(p-1)=\abs{S_p}/2$ by Lemma~\ref{lem:d-count}.
\end{proof}

\begin{corollary}\label{cor:trace-family}
	Let $p\equiv7\pmod 8$ with $\ord[p]{2}=(p-1)/2$.  Then
	$S_p=\emptyset$ and $d(p-1)=0$ \emph{unconditionally}: no
	power-residue nonvanishing hypothesis is needed.
	In particular $d(3510)=0$ is a theorem, and the same holds for
	$p=7,23,47,71,79,103,167,191,199,\ldots$
\end{corollary}

\begin{proof}
	The hypothesis gives $\ord[p]{2}=\varphi(p)/2$, odd because
	$p\equiv3\pmod4$.  Apply Theorem~\ref{thm:twoorbit} with $j=1$.
\end{proof}

\subsection{A global trace law}

The Artin--Schreier obstruction uses the trace of
$\bar\lambda_x$.  A different trace identity appears after rewriting the
collision as a weight-five ``box plus centre'' relation.

\begin{theorem}[Global trace law]\label{thm:global-trace}
	Let
	\[
		T(x)=\Tr_{\F_{2^f}/\F_2}(\theta^x).
	\]
	If $\bar\lambda_x\bar\lambda_y=1$, then
	\begin{equation}\label{eq:global-trace}
		T(x)+T(y)=f\pmod2.
	\end{equation}
\end{theorem}

\begin{proof}
	Multiplying
	$(\theta^x+\theta^{-x})(\theta^y+\theta^{-y})=1$
	by $\theta^{x+y}$ gives
	\[
		1+\theta^{2x}+\theta^{2y}
		+\theta^{2x+2y}+\theta^{x+y}=0.
	\]
	Apply the absolute trace.  Frobenius invariance gives
	$\Tr(\theta^{2x})=T(x)$,
	$\Tr(\theta^{2y})=T(y)$, and
	$\Tr(\theta^{2x+2y})=T(x+y)$; the last term contributes another
	$T(x+y)$ and cancels it.  Finally
	$\Tr(1)=f\bmod2$.
\end{proof}

\begin{corollary}[Three exact new cases]\label{cor:global-trace-cases}
	Exact trace enumeration gives
	\[
		d(72)=d(88)=d(232)=0,
	\]
	corresponding to $p=73,89,233$.
\end{corollary}

\begin{proof}
	For each prime, enumerate the indices satisfying the necessary
	Artin--Schreier condition of Theorem~\ref{thm:trace}, split them by
	$T(x)$, and apply \eqref{eq:global-trace}.  In all three cases no
	ordered pair survives both conditions.  The computation is exact
	bit-vector arithmetic in
	\path{code/ai_prototypes/plateau_endpoint.py}; direct evaluation of
	$S_p$ independently confirms the empty set.
\end{proof}

Further applications of the box-plus-centre identity produce additional
linear trace relations, for example
\[
	T(x-2y)+T(2x-y)+T(3x-3y)=0.
\]
They strengthen the filter in individual cases but have not yielded a
uniform endpoint theorem.  The formal constructions following
Corollary~\ref{cor:multiplicative-nogo} can be chosen to satisfy the
additive distribution law, Theorem~\ref{thm:trace},
Theorem~\ref{thm:global-trace}, and the known level sums while still
containing a primitive collision at both known Wieferich plateaus.  Thus
these trace laws are genuine progress, but any general proof must use the
full nonlinear multiplication law rather than only the currently extracted
linear consequences.

\subsection{Exact Kloosterman formulas at the full-torus extremes}
\label{sec:kloosterman}

For $r\geq1$ define the binary Kloosterman sum
\[
	K_r=\sum_{z\in\F_{2^r}}
	(-1)^{\Tr_{\F_{2^r}/\F_2}(z+z^{-1})},
	\qquad 0^{-1}=0.
\]

\begin{theorem}[Mersenne and Fermat formulas]\label{thm:kloosterman}
	\begin{enumerate}
		\item If $p=2^r-1$ is a Mersenne prime, $r\geq3$, then
		      \[
			      d(p-1)=\frac{p-3+K_r}{2}.
		      \]
		\item If $p=2^m+1$ is a Fermat prime, $m\geq2$, then
		      \[
			      d(p-1)=\frac{p-1+K_m}{2}.
		      \]
	\end{enumerate}
\end{theorem}

\begin{proof}
	In the Mersenne case $\mu_p=\F_{2^r}^{\times}$; in the Fermat case
	$\mu_p$ is the full norm-one torus in
	$\F_{2^{2m}}/\F_{2^m}$.  Therefore the trace-value set
	$\{\bar\lambda_x:x\neq0\}$ is the entire Artin--Schreier set described
	by Theorem~\ref{thm:trace}, so its necessary condition is also
	sufficient.

	Expand the two trace indicators for a pair
	$\alpha,\alpha+1$ and sum over $\alpha\neq0,1$.  The linear terms are
	evaluated by orthogonality of the additive character.  For the
	quadratic term use
	\[
		\alpha^{-1}+(\alpha+1)^{-1}
		=(\alpha^2+\alpha)^{-1}
	\]
	and the two-to-one map
	$\alpha\mapsto\alpha^2+\alpha$ onto the trace-zero subspace.
	The remaining character sum is exactly $K_r$ or $K_m$.
	Finally use Lemma~\ref{lem:d-count}.  This gives the two formulas.
\end{proof}

By the Weil bound, and more sharply by the congruences and recurrences for
binary Kloosterman sums \cite{LachaudWolfmann1990,Carlitz1969},
$d(p-1)>0$ for every Mersenne prime $p\geq31$ and every Fermat prime
$p\geq5$.  This explains the level-one solutions for
$p=17,31,127,257$.  When $\mu_{p^j}$ is a proper subgroup of the torus,
the same expansion becomes an average of twisted Kloosterman sums.  The
Weil error then exceeds the expected main term on the plateau, so the
exact full-torus evaluation does not extend to the endpoint problem.

\subsection{Why the method stops exactly at the plateau}

\begin{proposition}\label{prop:plateau-degenerate}
	Let $2\leq j\leq c_p$, so that $(p,j)$ lies on the plateau.  Then in
	the notation of Theorem~\ref{thm:pstructure}, $v=1$ and $C=c_p$, so
	$u_k=\langle2\rangle\equiv1\pmod{p^{C}}$ and the relation
	$(\sigma_v-u_k)\Psi_k=0$ is the trivial relation $0=0$.
	Beyond the eigenspace support of Theorem~\ref{thm:pstructure}(1) and
	the norm relation of Lemma~\ref{lem:lam-relations}, no constraint on
	$\Psi$ survives; the ambient eigenspace has $p^{\,j-1}-1$ free
	parameters in $\Z/p^{C}\Z$ after fixing its augmentation.
\end{proposition}

\begin{proof}
	On the plateau $\ord[p^j]{2}=\ord[p]{2}=f_1$ is prime to $p$, so the
	$p$-part $v$ of $2\in G$ is trivial, and
	$C=v_p(2^{f}-1)=c_p\geq j$.
	By the computation in the proof of
	Theorem~\ref{thm:pstructure}(3),
	$v_p(u_k-1)=c_p=C$, i.e.\ $u_k=1$ in $\Z/p^{C}\Z$; so the relation
	imposed by $\sigma_2$ is vacuous.  The eigencomponent $M_k$ is free of
	rank $p^{\,j-1}$ over $\Z/p^{C}\Z$ and the relation of
	Lemma~\ref{lem:lam-relations} fixes only its augmentation.
\end{proof}

\begin{heuristic}
	If the remaining free values of $\Psi$ behave randomly, then the
	necessary condition \eqref{eq:psi-collision} for $\ell=p$ should remove
	a proportion about $1-p^{-c_p}$ of pairs, leaving on the order of
	$\varphi(p^j)^2p^{-c_p}$ candidates. At $j=c_p$ this is on the order of
	$p^j$. This explains why the $p$-primary condition is not expected to
	decide Problem~\ref{prob:main}, but it is not a count of the actual
	values of $\Psi$.
\end{heuristic}

\begin{example}[the plateau case $p=3511$, $j=2$]\label{ex:3511-plateau}
	Here $f_1=1755$ is odd, $c_p=2$, $A=\{1756\}$, and
	Theorem~\ref{thm:qr} proves the case $j=1$ outright.  At $j=2$ the
	descent of Theorem~\ref{thm:pstructure}(3) says only that
	$\Psi\bmod p$ depends on $x$ modulo $p^{2}$, which is no information,
	and the $p$ values of $\Psi$ on $G_p$ are unconstrained.
	The computation \path{code/ai_prototypes/plateau_ppart_3511.py}
	evaluates $\Psi$ on $140$ elements of $G_p$ inside
	$\F_{2^{1755}}$, confirms the eigenrelation
	$\chi(\bar\lambda_{tx})=\chi(\bar\lambda_x)^{\tau(t)^{k_0}}$, and finds
	$4$ sampled pairs $(u_1,u_2)$ satisfying the $p$-primary necessary
	condition \eqref{eq:psi-collision} --- against $2.79$ expected at random.
	This is heuristic evidence that the obstruction, decisive at level $1$,
	does not decide level $2$; these are not five-term solutions.
\end{example}

\begin{remark}\label{rem:other-ell}
	The same eigenspace analysis applies to $\ell\ne p$: $\Psi_\ell$ is
	supported on those characters $\psi$ of $G$ with values in
	$\bar\F_\ell$ that satisfy $\psi(2)=2$ and $\psi(-1)=1$.  Such
	characters exist for many $\ell$ --- the constraint $\psi(2)=2$ alone
	is always solvable, because $\ord[\ell]{2}$ divides $\ord_G(2)=f$ ---
	so no analogue of Theorem~\ref{thm:pstructure}(2) is available in
	general.  Moreover, if $2^{i}\equiv-1\pmod{\ell^{C_\ell}}$ for some
	$i$, then every pair $(x,2^{i}x)$ satisfies \eqref{eq:psi-collision}
	for that $\ell$, so the $\ell$-primary condition alone can never be
	empty.  Those pairs are precisely the ones removed by
	Theorem~\ref{thm:frobratio}; the two
	obstructions of this section are therefore complementary, and neither
	is subsumed by the other.
	Individual $\ell\ne p$ can nevertheless be decisive: for $p=7$, $j=2$
	each of $\ell=7,127,337$ already excludes every admissible pair, and
	for $p=23$ both $\ell=23$ and $\ell=89$ do.
	This is not a usable route on the plateau, since it presupposes the
	factorization of $2^{364}-1$ or $2^{1755}-1$.
\end{remark}

\begin{remark}\label{rem:char-specific}
	Theorem~\ref{thm:qr} is genuinely characteristic-two input, and is
	therefore not in conflict with
	Corollary~\ref{cor:no-structural-proof}.
	Replacing $2$ by another prime $\ell$, the Frobenius relation becomes
	$\bar\lambda_{\ell x}=\bar\lambda_x^{\,\ell}$ and the eigenspace
	condition of Theorem~\ref{thm:pstructure} becomes
	$k\equiv1\pmod{\ord[p]{\ell}}$, so the admissible set $A$ has
	$(p-1)/(2\ord[p]{\ell})$ elements and the argument degrades as
	$\ord[p]{\ell}$ shrinks.
	For the counterexample $\ell=491$, $p=7$, $j=2$ of
	Table~\ref{tab:counterexamples} one has $\ord[7]{491}=1$, so $A$
	consists of all even $k$, the symbol is unconstrained, and no
	contradiction arises --- exactly as it must be.
\end{remark}

\section{Modular vanishing sums: what the theorems give}\label{sec:modular}

\subsection{Congruence modulo the rational prime $2$}

\begin{theorem}[Dvornicich--Zannier \cite{DvornicichZannier2002}]\label{thm:DZ}
	Let $\ell$ be a rational prime, let $\zeta_0=1,\zeta_1,\ldots,\zeta_{k-1}$
	be pairwise distinct roots of unity, and let $a_i\in\Q^{\times}$
	satisfy $\sum_{i}a_i\zeta_i\equiv0\pmod{\ell}$ while no proper subsum
	is $\equiv0\pmod{\ell}$.
	Then $Q=\operatorname{lcm}_i\ord{\zeta_i}$ is squarefree and
	\[
		\sum_{q\mid Q}(q-2)\leq k-2,
	\]
	the sum being over primes $q$ dividing $Q$.
\end{theorem}

We flag explicitly that the congruence in Theorem~\ref{thm:DZ} is modulo the
\emph{rational} prime $\ell$, that is, in $\Z[\zeta]/\ell\Z[\zeta]$, and not
modulo a prime ideal above $\ell$.
This is not a technicality.

\begin{proposition}\label{prop:DZ-fails-ideals}
	The literal analogue of Theorem~\ref{thm:DZ} for congruences modulo a
	prime ideal is false.
	In $\F_8=\F_2[x]/(x^3+x+1)$ the element $\alpha=x$ has order $7$ and
	$\alpha^{3}=\alpha+1$, so $1+\alpha+\alpha^{3}=0$.
	Lifting, there is a prime $\mathfrak{p}\mid2$ of $\Z[\zeta_7]$ with
	\[
		1+\zeta_7+\zeta_7^{3}\equiv0\pmod{\mathfrak{p}},
	\]
	and no proper subsum is $\equiv0\pmod{\mathfrak{p}}$, because
	$1+\zeta_7^{e}$ has norm $\Phi_7(-1)=1$ and is therefore a unit for
	$e\not\equiv0$.
	Here $k=3$ and $Q=7$, so the conclusion of Theorem~\ref{thm:DZ} would
	read $5\leq1$.
\end{proposition}

Nevertheless the congruence modulo the rational prime does apply to our
situation, and yields the following.
We give an elementary self-contained proof, which is stronger in that it
requires no hypothesis on subsums.

\begin{theorem}\label{thm:all-factors}
	Let $p>5$ be prime and $j\geq1$.
	Then for all $a,b$,
	\[
		\Lambda_{a,b}\not\equiv0\pmod{2\Z[\zeta_{p^j}]} .
	\]
	Equivalently, $\Phi_{p^j}(x)$ does not divide
	$1+x^{a}+x^{-a}+x^{b}+x^{-b}$ in $\F_2[x]/(x^{p^j}+1)$; that is,
	\eqref{eq:main} can never hold simultaneously for all primitive
	$p^j$-th roots of unity in $\overline{\F}_2$.
\end{theorem}

\begin{proof}
	Work in $R=\F_2[x]/(x^{p^j}+1)$.
	Since $x^{p^j}+1=\Phi_{p^j}(x)\cdot(x^{p^{j-1}}+1)$, every element of
	the ideal $\Phi_{p^j}R$ has a unique representative
	$\Phi_{p^j}(x)A(x)$ with $\deg A<p^{j-1}$.
	Now $\Phi_{p^j}(x)=\Phi_p(x^{p^{j-1}})=\sum_{k=0}^{p-1}x^{kp^{j-1}}$, so
	\[
		\Phi_{p^j}(x)A(x)=\sum_{k=0}^{p-1}x^{kp^{j-1}}A(x),
	\]
	and the $p$ summands have pairwise disjoint supports, contained in the
	intervals $[kp^{j-1},(k+1)p^{j-1})$.
	Hence
	\[
		\wt\!\left(\Phi_{p^j}A\right)=p\cdot\wt(A),
	\]
	so every nonzero element of $\Phi_{p^j}R$ has weight a positive
	multiple of $p$, in particular weight at least $p$.

	The image of $1+x^{a}+x^{-a}+x^{b}+x^{-b}$ in $R$ has weight at most
	$5$, and its coefficient at the exponent $0$ equals
	$1+2[a\equiv0]+2[b\equiv0]=1$ in $\F_2$, so it is nonzero.
	Since $5<p$, it is not in $\Phi_{p^j}R$.
\end{proof}

\begin{remark}
	Theorem~\ref{thm:DZ} yields the same conclusion, at the cost of
	verifying its hypotheses.
	Assume $\Lambda_{a,b}\equiv0\pmod{2}$.
	By Lemma~\ref{lem:degenerate} (whose proof applies verbatim modulo $2$,
	using that $1+\zeta^{e}$ is a unit for $e\not\equiv0$) the five roots
	are distinct.
	No subsum of size~$1$ or~$2$ is $\equiv0\pmod2$, again because
	$1+\zeta^{e}$ is a unit.
	For a subsum of size $k\in\{3,4\}$ with no vanishing proper subsum,
	Theorem~\ref{thm:DZ} forces its conductor $Q\mid p^j$ to be squarefree
	with $\sum_{q\mid Q}(q-2)\leq k-2\leq2$, so $Q=1$ and all its roots
	equal $1$, contradicting distinctness.
	Applying Theorem~\ref{thm:DZ} to the full five-term sum gives $Q$
	squarefree, hence $Q=p$, and $p-2\leq3$, i.e.\ $p\leq5$.
	The elementary weight argument is preferable because it is
	unconditional on the exact wording of the published theorem, which we
	were unable to consult directly, and because it gives the sharper
	statement that all weights in $\Phi_{p^j}R$ are multiples of $p$.
\end{remark}

\subsection{Why no such theorem can hold for a single prime ideal}

\begin{proposition}\label{prop:ubiquity}
	Let $\theta\in\F_{2^{f}}$ have exact order $N$, $f=\ord[N]{2}$, and let
	$w_0$ be minimal with $\sum_{i=0}^{w_0}\binom{N}{i}>2^{f}$.
	Then there are $w\leq2w_0$ pairwise distinct residues
	$k_1,\ldots,k_w$ modulo $N$ with
	$\theta^{k_1}+\cdots+\theta^{k_w}=0$.
\end{proposition}

\begin{proof}
	The $\F_2$-linear map $\F_2^{\Z/N\Z}\to\F_{2^{f}}$ sending
	$\delta\mapsto\sum_k\delta_k\theta^{k}$ is surjective.
	By the pigeonhole principle two distinct vectors of weight at most
	$w_0$ have the same image; their difference is a nonzero kernel vector
	of weight at most $2w_0$.
\end{proof}

Proposition~\ref{prop:ubiquity} shows that modular vanishing sums of small
weight are unavoidable as soon as $N$ is large compared with $2^{f}$; no
Mann-type or Conway--Jones-type inequality can survive reduction modulo a
single prime ideal.
In coding-theoretic language \cite[Ch.~7]{MacWilliamsSloane}, the kernel in
Proposition~\ref{prop:ubiquity} is the dual of the irreducible cyclic code
of length $N$ and dimension $f$ generated by the sequence
$\left(\Tr_{\F_{2^{f}}/\F_2}(\lambda\theta^{k})\right)_{k}$, and the
statement is the Gilbert--Varshamov-type upper bound on its minimum
distance; Proposition~\ref{prop:DZ-fails-ideals} is the observation that
this distance can be as small as $3$.
Lam and Leung reached the same conclusion from the opposite direction: they
determined the weight sets $W_\ell(m)$ of vanishing sums of $m$-th roots of
unity in $\overline{\F}_\ell$ and observed that the characteristic-zero
answer fails, remarking that they are ``left with no viable conjecture on
the structure of the weight set $W_p(m)$ in characteristic $p$''
\cite[\S1]{LamLeung1996}.

\subsection{The exact size of the gap}

\begin{theorem}\label{thm:gap}
	Let $p>5$ and $2\leq j\leq c_p$, so that $j$ lies on the plateau and
	$f=\ord[p^j]{2}=\ord[p]{2}$ divides $p-1$.
	Suppose $(a,b)$ is a counterexample to Problem~\ref{prob:main}, with
	$\gcd(a,p)=1$ after possibly interchanging $a$ and $b$.
	Then
	\[
		1\leq\abs{\Sigma_{a,b}}\leq g-1,
		\qquad
		g=\frac{\varphi(p^j)}{f}\geq p^{\,j-1}\geq p>5 ,
	\]
	and $\abs{\Sigma_{a,b}}\geq2$ whenever $-1\notin\langle2\rangle$ in
	$(\Z/p^j\Z)^{\times}$.
	Proposition~\ref{prop:sigma} cannot give a lower bound exceeding
	$\abs{H}\leq4$.
\end{theorem}

\begin{proof}
	The lower bounds are Proposition~\ref{prop:sigma} together with
	$\Sigma_{a,b}\ne\emptyset$, which is
	Proposition~\ref{prop:dictionary}(1).
	For the upper bound, $\abs{\Sigma_{a,b}}=g$ would mean that every prime
	above $2$ divides $\Lambda_{a,b}$, i.e.\ that
	$\Lambda_{a,b}\equiv0\pmod{2\Z[\zeta_{p^j}]}$, contradicting
	Theorem~\ref{thm:all-factors}.
	Finally $g=\varphi(p^j)/f=p^{j-1}(p-1)/f\geq p^{j-1}$ because
	$f\mid p-1$ on the plateau.
\end{proof}

This is the precise answer to the question in the title of this section.
\emph{The reciprocal, product-one symmetry of the five-term relation supplies
at most a factor of $4$, and generically a factor of $2$, toward the
$g\geq p^{j-1}$ primes above $2$ that a Dvornicich--Zannier style argument
needs.}
The deficit is not a constant that could be absorbed by sharpening the
inequality $\sum_{q\mid Q}(q-2)\leq k-2$; it grows like $p^{j-1}$.

There is also an archimedean constraint, which we record because it shows
how far the ``all primes above $2$'' case is from being tight.
By Proposition~\ref{prop:dictionary}(3)--(4),
a counterexample forces
$2^{f\abs{\Sigma}}\ \big|\ \Res(\Phi_{p^j},G_r)$, while
\[
	\frac{1}{\varphi(p^j)}\log\abs{\Res\!\left(\Phi_{p^j},G_r\right)}
	\longrightarrow
	m(G_r),
\]
$m$ denoting the logarithmic Mahler measure, and by Lawton's theorem
\cite{Lawton1983}
\[
	\lim_{r\to\infty}m(G_r)=m\!\left(1+x+x^{-1}+y+y^{-1}\right)=0.2513\ldots
	=0.3626\ldots\times\log2 .
\]
(The numerical value was computed by midpoint quadrature on a
$4096\times4096$ grid; the polynomial belongs to Boyd's family
$x+x^{-1}+y+y^{-1}+k$ at $k=1$ \cite{Boyd1998}.)
Therefore, asymptotically, $f\abs{\Sigma}\lesssim0.3626\,\varphi(p^j)$, i.e.
\[
	\abs{\Sigma_{a,b}}\ \lesssim\ 0.3626\,g .
\]
Here the roots-of-unity limit and the Lawton limit are coupled through the
choice of the residue class $r\bmod p^j$, so this estimate is heuristic rather
than uniform in $r$.
Thus the archimedean size heuristic is consistent with
Theorem~\ref{thm:all-factors}, while leaving the whole range
$2\leq\abs{\Sigma}\lesssim0.36\,g$ open.

\section{The cyclic-code and Zetterberg-code attack}\label{sec:coding}

Fix an odd $n$ and an element $\theta$ of exact order $n$.  The local
vanishing code is
\[
	\mathcal K_n
	=\left\{(c_t)_{t\in\Z/n\Z}\in\F_2^n:
		\sum_t c_t\theta^t=0\right\}.
\]
It is a binary cyclic code of dimension $n-\ord[n]{2}$.  A nondegenerate
five-term relation is the codeword with support
$\{0,\pm a,\pm b\}$.  This is a \emph{single-prime} code, unlike the ideal
$\Phi_{p^j}R$ in Theorem~\ref{thm:all-factors}, which imposes vanishing at
every prime above two.

\begin{theorem}[Interleaving off the plateau]\label{thm:code-interleaving}
	Let $n=p^j$ with $j>c_p$, let $\theta$ have order $p^j$, and put
	$\eta=\theta^p$.  Every polynomial of degree less than $p^j$ has a
	unique decomposition
	\[
		A(X)=\sum_{r=0}^{p-1}X^rA_r(X^p),
		\qquad \deg A_r<p^{j-1}.
	\]
	Then
	\[
		A(\theta)=0
		\quad\Longleftrightarrow\quad
		A_r(\eta)=0\ \text{for every }r.
	\]
	Thus $\mathcal K_{p^j}$ is the interleaving of $p$ copies of
	$\mathcal K_{p^{j-1}}$.  Every codeword of weight at most five is
	supported in one residue class modulo $p$.
\end{theorem}

\begin{proof}
	For $j>c_p$ one has
	\[
		[\F_2(\theta):\F_2(\eta)]
		=\frac{\ord[p^j]{2}}{\ord[p^{j-1}]{2}}=p.
	\]
	Hence $1,\theta,\ldots,\theta^{p-1}$ are linearly independent over
	$\F_2(\eta)$, and the equivalence follows.
	A nonzero component $A_r(\eta)$ cannot have weight one or two:
	$\eta^s\neq0$, and
	$\eta^s+\eta^t=0$ forces $s\equiv t\pmod{p^{j-1}}$.
	If two interleaved components were nonzero, the total weight would
	therefore be at least six.
\end{proof}

\begin{corollary}
	The interleaving theorem gives another proof of the five-term descent
	for $j>c_p$: since the support contains the exponent $0$, all five
	exponents lie in the zero residue class modulo $p$, so $p\mid a,b$.
\end{corollary}

On the plateau the extension degree in the proof is one, so the
interleaving collapses.  In the nonsplit case, with
$f=2\rho_2(p^j)$ and $q=2^{\rho_2(p^j)}$, the group $\mu_{p^j}$ lies in
the norm-one circle of order $q+1$ in $\F_{q^2}$.  The code
$\mathcal K_{p^j}$ is then a shortening of the binary Zetterberg code on
that circle.  The standard minimum-distance result for the ambient code
is five \cite{LachaudWolfmann1990}; shortening rules out weights below
five, but the forbidden relation has weight exactly five.  Thus the
general code bound reaches the boundary and stops.

Weight enumerators or classifications of minimum Zetterberg words do not
finish the problem either.  They classify words on the full norm-one
circle, whereas the endpoint asks whether a minimum word has the special
reciprocal support
\[
	\{1,\theta^{\pm a},\theta^{\pm b}\}
\]
inside one prescribed $p$-Sylow subgroup.  Retaining that support condition
turns the classification back into the original intersection problem.
In split cases the analogous irreducible cyclic-code bounds are weaker
still.  Coding theory therefore gives a clean structural proof off the
plateau and an exact explanation of why minimum-distance technology alone
cannot cross it.

\section{Counterexamples in other residue characteristics}\label{sec:counterexamples}

Every structural ingredient used so far --- the cyclotomic field
$\Q(\zeta_{p^j})$, its Galois group, the decomposition group
$\langle\ell\rangle$ at a prime above $\ell$, the classification of exact
vanishing sums, the Dvornicich--Zannier inequality, the reciprocal
product-one symmetry, and the plateau condition $c_p\geq j$ --- makes sense
verbatim with $2$ replaced by an arbitrary prime $\ell$.
We now show that the resulting statement is false.

For a prime $\ell\ne p$ define
\[
	f_1(\ell,p)=\ord[p]{\ell},
	\qquad
	c(\ell,p)=v_p\!\left(\ell^{f_1}-1\right),
\]
so that $c(\ell,p)\geq2$ says exactly that $p$ is a Wieferich prime to base
$\ell$, and $\mu_{p^{j}}\subseteq\F_{\ell^{f_1}}$ for $j\leq c(\ell,p)$.

\begin{theorem}\label{thm:counterexamples}
	The analogue of Problem~\ref{prob:main} with $2$ replaced by $\ell$ has
	a negative answer.
	Each row of Table~\ref{tab:counterexamples} gives a prime $\ell$, a
	prime $p>5$, an exponent $j\geq2$, an element $\theta$ of exact order
	$p^{j}$ in $\F_{\ell}$, and integers $a,b$ prime to $p$ with
	\[
		1+\theta^{a}+\theta^{-a}+\theta^{b}+\theta^{-b}\equiv0\pmod{\ell},
	\]
	and with $1,\theta^{\pm a},\theta^{\pm b}$ pairwise distinct.
	In every row $c(\ell,p)\geq j$, so the example lies on the base-$\ell$
	plateau.
\end{theorem}

\begin{table}[h]
	\centering
	\begin{tabular}{rrrrrrrr}
		\toprule
		$\ell$ & $p$ & $j$ & $p^j$ & $\theta$ & $a$ & $b$ & $c(\ell,p)$\\
		\midrule
		$491$   & $7$  & $2$ & $49$  & $42$   & $9$  & $15$  & $2$\\
		$1373$  & $7$  & $3$ & $343$ & $16$   & $29$ & $152$ & $3$\\
		$29009$ & $7$  & $2$ & $49$  & $7394$ & $4$  & $19$  & $2$\\
		$727$   & $11$ & $2$ & $121$ & $358$  & $2$  & $30$  & $2$\\
		$677$   & $13$ & $2$ & $169$ & $16$   & $28$ & $40$  & $2$\\
		\bottomrule
	\end{tabular}
	\caption{Verified counterexamples over prime fields.
	Each is checkable by five modular exponentiations.}
	\label{tab:counterexamples}
\end{table}

\begin{proof}
	Each row is a finite verification, carried out in
	\path{code/ai_prototypes/plateau_scan.py} and independently rechecked
	with plain modular arithmetic.
	For instance, in $\F_{491}$ the element $\theta=42$ satisfies
	$42^{49}\equiv1$ and $42^{7}\not\equiv1$, and
	\[
		1+42^{9}+42^{40}+42^{15}+42^{34}\equiv0\pmod{491},
	\]
	with $\gcd(9,7)=\gcd(15,7)=1$.
	Here $491\equiv1\pmod{49}$, so $f_1=1$ and $c(491,7)=v_7(490)=2$.
\end{proof}

Further counterexamples occur over non-prime fields, for example
$\ell=97$, $p=7$, $j=2$ in
$\F_{97}[x]/(x^2+18x+1)$ with $(a,b)=(11,17)$;
$\ell=241$, $p=11$, $j=2$ with $(a,b)=(17,18)$;
$\ell=337$, $p=13$, $j=2$ with $(a,b)=(9,24)$;
$\ell=577$, $p=17$, $j=2$ with $(a,b)=(21,38)$;
and $\ell=41$, $p=29$, $j=2$ in a quartic extension with $(a,b)=(4,7)$.

\begin{corollary}\label{cor:no-structural-proof}
	No proof of Problem~\ref{prob:main} can use only the following inputs:
	the arithmetic of $\Q(\zeta_{p^j})$ and its decomposition groups; the
	classification of exact vanishing sums of roots of unity
	(Mann~\cite{Mann1965}, Conway--Jones~\cite{ConwayJones1976},
	Lenstra~\cite{Lenstra1979}, Lam--Leung~\cite{LamLeung2000});
	congruence analogues of the Conway--Jones inequality
	\cite{DvornicichZannier2002} and its number-field-coefficient companion
	\cite{DvornicichZannier2000}; the
	reciprocal product-one symmetry of the five terms; and the plateau
	condition.
	All of these hold verbatim in the configurations of
	Table~\ref{tab:counterexamples}.
\end{corollary}

\begin{remark}[a characteristic-three accident]
	The smallest plateau configuration of all, $\ell=3$ and $p=11$ (with
	$3^{5}-1=242=2\cdot11^{2}$, so $c(3,11)=2$ and
	$\mu_{121}\subseteq\F_{243}$), has \emph{no} solutions, even though the
	naive count $p^{2j}/\ell^{f_1}=121^2/243\approx60$ predicts many.
	The reason is special to characteristic three.
	In $\F_{243}$ the group $\mu_{121}$ has index two, hence is the group of
	squares, so every $u\in\mu_{121}$ is $w^{2}$ and
	\[
		u+u^{-1}=\left(w+w^{-1}\right)^{2}-2=s^{2}+1,
		\qquad s=w+w^{-1}.
	\]
	Writing $v=z^{2}$ and $t=z+z^{-1}$ likewise, the equation
	$1+u+u^{-1}+v+v^{-1}=0$ becomes $3+s^{2}+t^{2}=0$, that is
	$s^{2}+t^{2}=0$.
	If $t=0$ then $v+v^{-1}=1$, so $v^{2}-v+1=(v+1)^{2}=0$ in
	characteristic three and $v=-1$, which does not lie in the odd-order
	group $\mu_{121}$; so $t\ne0$ and $(s/t)^{2}=-1$.
	But $-1$ has order two in $\F_{243}^{\times}$, a group of order
	$242=2\cdot11^{2}$, whereas every square has odd order; so $-1$ is not a
	square and there is no solution.
	This is precisely the kind of characteristic-specific accident that a
	proof for $\ell=2$ would have to exhibit, and it has no analogue at
	$\ell=2$.
\end{remark}

\section{The quantitative criterion}\label{sec:quantitative}

Table~\ref{tab:counterexamples} and the scan of
\path{code/ai_prototypes/plateau_scan.py} are explained by a single
heuristic.

\begin{heuristic}\label{heur}
	Let $n=p^j$ and
	\[
		\rho_\ell(n)=\min\{r>0:\ell^r\equiv\pm1\pmod n\},
		\qquad Q=\ell^{\rho_\ell(n)}.
	\]
	The number of pairs $(u,v)\in\mu_n^2$ with
	$1+u+u^{-1}+v+v^{-1}=0$ is about $n^{2}/Q$.
	Solutions are therefore expected exactly when
	\begin{equation}\label{eq:window}
		p^{\,j}\ \gtrsim\ \ell^{\,\rho_\ell(p^j)/2}.
	\end{equation}
	Equivalently, the $p$-power subgroup must be at least the square root
	of the split or norm-one torus containing it.
\end{heuristic}

The heuristic is simply that $\Gamma=\{v+v^{-1}:v\in\mu_n\}$, a set of
$(n+1)/2$ elements contained in $\F_Q$, should meet the translate
$-1-\Gamma$ in about $\abs{\Gamma}^2/Q$ points.  Each nondegenerate trace
intersection lifts to four ordered pairs $(u,v)$, producing the estimate
$n^2/Q$.  By Lemma~\ref{lem:dickson}, $\Gamma$ is the root set of
$sF_n(s)$.
It matches the data: over $\F_{491}$ the ratio $49^{2}/491\approx4.9$
predicts a handful of solutions and there are $8$.
Over $\F_{9409}=\F_{97^2}$ the group $\mu_{49}$ is nonsplit and its trace
values lie in $\F_{97}$, so the relevant ratio is
$49^2/97\approx24.8$, consistent with the observed $32$ solutions.
Over $\F_{2^{21}}$
(the case $p=7$, $j=2$, $\ell=2$) the ratio is $49^{2}/2^{21}\approx0.001$
and there are none.

On the base-$2$ plateau, put $\rho=\rho_2(p)$.
The group $\mu_{p^c}$ lies in a torus of size $2^\rho-1$ or
$2^\rho+1$, while \eqref{eq:window} requires $p^c\gtrsim2^{\rho/2}$.
This isolates the approximate window
\begin{equation}\label{eq:2window}
	c\log_2 p+O(p^{-c})\ \leq\ \rho_2(p)\
	\lesssim\ 2c\log_2 p .
\end{equation}
Equivalently, $2^\rho-1$ in the split case or $2^\rho+1$ in the nonsplit
case would need a plateau prime-power divisor at least as large as its
square root.

\begin{table}[h]
	\centering
	\begin{tabular}{rrrrrr}
		\toprule
		$p$ & $c_p$ & $\rho_2(p)$ & window \eqref{eq:2window} & $p^{2c}/2^{\rho_2(p)}$ & margin\\
		\midrule
		$1093$ & $2$ & $182$  & $(20.2,\,40.4]$ & $10^{-42.6}$  & $4.5\times$\\
		$3511$ & $2$ & $1755$ & $(23.6,\,47.1]$ & $10^{-514.1}$ & $37.3\times$\\
		\bottomrule
	\end{tabular}
	\caption{The two known base-$2$ Wieferich primes miss the counterexample
	window \eqref{eq:2window} by the stated factor.}
	\label{tab:window}
\end{table}

\begin{remark}
	A prime $p$ is Wieferich if and only if
	$p^2\mid2^{\ord[p]{2}}-1$.
	Hence the conjecture that every Mersenne number $2^{n}-1$ is squarefree
	would imply $c_p=1$ for all $p$ and, by
	Corollary~\ref{cor:nonwieferich}, would settle
	Problem~\ref{prob:main} outright.
	That conjecture is false, since $1093^2\mid2^{364}-1$.
	The heuristic asks for a different, still inaccessible statement:
	no plateau prime power $p^c$ dividing the relevant factor
	$2^{\rho_2(p)}\pm1$ should be as large as its square root.
	We are not aware of a technique that proves such a bound.
\end{remark}

\begin{remark}[why character sums cannot help]
	Writing $\psi$ for the canonical additive character of $\F_q$ and
	$S(t)=\sum_{u\in\mu_n}\psi(t(u+u^{-1}))$, the solution count is
	$N=q^{-1}\sum_{t}\psi(t)S(t)^{2}$.
	Each $S(t)$ with $t\ne0$ is an average of $(q-1)/n$ twisted
	Kloosterman sums, so $\abs{S(t)}\leq2\sqrt q$ and
	$\sum_{t}\abs{S(t)}^{2}=q(2n-1)$.
	These give only $N\leq2n-1$, which is the trivial bound.
	Even an optimal square-root cancellation estimate has error term
	$\gg\sqrt q$, which exceeds the main term $n^2/q$ whenever
	$n^{2}<q^{3/2}$.
	For $p=1093$ and $j=2$ this reads $1093^{4}\approx10^{12}$ against
	$2^{546}\approx10^{164}$.
	No equidistribution statement can therefore prove $N=0$; only an exact
	algebraic identity can.
\end{remark}

\section{Computations}\label{sec:computations}

All computations are exact except for two clearly marked places: the
numerical Mahler measure quoted in Section~\ref{sec:modular}, and the
complex-arithmetic sweep used to corroborate Theorem~\ref{thm:charzero}.
Every assertion about a finite field, a resultant, or a certificate is exact
integer arithmetic.

\begin{description}
	\item[\path{code/ai_prototypes/five_term_relations.py}]
	      Characteristic-agnostic solver.
	      Builds $\F_{\ell^{f}}$ from a Rabin-tested irreducible
	      polynomial, locates an element of exact order $p^j$, tabulates
	      $\Gamma=\{\theta^{b}+\theta^{-b}\}$ and intersects it with
	      $-1-\Gamma$.
	      Validated on $(\ell,p,j)=(2,5,1)$, where it recovers the eight
	      solutions coming from $\mathcal{R}_5$.
	\item[\path{code/ai_prototypes/plateau_scan.py}]
	      Enumerates all triples $(\ell,p,j)$ with $p>5$, $j\geq2$,
	      $c(\ell,p)\geq j$ and $\ell^{f_1}$ below a bound, runs the
	      solver on each, and emits and rechecks a certificate consisting
	      of the minimal polynomial of $\theta$ over $\F_\ell$ together
	      with $(a,b)$.
	      The scan over all bases $\ell\leq1000$ with
	      $\ell^{f_1}\leq2\times10^{7}$ produced the counterexamples cited
	      in Section~\ref{sec:counterexamples}.
	\item[\path{code/ai_prototypes/five_term_char2.py}]
	      Dedicated characteristic-two solver.
	      Elements of $\F_{2^{f}}$ are Python integers used as
	      $\F_2$-coefficient bit vectors; the field is re-coordinatized so
	      that $\theta=x$, making each successive power a shift and a
	      conditional \textsc{xor}.
	      Membership is decided on $64$-bit fingerprints --- legitimate
	      because addition is \textsc{xor}, so the fingerprint of a sum is
	      the \textsc{xor} of fingerprints --- and every fingerprint
	      collision is re-verified exactly.
	      For the Wieferich prime $p=1093$ and $j=2$, with
	      $n=1\,194\,649$ and $f=\ord[1093^2]{2}=364$, the search completes
	      in under two seconds and finds \emph{no} fingerprint candidates
	      at all, hence no solutions.
	      For $p=3511$ and $j=2$, with $n=12\,327\,121$ and
	      $f=\ord[3511^2]{2}=1755$, the search completes in under four
	      minutes and again finds no candidates.
	      Since $c_{1093}=c_{3511}=2$, these are the only plateau exponents
	      for the two known Wieferich primes, so together with
	      Theorem~\ref{thm:A} they settle Problem~\ref{prob:main} for both.
	      This is also an independent confirmation, by a method that never
	      forms a Fibonacci polynomial, of
	      $d(1093^2-1)=d(3511^2-1)=0$.
	\item[\path{code/ai_prototypes/five_term_resultants.py}]
	      Computes, for a given $p^j$, the complete table of relative norms
	      $\Nm_{K^{+}/\Q}(\Lambda_{1,r})$ over all classes of nonzero
	      ratios $r$, together with their factorizations, and reports
	      whether $2$ occurs.
	      This is the exact criterion of
	      Proposition~\ref{prop:dictionary}; running it on $(7,2)$
	      reproduces Table~\ref{tab:res49}.
	\item[\path{code/ai_prototypes/verify_p7_characteristics.py}]
	      Checks the prediction of Table~\ref{tab:res49}: for each prime
	      $\ell$ dividing one of the fourteen relative norms, the five-term
	      relation does have a solution over $\overline{\F}_\ell$ with $u$
	      or $v$ of exact order $49$.
	      All eight predicted characteristics were confirmed.
	\item[\path{code/ai_prototypes/sine_unit_obstruction.py}]
	      The computations of Section~\ref{sec:power}.
	      It tabulates the sine residues $\bar\lambda_x$, verifies the
	      fibre structure of Lemma~\ref{lem:fibre}, enumerates all
	      product-form collisions and their ratio classes (confirming
	      Theorem~\ref{thm:frobratio} for every $p\leq257$ with
	      $\ord[p]{2}\leq60$), and computes the power-residue symbols
	      $\Psi_\ell$ of Definition~\ref{def:psiell}.
	      For the $p$-primary symbol it checks
	      Theorem~\ref{thm:pstructure} by comparing the Fourier support of
	      $\Psi$ against the admissible eigencharacters $A$; over all
	      primes $7\leq p\leq199$ with $\ord[p]{2}\leq200$ the support was
	      contained in $A$ in every case, $\Psi$ vanished identically
	      exactly when $\ord[p]{2}$ was even, and the ratio classes
	      admitted by \eqref{eq:psi-collision} with $\ell=p$ always
	      contained the ratio classes of the actual solutions.
	      The obstruction was empty --- hence a proof of $d(p-1)=0$ ---
	      for $p=7,23,47,71,79,89,103,167,191,199$, and exactly sharp for
	      $p=31$ and $p=127$.
	      The same script verifies the eigenrelation of
	      Theorem~\ref{thm:pstructure}(1) and the descent of
	      Theorem~\ref{thm:pstructure}(3) at levels $j=2,3$, verifies
	      Corollary~\ref{cor:1093-ppart} for $p=1093$, $j=2$ inside
	      $\F_{2^{364}}$, and runs the test over all primes
	      $\ell\mid2^{f}-1$ of Lemma~\ref{lem:psi-rules}: for $p=17$ and
	      $p=31$, where solutions exist, every $\ell$ admits pairs, while
	      for $p=23$ and for $(p,j)=(7,2)$ every $\ell$ excludes them.
	\item[\path{code/ai_prototypes/plateau_ppart_3511.py}]
	      The plateau test of Example~\ref{ex:3511-plateau}.
	      It splits $\Phi_{3511}$ over $\F_2$ by a trace map to obtain
	      $\F_{2^{1755}}$, locates $\theta$ of exact order $3511^{2}$,
	      confirms the eigenrelation, and computes $\Psi$ on a sample of
	      the group $G_p$ of principal units by baby-step giant-step in
	      $\mu_{p^{2}}$.
	      With $140$ samples it exhibits $4$ pairs satisfying the
	      $p$-primary necessary condition, against $2.79$ expected at
	      random: the obstruction that settles $j=1$ does not settle
	      $j=2$.
	\item[\path{code/cn_experiments/verify_all.py}]
	      Verifies the shifted-Dickson recurrence and power-sum formula,
	      the Hasse-derivative generating function, the
	      Multiplicity-Two identity $C_n^{[2]}(Y)=Y^{-2}$, and the bridge
	      \eqref{eq:Cn-bridge}.  The supporting scripts
	      \path{code/cn_experiments/integer_verify.py} and
	      \path{code/cn_experiments/factorization_norms.py}
	      verify the integer lift, telescoping identity, real-cyclotomic
	      factorization, and discriminant/resultant certificates.
	\item[\path{code/cn_experiments/collision_parity.py}]
	      Checks the collision involution, the internal and mixed
	      resultant formulas, and the plateau inversion criterion.
	      \path{code/cn_experiments/relative_norm_identity.py} verifies
	      \eqref{eq:relative-collision-norm} and its resultant descent,
	      while \path{code/cn_experiments/dynatomic_orbits.py} verifies the covering-degree
	      bound and the doubling-orbit factorization of
	      Theorem~\ref{thm:dynamics-orbits}.
	\item[\path{code/ai_prototypes/circular_module.py}]
	      The computations of Sections~\ref{sec:module},
	      \ref{sec:traceobs}, and~\ref{sec:kloosterman}.
	      It verifies Theorem~\ref{thm:trace} for every index and every
	      $p\leq257$ with $\ord[p]{2}\leq40$ and for
	      $(p,j)=(7,2),(23,2)$; verifies the level pattern of
	      Theorem~\ref{thm:twoorbit} and its two proof ingredients for
	      $p=7,23,47,71,79,103,167,191,199$; confirms the group-ring
	      structure of the semi-local residue module by Smith normal form;
	      checks the freeness of the admissible-system parameter space
	      against the true system; constructs explicit injective phantom
	      systems with primitive collisions for $(1093,2)$ and
	      $(3511,2)$; and verifies the Kloosterman nullity formulas
	      against direct enumeration for all Mersenne primes up to
	      $131071$ and all Fermat primes up to $65537$.
	\item[\path{code/ai_prototypes/plateau_endpoint.py}]
	      Verifies the endpoint Sylow reduction, the box-plus-centre form,
	      the global trace law \eqref{eq:global-trace}, the cases
	      $p=73,89,233$, the elliptic model of
	      Proposition~\ref{prop:elliptic-translation}, the additive
	      distribution law, and the strengthened phantom systems at
	      $(1093,2)$ and $(3511,2)$.
	\item[\path{code/ai_prototypes/sp4_explicit.py}]
	      Checks the symplectic form, characteristic polynomial, and
	      element order in the explicit level-one $\Sp_4$ examples used to
	      test the torus classification.
\end{description}

\subsection{The complete resultant table for $p^j=49$}

By Proposition~\ref{prop:dictionary} the whole of Problem~\ref{prob:main}
for a given $(p,j)$ is the assertion that finitely many explicit integers
are odd.
For $p^j=49$ the ratio $r=b/a$ ranges over the nonzero residues modulo
$49$, and the relative norm $\Nm_{K^{+}/\Q}(\Lambda_{1,r})$ depends only on
the class $\{\pm r^{\pm1}\}$ (only $\{\pm r\}$ when $r$ is not a unit).
There are eleven unit classes and three non-unit classes.
Table~\ref{tab:res49} lists them all.

\begin{table}[h]
	\centering
	\begin{tabular}{lrl}
		\toprule
		class & $\Nm_{K^{+}/\Q}(\Lambda_{1,r})$ & factorization\\
		\midrule
		$\{1,48\}$              & $13555459$ & $97\cdot139747$\\
		$\{2,24,25,47\}$        & $1$        & unit (Corollary~\ref{cor:units})\\
		$\{3,16,33,46\}$        & $1$        & unit (Corollary~\ref{cor:units})\\
		$\{4,12,37,45\}$        & $3527$     & prime\\
		$\{5,10,39,44\}$        & $1373$     & prime\\
		$\{6,8,41,43\}$         & $97$       & prime\\
		$\{9,11,38,40\}$        & $881$      & prime\\
		$\{13,15,34,36\}$       & $1$        & unit (sporadic)\\
		$\{17,23,26,32\}$       & $29009$    & prime\\
		$\{18,19,30,31\}$       & $491$      & prime\\
		$\{20,22,27,29\}$       & $97$       & prime\\
		\midrule
		$\{7,42\}$              & $97$       & prime\\
		$\{14,35\}$             & $1$        & unit\\
		$\{21,28\}$             & $197$      & prime\\
		\bottomrule
	\end{tabular}
	\caption{The fourteen relative norms for $p^j=49$; the last three rows
	are the classes with $7\mid r$, i.e.\ with $\theta^{b}$ of order
	dividing $7$.
	Every entry is odd, which is exactly the assertion of
	Problem~\ref{prob:main} for $p=7$, $j=2$.}
	\label{tab:res49}
\end{table}

\begin{corollary}\label{cor:p7-complete}
	For $p=7$ and $j=2$, the residue characteristics $\ell$ in which
	\eqref{eq:main} has a solution with $\theta^{a}$ or $\theta^{b}$ of
	exact order $49$ are exactly
	\[
		97,\quad 197,\quad 491,\quad 881,\quad 1373,\quad 3527,\quad 29009,\quad 139747 ,
	\]
	and $2$ is not among them.
\end{corollary}

\begin{proof}
	By Proposition~\ref{prop:dictionary}(1) and (4), a prime $\ell$ admits
	a solution with ratio $r$ if and only if $\ell$ divides
	$\Res(\Phi_{49},G_r)$.
	Table~\ref{tab:res49} lists the complete factorizations.
	The verification that each of these primes really does admit a solution
	is the content of
	\path{code/ai_prototypes/verify_p7_characteristics.py} and of the scan
	in \path{code/ai_prototypes/plateau_scan.py}.
	The characteristics $\ell=139747$ and $\ell=197$ are degenerate: the
	first arises only from the class $r=\pm1$, where $v=u^{\pm1}$, and the
	second only from the class $r=\pm21$, where $7\mid b$.
\end{proof}

Corollary~\ref{cor:p7-complete} is the sharpest available description of the
situation.
For $p=7$, $j=2$ the truth of Problem~\ref{prob:main} is not a structural
fact; it is the statement that seven specific primes, produced by the
factorizations of eleven cyclotomic resultants, happen not to include $2$.

\section{Summary}\label{sec:summary}

\begin{enumerate}
	\item The shifted-Dickson sequence $C_n$ is the Newton power-sum
	      sequence of $Q_S$.  Its repeated roots are exactly the dividing
	      reciprocal quartics, every such root has multiplicity two, and
	      $d(n-1)=4\deg B_n$ (Theorems~\ref{thm:Cn-mult2}
	      and~\ref{thm:Cn-bridge}).  This removes all ambiguity about
	      multiplicities in the polynomial formulation.
	\item The degree argument proves
	      $\ord{\theta^a},\ord{\theta^b}\leq p^{c_p}$.
	      The universal theorem is therefore equivalent to the one
	      endpoint statement
	      \[
		      B_{p^{c_p}}=B_p,
	      \]
	      or equivalently to the gcd in \eqref{eq:endpoint-gcd} being one
	      (Theorem~\ref{thm:endpoint-equivalence}).  All primes below
	      $2^{64}$ satisfy the theorem.
	\item The integer lift turns one stabilization step into the oddness
	      of a discriminant and a resultant.  Real-cyclotomic
	      factorization shows that these are two independent obstructions:
	      an internal primitive collision norm and a primitive--lower
	      mixed norm (Theorem~\ref{thm:two-collision-norms}).  The mixed
	      norm is automatically odd off the plateau, but neither parity is
	      forced by cyclotomic algebra on the plateau.
	\item Arithmetic dynamics makes $D_p$ a critically fixed map
	      semiconjugate to $u\mapsto u^p$ through
	      $u\mapsto u+u^{-1}$.  It reproves mixed descent from the covering
	      degree and factors $B_{p^j}$ by doubling orbits
	      (Theorems~\ref{thm:dynamics-covering}
	      and~\ref{thm:dynamics-orbits}).  On the plateau the degree does
	      not grow, so the bound is vacuous and monodromy leaves the same
	      finite collision-orbit problem.
	\item The elliptic model identifies a relation with a point $P$ on
	      $Y^2+XY=X^3+X$ such that both $X(P)$ and $X(P+T)$ lie in
	      $\mu_{p^j}$, where $T$ has order four
	      (Proposition~\ref{prop:elliptic-translation}).  The condition is
	      on coordinates, not on a subgroup of the elliptic curve, so
	      isogenies and point counts do not imply endpoint avoidance.
	\item The $\Sp_4$ model identifies a counterexample with a regular
	      semisimple element of order $p^j$ and trace one.  Maximal-torus
	      theory restricts it to two like $\SL_2$ factors or a cyclic torus
	      of order $q^2\pm1$, but $p^j$-torsion is abundant on the plateau
	      and the trace-one condition is exactly the original incidence
	      (Theorems~\ref{thm:sp4-equivalence} and~\ref{thm:sp4-tori}).
	\item Characteristic zero has only the order patterns
	      $(5,5),(3,4),(2,6)$, so for $p>5$ the relevant algebraic integer
	      is nonzero.  Cyclotomic units eliminate exactly the structural
	      ratio classes $\{\pm2^{\pm1},\pm3^{\pm1}\}$ and no general
	      additional class (Theorem~\ref{thm:units}).
	\item Reduction modulo the rational ideal $(2)$ is impossible by the
	      Dvornicich--Zannier theorem or the elementary weight argument,
	      but the problem concerns one prime above two.  Reciprocal
	      symmetry reaches at most four of the
	      $g\geq p^{j-1}$ primes that an all-primes theorem needs
	      (Theorem~\ref{thm:gap}).
	\item The sharp local statement is the Hilbert--90 congruence
	      $s_{j,a}s_{j,b}\equiv1$ at one prime above two.  The semi-local
	      target is
	      $\Z[\bar G]/(\sigma_2-2)$, while the Frobenius, sign,
	      distribution, and norm axioms leave a free parameter space
	      $(E^\times)^{g^+}$ (Theorems~\ref{thm:semilocal-module}
	      and~\ref{thm:admissible-freeness}).  Explicit phantom systems
	      prove that those axioms do not imply collision avoidance.
	\item Power-residue symbols and Stickelberger eigenspaces give strong
	      necessary conditions and settle several level-one families.
	      Their descent becomes the identity $0=0$ on a Wieferich plateau
	      (Proposition~\ref{prop:plateau-degenerate}); this failure is
	      structural, not a defect of the calculation.
	\item Additive trace identities go farther.  The Artin--Schreier law
	      proves the index-two family, including $d(3510)=0$, and the
	      global trace law proves
	      $d(72)=d(88)=d(232)=0$
	      (Theorems~\ref{thm:trace}, \ref{thm:twoorbit},
	      and~\ref{thm:global-trace}).  At the full-torus extremes the
	      count is an exact binary Kloosterman sum
	      (Theorem~\ref{thm:kloosterman}).  The known linear trace
	      constraints still admit formal plateau collisions, so a general
	      proof must use more of the nonlinear addition law.
	\item Coding theory is decisive off the plateau: the local cyclic code
	      is an interleaving of lower-level codes, forcing every word of
	      weight at most five to descend
	      (Theorem~\ref{thm:code-interleaving}).  On a nonsplit plateau it
	      is a shortened Zetterberg code, whose general minimum-distance
	      theory stops at weight five, exactly the forbidden weight.
	\item The theorem is false after replacing characteristic two by other
	      residue characteristics (Theorem~\ref{thm:counterexamples}).
	      Hence characteristic-zero classification, cyclotomic Galois
	      structure, torus classification, and base-independent monodromy
	      cannot suffice.  A proof must exploit a special additive feature
	      of characteristic two.
	\item The size heuristic predicts a counterexample only when
	      $\rho_2(p)$ lies in the narrow window \eqref{eq:2window}.
	      The known Wieferich primes
	      $1093$ and $3511$ lie far outside it, and exact searches verify
	      their endpoints.  The complete $p^j=49$ resultant table likewise
	      contains no factor two.
\end{enumerate}

The minimal unresolved assertion is therefore precise:
\[
	\boxed{\text{For every }p>5\text{ with }c_p\geq2,\quad B_{p^{c_p}}=B_p.}
\]
Equivalently, every endpoint collision in $\mu_{p^{c_p}}$ must already lie
in $\mu_p$, or the gcd in \eqref{eq:endpoint-gcd} must be one.  Proving
this single statement proves $d(p^k-1)=d(p-1)$ for every prime $p$ and
every $k$.  The accumulated attacks show that the missing input must be
both characteristic-two and subgroup-sensitive; none of the currently
isolated degree, norm, trace, code, dynamics, or group-theoretic invariants
alone has that strength.

\bibliographystyle{amsalpha}
\bibliography{refs.bib}

\end{document}
